\documentclass[11pt]{amsart}
\usepackage[T1]{fontenc}
\usepackage{lmodern,amsmath,amssymb,amsthm,mathtools,booktabs,array,longtable}
\usepackage{mathrsfs}
\usepackage[a4paper,margin=28mm]{geometry}
\usepackage{xurl}
% Keep file and JSON path breaks at separators rather than inside words.
\newcommand{\filepath}[1]{\begingroup
  \def\UrlBreaks{\do\/\do\_}\def\UrlBigBreaks{}\Urlmuskip=0mu
  \path{#1}\endgroup}
\usepackage[colorlinks=true,linkcolor=blue,citecolor=blue,urlcolor=blue]{hyperref}
\hypersetup{pdftitle={Theorem (1+1.75) on the Goldbach Conjecture},pdfauthor={Slava Chan}}
\newtheorem{theorem}{Theorem}[section]
\newtheorem{proposition}[theorem]{Proposition}
\newtheorem{lemma}[theorem]{Lemma}
\newtheorem{corollary}[theorem]{Corollary}
\theoremstyle{definition}
\newtheorem{definition}[theorem]{Definition}
\newtheorem{remark}[theorem]{Remark}
\DeclareMathOperator{\supp}{supp}
\DeclareMathOperator{\li}{li}
\newcommand{\eps}{\varepsilon}
\newcommand{\one}{\mathbf 1}
\allowdisplaybreaks[1]
\setlength{\emergencystretch}{3em}
\title{Theorem (1+1.75) on the Goldbach Conjecture}
\author{Slava Chan}
\address{Independent researcher}
\email{ch4nslava@gmail.com}
\date{6 November 2026}
\subjclass[2020]{11P32, 11N36, 11N13}
\keywords{Chen's theorem, double sieve, semiprimes, primes in arithmetic progressions}
\begin{document}
\begin{abstract}
For fixed $1\le a\le2$, let Proposition $(1+a)$ assert that every
sufficiently large even integer $N$ can be written as
\[
 N=p+rq,\qquad r\le q^{a-1},
\]
where $p,q$ are primes and $r$ is either $1$ or a prime.
At $a=1$, this is the binary Goldbach conjecture for sufficiently large
even integers; at $a=2$, it is Chen's theorem. We prove unconditionally
that every sufficiently large even integer admits such a representation
with $a=1.75$, establishing Proposition $(1+1.75)$ and strengthening
the earlier $(1+1.9)$ theorem of Li and Liu~\cite{LiLiu}.
The proof subtracts a positive upper-sieve estimate for excluded
semiprimes from a Chen-count lower bound obtained by restricting paired
terms in an eleven-term Chen weight and applying the double sieve.
The threshold for $N$ is ineffective: the proof does not provide an
explicit computable threshold.
\end{abstract}
\maketitle
\section{Introduction}
\label{sec:current-statement}

For fixed $1\le a\le2$, we say that \emph{Proposition $(1+a)$} holds
if every sufficiently large even integer $N$ has a representation
\[
 N=p+rq,\qquad r\le q^{a-1},\qquad
 p,q\text{ prime},\qquad r=1\text{ or prime}.
\]
When $a=1$, the condition forces $r=1$, giving the binary Goldbach
conjecture for sufficiently large even integers. When $a=2$, ordering
the two factors gives Chen's theorem. The intermediate parameters ask
how much smaller one prime factor of a composite remainder can be
required to be.

We prove the following theorem.

\begin{theorem}
\label{cor:headline}
Proposition $(1+1.75)$ holds unconditionally. Thus every sufficiently
large even integer can be written as $N=p+rq$, where $p,q$ are primes,
$r$ is either $1$ or a prime, and $r\le q^{3/4}$.
\end{theorem}

The thresholds for $N$ in Theorems~\ref{cor:headline}--\ref{thm:current-main}
and Corollaries~\ref{cor:three-quarters}--\ref{cor:current-explicit} are ineffective.

The parameter $1.75$ can be reduced to $1.74816964$ and further to $1.74816865$,
with different quantitative lower bounds for the number of
representations; see Corollary~\ref{cor:current-explicit}.

Li and Liu~\cite[Theorem~1.1]{LiLiu} prove Proposition $(1+1.9)$
and remark that their parameter can be reduced to $1.894$.
Theorem~\ref{cor:headline} strengthens this factor restriction.
Both statements concern every sufficiently large even integer.

For the quantitative argument put $b=a-1$ and count prime summands
with $r\le q^b$. Prime remainders $r=1$ remain admissible, since
$1\le q^b$. The next theorem gives a lower bound for every fixed
$1/2<b<1$.

Throughout the paper, $N>2$ tends to infinity through the even integers,
unsubscripted logarithms are natural, $\log_N x=\log x/\log N$,
and $\Omega(n)$ counts prime factors with multiplicity. Put
\begin{equation}
 C_2=\prod_{\substack{\ell>2\\\ell\text{ prime}}}
 \left(1-\frac1{(\ell-1)^2}\right),\qquad
 C_N=C_2\prod_{\substack{\ell\mid N,\ \ell>2\\\ell\text{ prime}}}
 \frac{\ell-1}{\ell-2},\qquad
 H_N=C_N\frac{N}{\log^2 N}.
 \label{eq:current-normalization}
\end{equation}
For any positive even integer $M$, $C_M$ denotes the same product with $M$ in place of $N$.
We count distinct prime summands by
\begin{align}
 D_{1,2}(N)&=\#\{p<N:p\text{ prime},\ N-p>1,
                                      \ \Omega(N-p)\le2\},
                                      \label{eq:current-d12}\\
 D_b(N)&=\#\{p<N:p\text{ prime},\ N-p=rq,
       \ q\text{ prime},\ r=1\text{ or prime},\ r\le q^b\},
                                      \label{eq:current-db}\\
 B_b(N)&=\#\{p<N:p\text{ prime},\ N-p=rq,
                \ r,q\text{ prime},\ q^b<r\le q\}.
                                      \label{eq:current-bb}
\end{align}
The ordered factorization $r\le q$ is unique. In particular, a composite
remainder contributes once; a square contributes to $B_b$, whereas a prime
remainder contributes to $D_b$.

\begin{theorem}
\label{thm:current-main}
Let
\begin{equation}
 K=\frac{8452171}{1250000}=6.7617368,
 \qquad c=\frac{49044063}{25000000}=1.96176252.
 \label{eq:current-safe-constants}
\end{equation}
For each fixed $1/2<b<1$,
\begin{equation}
 \liminf_{\substack{N\to\infty\\2\mid N}}
      \frac{D_b(N)}{H_N}\ge c-K\log(1/b).
 \label{eq:current-general-theorem}
\end{equation}
\end{theorem}

Within $1/2<b<1$, this lower bound is positive precisely when
$b>\exp(-c/K)$; see Section~\ref{sec:current-transfer}.

\begin{corollary}
\label{cor:three-quarters}
For every sufficiently large even integer $N$,
\[
                  D_{3/4}(N)\ge0.0165H_N.
\]
\end{corollary}

\begin{corollary}
\label{cor:current-explicit}
Every sufficiently large even integer $N$ satisfies
\begin{align}
 D_{0.74816964}(N)&\ge10^{-5}H_N,
                                  \label{eq:current-main-count}\\
 D_{0.74816865}(N)&\ge10^{-6}H_N,
                                  \label{eq:current-secondary-count}\\
 D_{0.9}(N)&\ge1.24889H_N.
                                  \label{eq:current-common-exponent}
\end{align}
\end{corollary}

The first two inequalities establish Propositions $(1+1.74816964)$
and $(1+1.74816865)$ with respective count coefficients $10^{-5}$ and
$10^{-6}$ in the normalization~\eqref{eq:current-normalization}.
The estimate at $b=0.9$ permits a comparison under the same factor
restriction.

\begin{proof}[Proof of Theorem~\ref{cor:headline}]
Corollary~\ref{cor:three-quarters} gives $D_{3/4}(N)\ge0.0165H_N>0$
for every sufficiently large even $N$. Thus Proposition $(1+1.75)$ holds.
\end{proof}

\subsection{Related work and the method}
In the normalization~\eqref{eq:current-normalization}, the result of
Li and Liu~\cite[Theorem~1.1]{LiLiu} has coefficient $0.0004$ at
$b=0.9$; their stated refinement corresponds to $b=0.894$.
At $b=0.9$, our $D_b(N)$ coincides with Li--Liu's $D_{1,1.9}(N)$:
both include the prime-remainder case $r=1$, count each admissible
prime summand once, and use the same singular-series normalization.
Thus the coefficients refer to the same counting function.
Our coefficient follows from the combination $c-K\log(1/b)$: $c$ is
supplied by the lower Chen-count construction, and $K$ by the positive
upper majorant.
The smaller exponents in Corollary~\ref{cor:current-explicit} retain
an explicit positive count. Their discussion of lower and marked sieve
terms~\cite[Sections~3 and 8.2]{LiLiu} also identifies a constraint
that matters here: an improved upper-sieve distribution level must be
matched to the coefficient class of each term where it is used.

The double sieve transfers information between upper and lower sieve
remainders. We use Wu's formulation~\cite{Wu}, with finite positive
transitions for the two remainder functions.
Li~\cite[Theorem~1.1]{RunboLi} gives coefficient $1.9728$ for the total
Chen count in the same normalization. The coefficient $c=1.96176252$
used here is proved within the moving-residue coefficient class and
positive marked families of the factor-restricted argument. Each lower
transition retains its actual residual level and all prefix conditions.

The distribution estimate must be uniform as the residue varies with
$N$. Maynard's triply well-factorable estimate~\cite[Definition~2 and
Theorem~1.1]{MaynardII} has a constant depending on a fixed residue.
Lichtman's large-residue statement~\cite[Theorem~1.7]{Lichtman} has
level $(5-\vartheta)/8-\epsilon=153/256-\epsilon$ at
$\vartheta=7/32$. Here the residue is comparable with the complete
local product. The reciprocal phase
in~\eqref{eq:current-retained-phase} is retained, since replacing it
by $1$ in this range need not give a vanishing error.
Appendices~\ref{app:medium} and~\ref{app:classical}, together with the
factorization in Appendix~\ref{app:coefficients}, establish the
restricted distribution input at support exponent $0.599$ for the
finite upper and lower coefficient families used in this paper.
At each application we retain the selected factorization, full and
complementary square conditions, and positivity of the marked family.

The proof subtracts the excluded semiprimes from the full Chen count.
For the upper estimate we construct a nonnegative divisor majorant,
compute its density, and prove its distribution estimate for the same
coefficient. For the lower estimate we restrict the two-prime and
four-prime terms of an eleven-term Chen weight in pairs. An injection
preserves the pointwise weight inequality, while a lower bound for
Buchstab's function on the restricted fibres improves the coefficient.

The nonnegative upper majorant and its coefficient realization are
constructed here. For the lower bound, we use Wu's double-sieve framework
with the paired deletion and coefficient realizations established here.
Together with the finite transitions and directed enclosures, these
constructions determine the constants $K,c$. Their mathematical
definitions and certificate correspondence are given in
Appendix~\ref{app:numeric}.

\section{Reduction and outline of the proof}
\label{sec:outline}

\paragraph{Proof outline and verification.}
On a first reading, Lemma~\ref{lem:numerical-endpoints} may be taken as
the finite numerical input: follow the upper construction, the eleven-term
lower construction and paired deletion, and then the final subtraction in
Section~\ref{sec:current-transfer}. For a complete verification, return to
the appendices for the distribution estimates, coefficient realizations,
finite transitions and directed enclosures, and check their correspondence
with the Numerical Supplement described in Appendix~\ref{app:numeric}.

Unique factorization gives the exact identity
\begin{equation}
                 D_b(N)=D_{1,2}(N)-B_b(N).
                 \label{eq:reduction}
\end{equation}
Thus it suffices to prove
\begin{equation}
 \liminf_{\substack{N\to\infty\\2\mid N}}
       \frac{D_{1,2}(N)}{H_N}\ge c,
 \qquad
 \limsup_{\substack{N\to\infty\\2\mid N}}
       \frac{B_b(N)}{H_N}\le K\log(1/b).
 \label{eq:two-estimates}
\end{equation}
The factor $\log(1/b)$ comes from the excluded factor region. For
$\beta=b/(1+b)$, the limiting logarithmic region of the ordered bad
semiprimes is $\beta<t\le1/2$, and
\[
 \int_\beta^{1/2}\frac{dt}{t(1-t)}=\log(1/b).
\]
To turn this mass into a count with a prime complement, we apply a
positive upper sieve whose distribution and density have both been
established for its actual coefficients.

Section~\ref{sec:current-distribution}
establishes the distribution estimate for the actual upper and lower
coefficients. Section~\ref{sec:current-upper} constructs the positive
upper sieve and obtains $K$. Sections~\ref{sec:current-chen-weight}
and~\ref{sec:current-lower} give the lower construction; its switched
costs use the upper bound already proved. Section~\ref{sec:current-numeric}
states the finite numerical inequalities and their error allowances.
Finally Section~\ref{sec:current-transfer} proves the second estimate
in~\eqref{eq:two-estimates} and completes the theorem. In particular,
the construction of $K$ uses neither the lower double-sieve tables nor
the Chen count obtained from them.

\paragraph{Choice of parameters.}
At each fixed accuracy, the finite collection of policies and marked queries,
the number of generations, cutoffs defining fixed compact parameter domains,
approximation tolerances and
support reserves are chosen before $N$ tends to infinity. Prime intervals,
weights and arithmetic coefficients may still depend on $N$ through these
fixed rules. The common choices and error propagation are specified in
Section~\ref{sec:lower-limits} and \eqref{coef:parameter-order}; the strict
scalar surpluses in Section~\ref{sec:current-numeric} allow one fixed accuracy
to be chosen for the final count.

\paragraph{Notation index.}
The following symbols keep the same meaning in the main text and appendices.
\begin{center}
\small
\begin{tabular}{@{}>{\raggedright\arraybackslash}p{.23\textwidth}>{\raggedright\arraybackslash}p{.49\textwidth}>{\raggedright\arraybackslash}p{.22\textwidth}@{}}
\toprule Symbol & Meaning & Location\\
\midrule
$E_{\rm up}$; $R_{\rm low}$ & Upper and lower scalar endpoints & \eqref{eq:current-E}; \eqref{eq:current-R}\\[2pt]
$\mathscr E_0,\mathscr E_1,\mathscr E$ & Paired rectangle families and their union & Section~\ref{sec:current-chen-weight}\\[2pt]
$a_{\rm s},h_{\rm s},d_{\rm s}$ & Fixed sieve cutoffs & Section~\ref{sec:current-chen-weight}\\[2pt]
$H(s),h(s)$; $H_m,h_m$ & Upper and lower defects; generation-$m$ envelopes & Section~\ref{sec:current-lower}\\[2pt]
$B$; $\mathscr B$ & Arithmetic prime prefix; geometric rectangle or box & Sections~\ref{sec:current-distribution}, \ref{sec:current-lower}\\[2pt]
$H_N$; $B_{\rm low}$ & Count normalization; each lower error allowance & \eqref{eq:current-normalization}; \eqref{eq:current-B}\\[2pt]
$A_+,A_-$ & Normalized linear-sieve numerators $sF/(2e^\gamma),sf/(2e^\gamma)$ & Section~\ref{sec:current-lower}\\[2pt]
$\Omega$; $\omega_{\rm env}$ & Prime-factor count with multiplicity; Buchstab upper envelope & Section~\ref{sec:current-statement}; \eqref{sv:omega-envelope}\\[2pt]
$\rho_N^{\rm up},\rho_N^{\rm low}$ & Actual upper and lower divisor coefficients; $\rho_N$ denotes a generic admitted coefficient & Sections~\ref{sec:current-upper}, \ref{sec:current-lower}\\[2pt]
\bottomrule
\end{tabular}
\end{center}

We use the standard linear sieve functions $F,f$, normalized by
\[
 sF(s)=2e^\gamma\quad(1\le s\le3),\qquad
 sf(s)=0\quad(1\le s\le2),
\]
and subsequently by
$(sF(s))'=f(s-1)$ and $(sf(s))'=F(s-1)$ on their respective
ranges. Buchstab's function satisfies
\[
 \omega(s)=1/s\quad(1\le s\le2),\qquad
 (s\omega(s))'=\omega(s-1)\quad(s>2).
\]
Here $P^-(n)$ is the least prime factor for $n\ge2$, and $P^-(1)=+\infty$. The finite interval data used for these
functions are discussed in Appendix~\ref{app:numeric}.

\subsection{A positive sieve coefficient}
The upper estimate uses the following single coefficient theorem.
Its three conclusions concern the same divisor coefficient.
Here ``positive'' means that the resulting divisor sum $R_N(n)$ is
a nonnegative majorant. The individual coefficients $\rho_N^{\rm up}(q)$
may have either sign.

\begin{theorem}
\label{thm:positive-coefficient}
Fix $a_*>0$, a bound on the number of labelled prime factors, and a
finite collection of smooth weights $W_N(p_1,\ldots,p_k)$ supported on
\[
 p_i\ge c_*N^{a_*},\qquad p_1\cdots p_k\le C_*N,
\]
with fixed constants $c_*,C_*>0$, uniformly bounded size, and fixed
polylogarithmic bounds for their mixed Euler derivatives. Apart from
these smooth geometric conditions, impose no additional arithmetic
filter. Put $\rho_0=19/96$, $\theta_*=599/1000$, and
\[
 P_N(y)=\prod_{\substack{p<y\\p\nmid N}}p,\qquad
 X_W^{(q)}=\sum_{(p_1\cdots p_k,q)=1}W_N(p_1,\ldots,p_k).
\]
There are real coefficients $\rho_N^{\rm up}(q)$ and a constant $M$, independent
of $N$, with squarefree support $q\le N^{\theta_*}$, $(q,N)=1$,
and $|\rho_N^{\rm up}(q)|\le M$, such that
\begin{equation}
 R_N(n):=\sum_{q\mid n}\rho_N^{\rm up}(q)
   \ge\one_{(n,P_N(N^{\rho_0}))=1}\ge0
                           \qquad(n\in\mathbb Z),\label{eq:core-positivity}
\end{equation}
\begin{equation}
 \sum_q\frac{\rho_N^{\rm up}(q)}{\varphi(q)}
   \le(K+o(1))\frac{C_N}{\log N},\label{eq:core-density}
\end{equation}
\begin{multline}
 \sum_q\rho_N^{\rm up}(q)\left\{
 \sum_{p_1\cdots p_k\equiv N\ (q)}W_N(p_1,\ldots,p_k)
       -\frac{X_W^{(q)}}{\varphi(q)}\right\}\\
   =O_{A,W}(N(\log N)^{-A})\qquad(A>0).
                           \label{eq:core-distribution}
\end{multline}
The finite sum defines $R_N(0)$. The finite policy, approximation
parameters and support margins are fixed before $N$ tends to infinity
through all even integers. The constants may depend on this fixed
collection and its derivative bounds, but are uniform in $N$.
\end{theorem}

The coefficient is constructed in Section~\ref{sec:current-upper}.
The distribution statement is proved first in
Section~\ref{sec:current-distribution} for the precise coefficient
representations needed by that construction.

\section{Distribution estimates for the sieve coefficients}
\label{sec:current-distribution}

\begin{definition}
A finite prefix policy specifies finitely many regions in the
coordinates $t_i=\log_Np_i$ of ordered prime prefixes and a sieve level
on each region. The decision at a prefix is fixed before its children
are refined. A selected branch is admissible when its full-square and
selected-complement conditions in
\eqref{coef:full-squares}--\eqref{coef:selected-squares} have the required
strict margins; an inherited branch keeps its parent decision.
A marked family is a nonnegative finite weighted family of prefixes
$B$ at a common total level $Q$, with residual level $Q/B$ and the
original prime ceiling retained. A marked query is a sifted count for
one such family. References below to an admitted coefficient mean a
coefficient obtained from these policies and finite marked transitions,
with the three-factor representations of
Appendix~\ref{app:coefficients}.
\end{definition}

\subsection{A local estimate with a moving residue}

The distribution problem is local in the complete product. We denote
the convolution lengths by $U,V$, write $x=UV$, and use $\mathfrak a$
for the residue; $N$ remains the global even target.
The estimate proved in Appendix~\ref{app:medium} is the following.

\begin{proposition}[Restricted medium convolution]
\label{prop:current-medium}
Fix $\vartheta=7/32$, $0<\sigma<1/100$, $c_0>0$, a divisor order
$B$, and $0\le\kappa\le\sigma/100$. Let $U,V,Q_1,Q_2,Q_3\ge1$. Suppose
\begin{align}
 c_0x&\le |\mathfrak a|\le x^{1+\kappa},&
 Q_1&<Ux^{-\sigma}, \notag\\
 U^2Q_2Q_3^2&<x^{1-8\sigma},&
 |\mathfrak a|^{\vartheta}UQ_1Q_2^5Q_3^2&<x^{2-15\sigma},
       \label{eq:current-medium-ranges}\\
 U^{1+2\vartheta}Q_1Q_2^{5-2\vartheta}Q_3^{2+2\vartheta}
      &<x^{2-15\sigma}.&&\notag
\end{align}
Let $\alpha$ be one-bounded, supported on an individual interval in
$[U,2U]$ and on
$P^-(n)>x^{1/(\log\log x)^3}$.  Assume its auxiliary-coprimality
Siegel--Walfisz estimate: for every $A>0$, every $k,q\ge1$ and
$(b_0,q)=1$,
\begin{equation}
 \left|\sum_{\substack{n\equiv b_0\pmod q\\(n,k)=1}}\alpha_n
 -\frac1{\varphi(q)}\sum_{(n,kq)=1}\alpha_n\right|
 \le C_AU\tau(k)^{C_0}(\log x)^{-A},
 \label{eq:current-aux-sw}
\end{equation}
with fixed $C_0$ and uniform family constants.  Let
$|\beta_m|\le\tau(m)^B$, $m\in[V,2V]$, and let each $\gamma_i$ be a
separate one-bounded sequence on $q_i\asymp Q_i$, supported on
$(q_i,\mathfrak a)=1$.  For
\[
 \Delta_{\ell}=\sum_{n,m}\alpha_n\beta_m
 \left({\one}_{nm\equiv\mathfrak a\pmod \ell}
       -\frac{{\one}_{(nm,\ell)=1}}{\varphi(\ell)}\right),
\]
the stated classical inputs of Appendix~\ref{app:classical} give
\begin{equation}
 \sum_{q_1,q_2,q_3}\gamma_1(q_1)\gamma_2(q_2)\gamma_3(q_3)
                    \Delta_{q_1q_2q_3}
       =O_A\bigl(x(\log x)^{-A}\bigr).
 \label{eq:current-medium-conclusion}
\end{equation}
The constant is uniform in the residue and families satisfying these
conditions.  Fixed endpoint multiples are absorbed by the strict
power margins.
\end{proposition}

\paragraph{Application of the local estimate.}
Fix one finite admitted coefficient family. Appendix~\ref{app:coefficients}
uses its full and selected-complement square conditions to construct three
separate modulus factors on whole atom boxes. Appendix~\ref{app:prime-products}
verifies the remaining hypotheses of Proposition~\ref{prop:current-medium}
and applies it to the finite prime-product families. The
\hyperlink{coef:medium-checks}{hypothesis checks in Appendix~\ref{app:coefficients}}
identify these steps. Low-rank pieces and the branches without a medium slot
have their separate estimates. The resulting common input for the sieve
constructions is the prime-product distribution estimate in the next subsection.

The three modulus coefficients in
Proposition~\ref{prop:current-medium} remain separate, and the estimate
allows their signs. In the application the medium slot is a rough
$1$ or M\"obius sequence on an individual chart, with a permitted
Mellin twist. Appendix~\ref{app:prime-products} proves the auxiliary
Siegel--Walfisz condition~\eqref{eq:current-aux-sw} for these sequences.

When the residue has size comparable to $x$, reciprocity leaves an
oscillatory phase. Appendix~\ref{app:medium} keeps this phase and bounds
its variation within the strict power margins
of~\eqref{eq:current-medium-ranges}. The spectral estimates, including
their Fourier normalizations, are stated in
Appendix~\ref{app:classical}.

\subsection{From the local estimate to prime products}

For a fixed finite smooth weight $W_N$ on $k$ labelled original primes,
all at least $cN^{a_*}$ for fixed $a_*>0$, put
\[
 X_W^{(q)}=\sum_{(p_1\cdots p_k,q)=1}W_N(p_1,\ldots,p_k),
 \qquad X_W=\sum W_N(p_1,\ldots,p_k).
\]
The product is $O(N)$.  The allowed geometric families have fixed
polylogarithmic mixed Euler-derivative bounds and no additional
arithmetic filter.  The conclusion used throughout the paper is
\begin{equation}
 \sum_q\rho_N(q)
 \left\{\sum_{p_1\cdots p_k\equiv N\pmod q}W_N(p_1,\ldots,p_k)
            -\frac{X_W^{(q)}}{\varphi(q)}\right\}
       =O_{A,W}\bigl(N(\log N)^{-A}\bigr).
 \label{eq:current-PP}
\end{equation}
The same signed coefficient $\rho_N$ occurs in both terms.  Its
admitted classes are the actual upper tree, the actual lower trees,
the marked queries and ordinary descendants of the finite $H/h$
construction, and the positive switched families on which the upper
coefficient is used.

The coefficient proof in Appendix~\ref{app:coefficients} first freezes
one finite causal policy.  The physical level $N^\lambda$ determines
its density and its residual $Q/B$.  A slightly larger auxiliary
support cover $N^{\ell_{\rm aux}}$ may be used to prove its
factorization; it does not change $\lambda$ or the integer coefficient.
For an ordered prime prefix write $t_i=\log_Np_i$ and
$S_j=\sum_{i=1}^jt_i$, and let $I$ be the subset of slots selected
for one factor of the modulus. The full-square conditions
\[
 S_{j-1}+2t_j\le\ell_{\rm aux}
\]
and every selected-complement square
\[
 \sum_{i\in I}t_i+
       \sum_{\substack{i<j\\i\notin I}}t_i+2t_j
       \le\ell_{\rm aux}\quad(j\notin I)
\]
are retained. After all selected slots have appeared, the second test
becomes the full-square test.  Descendant closure and the physical
prime ceiling therefore persist through the finite ordinary queries.

There is an exact tensor identity before estimating discrepancies.
For a common prime-label alphabet with multiplicities $m_I$, split
the slots into three groups and divide each group convolution by
$M_j=\prod_I m_{I,j}!$.  The resulting $\gamma_j$ is one-bounded.
With $M_{\rm all}=\prod_I m_I!$, the full normalized convolution is
\[
 \frac{M_1M_2M_3}{M_{\rm all}}
                       (\gamma_1*\gamma_2*\gamma_3).
\]
Marked multiplicities remain in the external scalar.  The small-prime
beta factor attaches by the unique small/rough factorization.  The
partition is fixed on a whole atom box before any arithmetic
discrepancy is inspected.  Global squarefreeness remains in the
physical coefficient; its removal only for analytic comparison has
the fixed-power error
$O((N+N^{\theta_*})N^{-\eta_0+o(1)})$, where $N^{\eta_0}$ is
the small-prime presieve cutoff and $\theta_*<3/5$ is the support
exponent of the finite coefficient construction.

Appendix~\ref{app:prime-products} makes a smooth complete-product trim
below $2N^{1-\kappa_0}$ and then an exact smooth dyadic partition.
The local scale is the complete product, with
\[
 cN^{1-\kappa_0}\le x\le CN,\qquad
 c_0x\le N\le x^{1+\kappa}.
\]
It is never a residual quotient.  On an original group at scale $R$,
the exact order-three Heath--Brown identity is obtained from
$\Lambda*(\delta-1*\mu_Z)^{*3}=0$ below $2R$, with
$Z=(2R)^{1/3}$:
\begin{equation}
 \Lambda=\sum_{j=1}^3(-1)^{j-1}\binom3j
                 \mu_Z^{*j}*1^{*(j-1)}*\log.
 \label{eq:current-HB}
\end{equation}
The smooth original chart is retained; no sharp coupled group boundary
is inserted.  The Fourier separation of the geometric profile is
charged by its integrable Fourier norm.  The actual rough slots have
the auxiliary SW estimate.  Medium branches use
Proposition~\ref{prop:current-medium}; low ranks use below-half
Bombieri--Vinogradov.  The one-long and two-long alternatives are the
direct Brun arguments of Appendix~\ref{app:brun}, keeping the exact
coprime mean and both Poisson zero-axis cancellations.  In particular,
negative lower Brun polynomials are not multiplied as though they
were positive.

Choose the fixed moment order $K_{\rm mom}$ before narrowing
\[
 0<\kappa\le\min\left(\kappa_*,\frac{\sigma}{200},
                           \frac1{2K_{\rm mom}}\right),\qquad
 0<\kappa_0<\frac{\kappa}{4(1+\kappa)}.
\]
Here $\kappa_*>0$ is fixed by the strict support margins in
Appendix~\ref{app:coefficients}. The complete local support has exponent $\Theta_{\rm loc}<3/5$;
the long-branch parameter $\zeta$ is selected with
$20\zeta+\Theta_{\rm loc}<1$ and
$40\zeta+3\Theta_{\rm loc}/2<1$.  All finite atom scalar masses and
the fixed logarithmic number of charts are paid by stronger local
logarithmic accuracy.  This proves \eqref{eq:current-PP} for the
admissible families and all sufficiently large even targets.  

\section{A positive upper sieve with its own density}
\label{sec:current-upper}

The upper construction uses
\[
 \rho_0=\frac{19}{96},\quad \mu_0=\frac{25}{128},\quad
 \Theta=\frac{299}{500},\quad \Lambda_{\rm cap}=\frac{153}{256},\quad
 v(t)=\min\left(\frac{1+t}{2},\frac{57}{64}-\frac{3t}{2}\right).
\]
Its causal cap changes are made on the actual prefix; every retained
ancestor decision is preserved.  Fourth, fifth and sixth changes are
successive replacements, whose raw density differences telescope.
Later eighth-rank changes occupy $30+15+13+29=87$ specified domains.
Same-stage domains have disjoint interiors; successive stages use
differences from their respective parent caps.  
Appendix~\ref{app:upper} supplies the cap formulas and separation proofs.

\subsection{Exit identity and literal residual densities}

For a finite allowed prime universe, let $\xi$ be multiplicative with
$\xi(1)=1$ and $0\le \xi(p)\le1$. The upper
tree, whose first exits have odd rank, satisfies the exact identity
\begin{equation}
 \sum_d\mu(d){\one}_{d\ {\rm retained}}\xi(d)
 =\prod_p(1-\xi(p))
   +\sum_{e\ {\rm odd\ first\ exit}}\xi(e)
                      \prod_{p<p_{\rm last}(e)}(1-\xi(p)).
 \label{eq:current-upper-exit}
\end{equation}
Partition the excluded divisors by their first excluded prefix; the
remaining smaller-prime subsets give the displayed Euler product.
Taking $\xi(p)={\one}_{p\mid n}$ proves a rough-indicator majorant for
every integer $n$, including zero and negative integers.

For the raw density, take
$g_N(p)={\one}_{p>2,\,p\nmid N}/(p-1)$ and
$V_N(y)=\prod_{2<p<y,\,p\nmid N}(1-1/(p-1))$.
At a fixed rank the normalized exit for $g_N$ lies between zero and
the corresponding exit for $g_0(p)={\one}_{p>2}/(p-1)$: a surviving
prefix acquires the multiplier
\[
 \prod_{\substack{p_{\rm last}\le p<y\\p\mid N}}(1-g_0(p))\le1.
\]
The preceding square tests bound the last prime of a rank-$j$ exit
below by $N^{\beta_{\min}/(3\cdot2^{j-1})}$ on positive compact
parameter sets.  Deleting target primes thus changes each fixed-rank
integral by $o(1)$ uniformly in $N$.  Positive exit-series convergence,
the one-sided bounds for the same raw Rosser weights, and then a
finite-rank sandwich yield
\begin{equation}
 \frac{V_N^+(R,y)}{V_N(y)}=F(s)+o(1)\quad(s>1),\qquad
 \frac{V_N^-(R,y)}{V_N(y)}=f(s)+o(1)\quad(s\ge2),
 \quad s=\frac{\log R}{\log y}.
 \label{eq:current-raw-density}
\end{equation}
The upper identity for $1<s\le3$ also follows directly by moving its
first cube to $R^{1/3}$.  This computes the density of the divisor coefficient itself.  The ordinary exit-series inputs and their normalization are
specified in Appendix~\ref{app:upper}.

Mertens gives, uniformly in the target,
\begin{equation}
 V_N(N^u)=\frac{2e^{-\gamma}C_N}{u\log N}(1+o(1)).
 \label{eq:current-mertens}
\end{equation}
For example the exact ratio to $2C_N\prod_{p<z}(1-1/p)$ is
\[
 \prod_{p\ge z}\left(1-\frac1{(p-1)^2}\right)^{-1}
 \prod_{\substack{p\mid N\\p\ge z}}\left(1-\frac1{p-1}\right),
\]
whose logarithm tends uniformly to zero at $z=N^u$, $u>0$ fixed.
With $\overline F=F/(2e^\gamma)$ and
$\overline f=f/(2e^\gamma)$, this is the normalization producing
$4\overline F(s)/u$ and $4\overline f(s)/u$ in the density functionals.

For a selected eighth cell, if $S_8=t_1+\cdots+t_8$ and the parent is bounded above by $R_6$, its cap $M$ satisfies
\[
 1<\frac{\lambda_{6,\rm parent}-S_8}{t_8}
  \le\frac{R_6-S_8}{t_8}<\frac{M-S_8}{t_8}<3.
\]
The same residual identity therefore gives the lower saving
\begin{equation}
 \frac{4(M-R_6)}{(R_6-S_8)(M-S_8)\prod_{i=1}^8t_i}.
 \label{eq:current-eighth-saving}
\end{equation}
Using the larger parent $R_6$ is conservative.  This density has no
Chen quarter or ordering factorial.  The selected domains have
$t_8>53/13824$, and all full and selected-complement squares remain
in force.

\subsection{Finite realization and density errors}

The ideal policy is approximated on a finite causal partition before
$N$ grows.  Small-last-prime tails tend to zero; proper affine walls
have vanishing mass, whereas open inherit/equality regions are
retained.  At even rough rank take outer endpoint tests and at odd
rank inner tests, preserving ancestry.  This produces a coefficient
$a_N(d)$ above the raw upper coefficient.  Attach the small-prime
beta upper polynomial to positive $a_N(d)$ and the beta lower
polynomial to negative $a_N(d)$.  If $B^-\le I_w\le B^+$, then
$a_N(d)B_{\rm sign}(n)\ge a_N(d)I_w(n)$ term by term.
The resulting physical $\rho_N^{\rm up}$ has unique small/rough decomposition,
squarefree support $q\le N^{\theta_*}$, a uniform bound
$|\rho_N^{\rm up}(q)|\le M$ for fixed $M$, $(q,N)=1$, and
\begin{equation}
 R_N(n):=\sum_{q\mid n}\rho_N^{\rm up}(q)
   \ge {\one}_{(n,P_N(N^{\rho_0}))=1}\ge0,
 \qquad P_N(y)=\prod_{\substack{p<y\\p\nmid N}}p.
 \label{eq:current-physical-upper}
\end{equation}
The finite sum also defines $R_N(0)$.

Restricting \eqref{eq:current-upper-exit} to primes above the presieve
gives the upper projection inequality
\[
 V_N(w)\sum_{d\ {\rm rough}}\lambda_N(d)g_N(d)
       \le\sum_{d\ {\rm full}}\lambda_N(d)g_N(d).
\]
Together with the absolute beta and box-error estimates this controls
the density of the same $\rho_N^{\rm up}$ from above.  Support inflation can
be kept below $1/1000$, giving $\theta_*=599/1000$ before the small
complete-scale adjustment.  The auxiliary support cover used in
Section~\ref{sec:current-distribution} does not change this density.

Let $E_{\rm up}$ be the raw upper-density endpoint and let $u$ be each of the
two allowances for finite approximation and coefficient realization.
Their exact rational values are given in
\eqref{eq:current-E}--\eqref{eq:current-u}. The directed estimates in
Appendix~\ref{app:numeric} give
\begin{equation}
                     K-E_{\rm up}-2u>9.6522\cdot10^{-8}>0.
                     \label{eq:current-upper-slack}
\end{equation}
Thus the physical coefficient in \eqref{eq:current-physical-upper} satisfies
\begin{equation}
 S_N:=\sum_q\frac{\rho_N^{\rm up}(q)}{\varphi(q)}
       \le(K+o(1))\frac{C_N}{\log N},
 \label{eq:current-upper-density}
\end{equation}
and it is precisely the coefficient to which
\eqref{eq:current-PP} applies.

\begin{proof}[Proof of Theorem~\ref{thm:positive-coefficient}]
Use the finite coefficient constructed above.
Equation~\eqref{eq:current-physical-upper} gives its bounded squarefree
support and~\eqref{eq:core-positivity}, and
\eqref{eq:current-upper-density} gives its density bound.
The factorization and parameter choices in
Appendix~\ref{app:coefficients}, followed by the prime-product estimate
in Appendix~\ref{app:prime-products}, establish
\eqref{eq:current-PP} for this coefficient and each weight in the
fixed finite collection. Choose their common refinement before taking
$N$ to infinity. Taking maxima of the finitely many constants and
thresholds gives the asserted uniformity.
\end{proof}

\begin{proposition}[Positive product consequence]
\label{prop:current-positive-product}
Fix a finite nonnegative labelled prime-product family of bounded
rank and uniformly bounded weights, with every original prime at least $N^{a_*}$ for some fixed
$a_*>0$, product $O(N)$, and only the admitted fixed geometric
restrictions.  Suppose its unsifted weighted mass $X_W$ satisfies
\[
 \limsup_{\substack{N\to\infty\\2\mid N}}
             \frac{\log N}{N}X_W\le J<\infty.
\]
Let $W_N(m)$ denote the induced weight on the product $m$. Then
\[
 \limsup_{\substack{N\to\infty\\2\mid N}}\frac1{H_N}
       \sum_{\substack{m<N\\N-m\text{ prime}}}W_N(m)\le KJ.
\]
The error estimate is absolute also when $J=0$.
\end{proposition}

\begin{proof}
For a prime $\ell$, squarefree support and boundedness give
\[
 \sum_{q:\ell\mid q}\frac{|\rho_N^{\rm up}(q)|}{\varphi(q)}
       \ll\frac{\log N}{\ell-1}.
\]
Hence, by a union bound over the fixed number of original prime slots,
\begin{equation}
 \left|\sum_q\frac{\rho_N^{\rm up}(q)}{\varphi(q)}
                         (X_W^{(q)}-X_W)\right|
       \ll X_WN^{-a_*}\log N=o(H_N).
 \label{eq:current-coprime-correction}
\end{equation}
No factor $N^{\theta_*}$ is introduced.  Small prime remainders cost
$N^{\rho_0+o(1)}$.  The smooth complete-product trim costs
$N^{1-\kappa_0+o(1)}$ after the full modulus sum; its product is
strictly below $N$, so no $\tau(0)$ occurs.  Positive majorization
before distribution gives
\[
 \sum_{m<N,\,N-m\ {\rm prime}}W_N(m)
 \le\sum_m W_N(m)R_N(N-m)+o(H_N)
 \le X_WS_N+o(H_N).
\]
Outer approximations with $m\ge N$ are harmless because
$R_N(n)\ge0$ for every integer.  Now apply
\eqref{eq:current-upper-density}.  Fixed geometric boundaries are
handled by positive smooth outer approximation, with $N$ tending to
infinity before the approximation error tends to zero.  At each
precision the same $\rho_N^{\rm up}$ is used in the pointwise inequality,
density and distribution formula.  It need not be the same
coefficient at different precisions.
\end{proof}

\section{The full Chen weight and paired deletion}
\label{sec:current-chen-weight}

We use $a_{\rm s},h_{\rm s},d_{\rm s}$ for the fixed sieve cutoffs;
the parameter $a$ in Proposition $(1+a)$ retains its introductory meaning. Put
\[
 {a_{\rm s}}=\frac{417}{5000},\qquad {h_{\rm s}}=\frac{19}{125},\qquad
 {d_{\rm s}}=\frac{59}{250},\qquad \mu_0=\frac{25}{128}.
\]
For $n\le N$ define
$Q_d(y;n)=\one_{d\mid n}\one_{P^-(n/d)\ge y}$.
All prime prefixes below have distinct entries.  Suppressing the
argument $n$, the eleven terms are
\begingroup\allowdisplaybreaks[4]
\begin{align*}
 Y_1&=Q_1(N^{a_{\rm s}}),\\
 Y_2&=Q_1(N^{h_{\rm s}}),\\
 Y_3&=\sum_{N^{a_{\rm s}}\le p<N^{1/3}}Q_p(N^{a_{\rm s}}),\\
 Y_4&=\sum_{N^{a_{\rm s}}\le p<N^{d_{\rm s}}}Q_p(N^{a_{\rm s}}),\\
 Y_5&=\sum_{N^{a_{\rm s}}\le p_1<p_2<N^{h_{\rm s}}}Q_{p_1p_2}(N^{a_{\rm s}}),\\
 Y_6&=\sum_{N^{a_{\rm s}}\le p_1<N^{h_{\rm s}}\le p_2<N^{d_{\rm s}}}Q_{p_1p_2}(N^{a_{\rm s}}),\\
 Y_7&=\sum_{N^{d_{\rm s}}\le p_1<p_2<(N/p_1)^{1/2}}Q_{p_1p_2}(p_2),\\
 Y_8&=\sum_{N^{a_{\rm s}}\le p_1<N^{1/3}\le p_2<(N/p_1)^{1/2}}
                                          Q_{p_1p_2}(p_2),\\
 Y_9&=\sum_{N^{h_{\rm s}}\le p_1<N^{d_{\rm s}}\le p_2<(N/p_1)^{1/2}}
                   Q_{p_1p_2}\bigl((N/(p_1p_2))^{1/2}\bigr),\\
 Y_{10}&=\sum_{N^{a_{\rm s}}\le p_1<p_2<p_3<p_4<N^{h_{\rm s}}}Q_{p_1p_2p_3p_4}(p_2),\\
 Y_{11}&=\sum_{N^{a_{\rm s}}\le p_1<p_2<p_3<N^{h_{\rm s}}\le p_4<N^{d_{\rm s}}}
                                          Q_{p_1p_2p_3p_4}(p_2).
\end{align*}
\endgroup
In particular $Y_{11}$ has its full last range.  Set
\begin{equation}
 W=3Y_1+Y_2-Y_3-Y_4+Y_5+Y_6-2Y_7-Y_8-Y_9-Y_{10}-Y_{11}.
 \label{eq:current-full-weight}
\end{equation}

The terms $3Y_1+Y_2$ supply the rough-count baseline at the two
cutoffs. The one-prefix subtractions $Y_3,Y_4$ remove overcounting,
while $Y_5,Y_6$ restore the selected two-prefix contributions.
The terms $2Y_7,Y_8,Y_9$ subtract the remaining three-or-more-factor
patterns, and the four-prefix terms $Y_{10},Y_{11}$ compensate the
growth of the positive pair counts when several small factors occur.
Their combination gives the pointwise inequality below. The later
paired deletion links $Y_6$ and $Y_{11}$ by an injection, so that the
change preserves that inequality.

\begin{lemma}[Pointwise full-weight inequality]
\label{lem:current-full-weight}
For squarefree $n\le N$,
\[
                  W(n)\le4{\one}_{\Omega(n)\le2}.
\]
\end{lemma}

\begin{proof}
If $n$ is not $N^{a_{\rm s}}$-rough, all eleven terms vanish.  For $Y_9$, the
only less immediate case, its quotient cutoff exceeds
$N^{(1-{d_{\rm s}})/4}>N^{a_{\rm s}}$.  Suppose therefore that $n$ is squarefree and
$N^{a_{\rm s}}$-rough.  Let $r,s,t$ count its prime factors below $N^{h_{\rm s}},N^{d_{\rm s}},N^{1/3}$,
respectively.  Its first six counts are
\[
 (Y_1,Y_2,Y_3,Y_4,Y_5,Y_6)
      =\left(1,{\one}_{r=0},t,s,\binom r2,r(s-r)\right).
\]
If $r\ge2$, the quotient sifted at the second selected prime forces
the first two primes in $Y_{10}$ and $Y_{11}$ to be the two least
prime factors of $n$.  Thus
\[
 Y_{10}=\binom{r-2}{2},\qquad Y_{11}=(r-2)(s-r),
\]
with the binomial zero at $r=2,3$.  Substitution in the first six
weighted terms minus $Y_{10}+Y_{11}$ gives $s-t\le0$; the other
subtractions help.

For $k=\Omega(n)\ge3$, let $q_1<q_2<\cdots$ be the factors of $n$.
Then $q_1q_2^2<n\le N$, so $t\ge1$,
$Y_7={\one}_{s=0}$ and $Y_8={\one}_{t=1}$.
When $r\le1$, the remaining upper bounds for $W$ are
\[
\begin{array}{c|c}
 \text{case}&\text{upper bound}\\ \hline
 r=1&2-t-{\one}_{t=1}\\
 r=0,\ s\ge2&4-s-t\\
 r=s=0&2-t-{\one}_{t=1}\\
 r=0,\ s=1,\ t\ge3&3-t\\
 r=0,\ s=1,\ t\le2&3-t-1-{\one}_{t=1}.
\end{array}
\]
In the last case, $k\ge4$ would imply
$n\ge N^{{h_{\rm s}}+{d_{\rm s}}+2/3}>N$, so $k=3$; moreover
$q_1q_2q_3^2\ge N^{{h_{\rm s}}+{d_{\rm s}}+2/3}>N$.  The pair $(q_1,q_2)$ therefore
contributes to $Y_9$, giving the extra subtraction displayed there.
All the bounds are nonpositive.  For $k\le2$ direct substitution
gives $W\le4$; the unit has $W(1)=4$.
\end{proof}

Let $\mathscr E=\mathscr E_0\cup \mathscr E_1$ be the disjoint union of the rational rectangles
specified in Appendix~\ref{app:numeric}. The two families have $6091$ and
$11481$ members, respectively, in the coordinates
\[
 (t,u)=(\log_Np_2,\log_Np_1),\qquad {a_{\rm s}}\le u<{h_{\rm s}}\le t<{d_{\rm s}}.
\]
The rectangles have disjoint interiors; shared edges have zero limiting
mass and are assigned once in the finite partition.
Write $Y_6(\mathscr E)$
for its pair terms and $Y_{11}(\mathscr E)$ for the four-prefix terms whose
last two selected primes $(p_3,p_4)$ have
$(\log_Np_4,\log_Np_3)\in \mathscr E$.  The weight is
\begin{equation}
                         W_{\mathscr E}=W-Y_6(\mathscr E)+Y_{11}(\mathscr E).
 \label{eq:current-paired-weight}
\end{equation}

\begin{lemma}[Paired injection]
\label{lem:current-paired-injection}
For squarefree $n\le N$, $Y_{11}(\mathscr E)\le Y_6(\mathscr E)$, and hence
$W_{\mathscr E}(n)\le4{\one}_{\Omega(n)\le2}$.
\end{lemma}

\begin{proof}
Map a contributing four-prefix tuple to its last pair $(p_3,p_4)$.
The quotient roughness forces its first two selected primes to be
the two least factors of $n$. The image lies in the prescribed
rectangle, and removing $p_3p_4$ leaves an $N^{a_{\rm s}}$-rough residual,
so this pair contributes to $Y_6(\mathscr E)$. Since the first two factors
are determined by $n$, two tuples with the same last pair coincide.
The map is therefore injective.
\end{proof}

The fixed prime-factor bound $\lfloor1/{a_{\rm s}}\rfloor=11$ bounds
$|W_{\mathscr E}|$ by $1016$ on rough $n$.  The number of $n\le N$ with a
repeated prime at least $N^{a_{\rm s}}$ is at most $2N^{1-{a_{\rm s}}}$ for large $N$.
After restricting to primes $p<N$ with $(p,N)=1$ and removing the
unit, the squareful and unit costs are at most
$2032N^{1-{a_{\rm s}}}+4=o(H_N)$.  It follows that
\begin{equation}
 4D_{1,2}(N)\ge
       \sum_{\substack{p<N\\p\ {\rm prime}\\(p,N)=1}}W_{\mathscr E}(N-p)
        -o(H_N).
 \label{eq:current-weight-to-count}
\end{equation}

\section{The lower sieve and finite double-sieve transitions}
\label{sec:current-lower}

\subsection{A lower estimate for the original prime sequence}

The ordinary lower trees start at $w={a_{\rm s}}$ and $w={h_{\rm s}}$, with cap
$\nu=249/448$ and fixed small buffers.  Their first exits have even
rank.  The analogue of \eqref{eq:current-upper-exit} is
\begin{equation}
 \sum_d\mu(d){\one}_{d\ {\rm retained}}\xi(d)
 =\prod_p(1-\xi(p))
  -\sum_{e\ {\rm even\ first\ exit}}\xi(e)
                       \prod_{p<p_{\rm last}(e)}(1-\xi(p)).
 \label{eq:current-lower-exit}
\end{equation}
This identity proves both the pointwise lower inequality and the
literal raw-density formula. The selected fourth-rank restorations
preserve even first exits: each cap is chosen before the fourth cube
test $S_3+3t_4<\lambda_4$, and the ordinary residual level exceeds
$p_4^2$. A restored slab and a fifth-rank saving are
assigned disjoint regions in the lower functional.

For a presieve $z_0=N^{\eta_0}$, even exits give the exact projection
direction
\begin{equation}
 V_N(z_0)R_{\rm rough}-V_{\rm full}
  =\sum_{\substack{e\ {\rm full\ even\ exit}\\p_{\rm last}(e)<z_0}}
                              g_N(e)V_N(p_{\rm last}(e))\ge0.
 \label{eq:current-lower-projection}
\end{equation}
Round positive even coefficients inward and negative odd coefficients
outward, retaining the prefix decisions.  Attach beta-minus to positive
coefficients and beta-plus to negative coefficients.  If the rounded
coefficient is $a_{\rm rnd}(d)$, then
$a_{\rm rnd}(d)B_{\rm sign}(n)\le a_{\rm rnd}(d)I_0(n)$ term by term.  Therefore the
physical lower coefficient satisfies
\[
 \sum_{q\mid n}\rho_N^{\rm low}(q)\le
           {\one}_{(n,P_N(N^w))=1},\qquad
 \sum_q\frac{\rho_N^{\rm low}(q)}{\varphi(q)}
              \ge V_{\rm full}-\text{coefficient-realization error}.
\]
Beta-minus need not be nonnegative.  The second inequality uses
\eqref{eq:current-lower-projection} and absolute error control, not
coefficient order alone.  The coefficient theorem supplies
\eqref{eq:current-PP} for this very lower coefficient.

The original prime sequence has a sharp endpoint at $N$. To use the
smooth distribution estimate for a signed lower coefficient, we first
control the endpoint error in absolute value.

\begin{lemma}[Sharp original-prime endpoint]
\label{lem:current-sharp-prime}
Suppose the actual coefficient satisfies $|\rho_N^{\rm low}(q)|\le M$ with
fixed $M$, is supported on squarefree $q\le N^\theta$ with
$\theta<1$ and $(q,N)=1$, and satisfies the permitted
smooth-singleton instance of \eqref{eq:current-PP}.  Then
\begin{equation}
 \sum_q\rho_N^{\rm low}(q)
    \left(\pi(N;q,N)-\frac{\pi(N)}{\varphi(q)}\right)
      =O_A\bigl(N(\log N)^{-A}\bigr).
 \label{eq:current-sharp-pi}
\end{equation}
\end{lemma}

\begin{proof}
Smooth the endpoint over $[N-\delta N,N+\delta N]$, with
$\delta=(\log N)^{-B}$ and $B>A+3$, and apply the smooth singleton
estimate. The absolute boundary and exact-coprime mean estimates in
Appendix~\ref{app:cross-review} bound the difference by
$O(\delta N\log N+\log^2N)$, together with the fixed-power complete
product trim. The equality $p=N$ is impossible for even $N>2$.
These errors are absorbed by the chosen logarithmic accuracy.
\end{proof}

Apply the pointwise lower coefficient to the original prime sequence.
Lemma~\ref{lem:current-sharp-prime} bounds its remainder, while
\eqref{eq:current-raw-density} and the prime number theorem evaluate
its main term. This gives the ordinary lower prime counts used below.

\subsection{Finite defects and actual descendants}

Put
$A_+(s)=sF(s)/(2e^\gamma)$ and
$A_-(s)=sf(s)/(2e^\gamma)$.
For a nonnegative finite marked family $\sigma(B)$ and one common
$Q=N^\lambda$, normalize by
\[
 \Theta_Q=4\operatorname{li}(N)
       \sum_B\frac{\sigma(B)C_{BN}}{\varphi(B)\log(Q/B)}.
\]
Here $C_{BN}$ is the same singular-factor function as in
\eqref{eq:current-normalization}.  A defect $H(s)$ means an upper
bound $(A_+(s)-H(s)+o(1))\Theta_Q$; a defect $h(s)$ means a lower
bound $(A_-(s)+h(s)+o(1))\Theta_Q$.  These are estimates
averaged over the positive marked family, not pointwise in $B$ and
not assertions for a signed outer measure.

For residual exponent $L=\lambda-\log_NB$, write
$\alpha=(1-\log_NB)/L\ge2$.  A singleton child with exponent $Lu$
has
\[
 L'=L(1-u),\qquad \alpha'=1+\frac{\alpha-1}{1-u};
\]
a two-prime child replaces $u$ by $u+v$.  All children retain the
same total $Q$, every full and selected square, and their physical
prime ceiling.  The parameter $\alpha$ increases and is never
reset to $2$ in a switched mass.

Lemma~\ref{sv:finite-seed} proves the initial defect $J_0=1/50$
and the positive transition
$J\mapsto0.020249+0.13J$ on the permitted marked states with
$1\le s\le11/5$. Its hypotheses retain the residual level, the
prefix square conditions and a common comparison cutoff. Four
applications give $J_4=0.023273777353>2327/100000$; this is the seed
for the finite recurrence below.
The finite collection consists of $70$ one-step and $36$ two-step
inequalities, together with the positive Buchstab transitions
\[
 H_{\rm new}(s)=\int_{\max(2,s-1)}^5\frac{h(w)}w\,dw,
 \qquad
 h_{\rm new}(s)=\int_{s-1}^5\frac{H(w)}w\,dw\quad(s\ge2).
\]
The one-step switched cost is $KI/8$, where $I$ is the
nonnegative switched mass of the action specified in
Appendix~\ref{app:lower}.  A two-step action has the
positive defect contribution
\begin{equation}
 \frac15\left[4H(s_0)+H(\phi_1)
  +\sum_{\mathcal D_1}\int\frac{h(\rho)}{u(1-u)}\,du
  +\sum_{\mathcal D_2}\iint\frac{H(\rho)}{uv(1-u-v)}\,du\,dv\right]
 \label{eq:current-w2-feedback}
\end{equation}
in addition to its proved ordinary and positive-switching base.
There are three singleton and four good pair domains.  With
$a_i=1/\phi_i$ and $a_0=1/s_0$, the pair domains are
\[
 \begin{array}{c|c|c}
 &\text{domain}&\rho_i(u,v)\\ \hline
 D_5^+&a_0<u<v<a_2,\ v+2u<1&s_0(1-u-v)\\
 D_6&(a_0,a_1)\mathbin{\times}(a_2,a_3)&s_0(1-u-v)\\
 D_7&a_0<u<v<a_1&(1-u-v)/u\\
 D_8&(a_0,a_1)\mathbin{\times}(a_1,a_2)&(1-u-v)/u.
 \end{array}
\]
The admissibility conditions include $2v<1$ and $v+2u<1$.
The excluded portion of $D_5$ receives only a positive product cost,
not ordinary feedback.  Unit surfaces in the switching terms stay
in their masses.  The envelope used for every actual $\alpha\ge2$ is
$\omega(s)\le\max(1/s,70893/125000)$ for $s\ge1$, with the
conservative extension below $1$ specified in Appendix~\ref{app:lower}.

\paragraph{Two levels of iteration.}
The four scalar steps produce the certified constant used to initialize
the arrays; they precede generation zero of the array recurrence.
The indices in $H_m,h_m$, including $m=248$ and $430$, count only
simultaneous updates of the two cell arrays from that initialization.
Thus the scalar seed iteration and the envelope generations have
different state spaces and iteration counts.

We define the envelopes by simultaneous finite induction. Each
right-hand side uses the preceding generation, takes the maximum of
alternative estimates, and uses a minimum over the full closed support
when an argument crosses a cell boundary. The $200$ closed cells have
endpoints $1+i/50$, $0\le i\le200$. Let $H_m,h_m$ be the step-function
lower envelopes after $m$ generations; at a shared endpoint take the
smaller of the adjacent values. Starting directly from the seed,
$430$ rounds determine all $400$ final entries. These
generation-$430$ arrays are used only in $U_3,U_4$; the remaining
occurrences use generation $248$.

The finite induction also needs a common approximation, not just
finite depth.  Positive transition mass is at most $12$ and the
root coefficient norm at most $7/2$.  A local oriented error
$\delta$ propagates by at most
$\delta(12^{441}-1)/11$ before normalization.  Choose
$\delta$ inside the prescribed final reserve and use one nested
finite refinement for all queries.  Residual exponents after the
finite transitions remain above a fixed positive multiple of
$101^{-440}$, so every positive leaf satisfies the fixed-cutoff
hypothesis of Proposition~\ref{prop:current-positive-product}.
Zero-mass cells are controlled by absolute errors.

\subsection{Evaluation of the eleven-term weight}

The lower functional is
\begin{equation}
 \mathcal C=\frac14
   (3L_1+L_2+L_5+L_6-U_3-U_4-2U_7-U_8-U_9-U_{10}-U_{11}).
 \label{eq:current-root-functional}
\end{equation}
The lower starts are the causal starts at ${a_{\rm s}},{h_{\rm s}}$. The ordinary estimate
at ${h_{\rm s}}$ is inserted after its full least-prime correction is removed
from the reference expression. In $U_3,U_4$ the
$H_{430}$ argument is $\max(2,L/{a_{\rm s}})$ with $L=v(t)-t$ in the enhanced
range. The region $L/{a_{\rm s}}\le2$ has positive measure and uses the value
of the lower envelope at $2$.
The pair lower terms use $h_{248}$ where $L/{a_{\rm s}}\ge2$ and zero below
$2$.  The positive costs $U_7,U_8,U_9$ retain their labelled
multiplicities and are estimated by
Proposition~\ref{prop:current-positive-product}.

For $U_{10},U_{11}$, work on one common positive partition of the
four-prefix domain, with decreasing coordinates $x>y>z>w$.
At each cell use the minimum of the switched kernel
$K\omega((1-S_4)/z)/z$ and the admissible ordinary kernels
$4(1-H_{248}(\max(2,L/w)))/L$. Each ordinary alternative retains
all full and selected square conditions. A refined estimate replaces
its parent estimate, and overlapping alternatives are compared rather
than added. Appendix~\ref{app:lower} gives the states and their common
prefix-preserving approximation.

The restriction in~\eqref{eq:current-paired-weight} requires an estimate
of the lower functional as well as the pointwise injection. On a pair
rectangle $\mathscr B=[t_l,t_r]\times[u_l,u_r]$, put
\[
 R(t,u)=\frac{57}{64}-\frac52t-u,\qquad
 R_{\max}=R(t_l,u_l),\quad s_{\max}=R_{\max}/{a_{\rm s}},\quad
 M_{\mathscr B}=\log(t_r/t_l)\log(u_r/u_l).
\]
Write $\mathcal L_{\mathscr B}$ for the cellwise lower estimator assigned to a
rectangle $\mathscr B$. It vanishes on $\mathscr E_0$. In $\mathscr E_1$ there are
$227$ crossing cells on which it vanishes and $11254$
positive cells with $2\le s_{\min}\le s_{\max}\le9/4$.
Since $\log(s-1)/s$ is increasing there, and $h_{\mathscr B}^{\min}$ denotes the
cellwise $h_{248}$ minimum over the full closed argument support,
the discarded lower estimator, in Chen-quarter units, is
bounded above by
\begin{equation}
 \sum_{\mathscr B\in \mathscr E_1\ {\rm positive}}
       \frac{M_{\mathscr B}}{R_{\max}}\bigl(\log(s_{\max}-1)+h_{\mathscr B}^{\min}\bigr)
       \le 0.002372094.
 \label{eq:current-crop-loss}
\end{equation}

For the negative term, the same selected parents are switched in
the $48$-division four-prefix cover and avoid the refinement on
$40000$ parent cells.
If $A_{\mathscr B}$ is the reference density mass and $s_{\mathscr B}$ its fixed $K$
feedback slope, the cellwise bounds in Appendix~\ref{app:numeric} give
$0\le s_{\mathscr B}\le A_{\mathscr B}/K_0$, $K_0=6.7619$.  Therefore
\[
 A'_{\mathscr B}=A_{\mathscr B}-(K_0-K)s_{\mathscr B}\ge(K/K_0)A_{\mathscr B}.
\]
Dividing by its positive reference mass gives an explicit
pointwise density majorant of the switched kernel.
Restricting this majorant to the complementary positive family
removes at least $KI_{\mathscr E}/8$ in Chen units, where
\begin{equation}
 I_{\mathscr E}=\int_{(t,u)\in \mathscr E}\int_{{a_{\rm s}}<w<z<u}
                      \frac{dt\,du\,dz\,dw}{tu z^2w}.
 \label{eq:current-crop-integral}
\end{equation}
The use of $\omega\ge1/2$ is valid since
$(1-S_4)/z\ge77/38>2$ throughout this family.  For one rectangle
the integral is exactly
\[
 \log(t_r/t_l)\left[
 \frac{\log(u_r/u_l)}{a_{\rm s}}+\frac{\log(u_r/{a_{\rm s}})+2}{u_r}
                         -\frac{\log(u_l/{a_{\rm s}})+2}{u_l}\right].
\]
The finite lower bounds for the gross improvements are
$0.002095161$ on $\mathscr E_0$ and $0.003799514$ on $\mathscr E_1$.
Subtracting \eqref{eq:current-crop-loss}, the additional raw
improvement from $\mathscr E_1$ is
\begin{equation}
          0.003799514-0.002372094
             =0.001427420=\frac{71371}{50000000}.
 \label{eq:current-new-crop-gain}
\end{equation}
The paired restriction removes mass from the positive majorant and
pays for the assigned lower estimators on the same deleted region.
The masks $\mathscr E_0,\mathscr E_1$ are disjoint.

\subsection{A preliminary lower functional}

Denote the current lower endpoint by $R_{\rm base}$ and the bound
for each finite-density and coefficient-realization error by
$B_{\rm base}$; see~\eqref{eq:base-R}--\eqref{eq:base-B}.
There are two further improvements on $\mathscr E$: a larger lower bound
for $\omega$ on an inner fibre, and a smaller payment for the
discarded lower estimator. We treat them separately before combining
their contributions.

\subsection{The paired-domain improvement}
\label{sec:current-refinement}

\begin{lemma}[A global Buchstab lower bound]
\label{lem:current-omega-refinement}
For $s\ge1$ one has $\omega(s)\ge1/2$, and for $s\ge7/3$ one has
$\omega(s)\ge11/20$.
\end{lemma}

\begin{proof}
The bound $1/2$ holds on $[1,2]$ and propagates by
$(s(\omega(s)-1/2))'=\omega(s-1)-1/2$.
For $7/3\le s\le3$, use $\log x\ge2(x-1)/(x+1)$ for $x\ge1$ to get
\[
 \omega(s)=\frac{1+\log(s-1)}s\ge\frac{3s-4}{s^2}\ge\frac{11}{20}.
\]
The last inequality follows because the concave polynomial
$60s-80-11s^2$ is positive at both endpoints.
For $3\le s\le10/3$, the first bound and $\log2\ge2/3$ give
\[
 s\omega(s)\ge1+\log2+(s-3)/2\ge s/2+1/6\ge11s/20.
\]
This establishes the stronger bound on the unit interval
$[7/3,10/3]$. The identity
$(s(\omega(s)-11/20))'=\omega(s-1)-11/20$ then propagates it to
successive unit intervals.
\end{proof}

\paragraph{The extra positive fibre.}
The already selected parent majorant dominates
$K\omega((1-t-u-z-w)/z)/z$ on ${a_{\rm s}}<w<z<u<{h_{\rm s}}<t<{d_{\rm s}}$.
For each rectangle $\mathscr B=[t_l,t_r]\times[u_l,u_r]$ of the same mask,
put
\begin{equation}
 Z_{\mathscr B}=\min\left(u_l,\frac{3(1-t_r-u_r)}{13}\right).
 \label{eq:current-Z}
\end{equation}
The finite exact geometry check gives $Z_{\mathscr B}>{a_{\rm s}}$ for every cell.
Moreover the definition gives
$1-t_r-u_r-2Z_{\mathscr B}\ge(7/3)Z_{\mathscr B}$.  Hence on the inner fibre
${a_{\rm s}}<w<z<Z_{\mathscr B}$ one has $z<u$ and
\[
 1-t-u-z-w\ge1-t_r-u_r-2z\ge\frac73z.
\]
Lemma~\ref{lem:current-omega-refinement} supplies an additional
$1/20$ in $\omega$, beyond the $1/2$ already paid in
\eqref{eq:current-crop-integral}.  The remainder of every fibre
keeps its original bound.  The two nonnegative contributions hold simultaneously.

Integrating the added $1/20$ and including the Chen factor $1/4$
gives at least the following additional coefficient from $\mathscr B$:
\begin{align}
 G_{\mathscr B}
 &=\frac{K}{80}
   \log\frac{t_r}{t_l}\log\frac{u_r}{u_l}
      \int_{a_{\rm s}}^{Z_{\mathscr B}}\frac{\log(z/{a_{\rm s}})}{z^2}\,dz \notag\\
 &=\frac{K}{80}
   \log\frac{t_r}{t_l}\log\frac{u_r}{u_l}
   \left[\frac1{a_{\rm s}}-\frac{1+\log(Z_{\mathscr B}/{a_{\rm s}})}{Z_{\mathscr B}}\right].
 \label{eq:current-extra-fibre}
\end{align}
The inner $w$ integral is $\log(z/{a_{\rm s}})$, and the remaining primitive
is $-(1+\log(z/{a_{\rm s}}))/z$.  Thus there is no permutation factor and
no extra factor $2$.  The fixed rational walls in
\eqref{eq:current-Z} remain within the same admissible finite
positive family.  Restriction uses the positive parent majorant and its assigned cell decomposition.

The directed interval estimate of Appendix~\ref{app:refinement} gives
\begin{equation}
 \sum_{\mathscr B\in \mathscr E_0\cup \mathscr E_1}G_{\mathscr B}>
 \frac{558583}{10^9}=0.000558583=:g.
 \label{eq:current-extra-gross}
\end{equation}
Only this rational lower endpoint is used below.

\paragraph{The discarded positive lower estimator.}
The lower estimator assigned to a rectangle $\mathscr B$ is the maximum of its
constant and chord estimates for the ordinary term, plus its nonnegative
double-sieve defect. The complete sum on $\mathscr E$, including the Chen factor
$1/4$, is denoted by $L_{\mathscr E}$. Its rational fixed-point evaluation gives
\begin{equation}
 L_{\mathscr E}=\frac{5645403422587826010135}{2417851639229258349412352}
 <\frac{466977}{200000000}=0.002334885=:\lambda_{\rm loss}.
 \label{eq:current-literal-loss}
\end{equation}
We next bound this discarded lower functional directly by rational intervals.

To specify the enclosure, on a positive cell use the residual
function $R(t,u)=57/64-(5/2)t-u$ used in the preliminary functional,
and put
\begin{align*}
 r_{\mathscr B}&=R(t_r,u_r),\quad R_{\mathscr B}=R(t_l,u_l),\quad
 s_l=r_{\mathscr B}/{a_{\rm s}},\quad s_h=R_{\mathscr B}/{a_{\rm s}},\\
 f_l&=\frac{\log(s_l-1)}{s_l},\qquad
 f_h=\frac{\log(s_h-1)}{s_h}.
\end{align*}
Here $2\le s_l\le s_h\le9/4$.  Define
\begin{align*}
 M_{\mathscr B}&=\log(t_r/t_l)\log(u_r/u_l),\\
 M_{R,\mathscr B}
 &=\int_{\mathscr B} R(t,u)\frac{dt\,du}{tu}\\
 &=\frac{57}{64}M_{\mathscr B}
   -\frac52(t_r-t_l)\log(u_r/u_l)
   -(u_r-u_l)\log(t_r/t_l).
\end{align*}
The two ordinary alternatives, after division by the Chen
factor $4$, are bounded above by
\begin{equation}
 \frac{f_lM_{\mathscr B}}{{a_{\rm s}}},
 \qquad
 \frac{f_l(R_{\mathscr B}M_{\mathscr B}-M_{R,\mathscr B})+
       f_h(M_{R,\mathscr B}-r_{\mathscr B}M_{\mathscr B})}{a_{\rm s}(R_{\mathscr B}-r_{\mathscr B})}.
 \label{eq:current-floor-proxy}
\end{equation}
Indeed the exact endpoint values $f_l,f_h$ dominate the
lower endpoints used in the functional.  Both coefficient integrals are
nonnegative because $r_{\mathscr B}\le R(t,u)\le R_{\mathscr B}$, and the fixed-point lower roundings can only decrease the estimator value.
The cellwise estimator takes a maximum of its constant and
chord alternatives; taking the maximum of the two upper bounds
in \eqref{eq:current-floor-proxy} preserves the direction.

For the separate defect term, take $h_{\mathscr B}^{\min}$ to be the minimum of
the $h_{248}$ closed-cell values meeting the exact interval
$[s_l,s_h]$, including both adjoining cells at a knot.  This
minimum is at least the one over the outward-rounded lookup
interval.  Therefore $h_{\mathscr B}^{\min}M_{\mathscr B}/R_{\mathscr B}$ is an upper bound for the
cellwise defect estimator, and can be added to the preceding maximum.
All relevant $h_{248}$ entries are nonnegative.  If the outward lookup crosses below $2$ and returns zero, the same
upper bound remains valid.  On the $227$ crossing cells the
constant, eligible chord alternatives and lookup all give zero.
Every $\mathscr E_0$ cell has $s_h\le2$ and likewise has zero cellwise lower estimator.

The proxy sum has a rational upper enclosure strictly below
$\lambda_{\rm loss}$. Appendix~\ref{app:refinement-review} obtains it by range
reduction and alternating logarithm series with explicit remainders.
Thus $\lambda_{\rm loss}$ can replace the payment
in~\eqref{eq:current-crop-loss} without weakening the lower inequality.

The recovered payment is exactly
\begin{equation}
 0.002372094-\lambda_{\rm loss}=\frac{37209}{10^9}=0.000037209.
 \label{eq:current-recovered-loss}
\end{equation}
The retained $L_6$ lower endpoint is the whole $L_6$ lower
endpoint minus $4\lambda_{\rm loss}$, namely
\[
 \frac{1421673188323272382945058469141}
      {472236648286964521369600000000}.
\]
The total removed negative gross in Chen units is now the safe
value $0.002095161+0.003799514+g=0.006453258$.
Combining~\eqref{eq:current-extra-gross} and
\eqref{eq:current-recovered-loss} gives the additional lower coefficient
\begin{equation}
 \Delta=g+(0.002372094-\lambda_{\rm loss})
 =\frac{37237}{62500000}=0.000595792.
 \label{eq:current-refinement-delta}
\end{equation}

We add $\Delta/8=37237/500000000=0.000074474$
to each of the two error allowances.  The two enlarged allowances cover finite density approximation and
coefficient realization, respectively.
The net increment after both extra reserves is
\begin{equation}
 \Delta-2(\Delta/8)=\frac{3\Delta}{4}
                  =\frac{111711}{250000000}=0.000446844.
 \label{eq:current-refinement-net}
\end{equation}
The walls and queries are finite, so their local accuracy
can be chosen inside these positive additional allowances by
the same common approximation and coefficient realization argument used
for the preliminary lower functional.

\subsection{The lower endpoint and the order of limits}
\label{sec:lower-limits}

With the paired-domain improvement, the raw endpoint and each
realization allowance are
\begin{align}
 R_{\rm low}&=R_{\rm base}+\Delta,\label{eq:current-R}\\
 B_{\rm low}&=B_{\rm base}+\Delta/8.\label{eq:current-B}
\end{align}
Their exact values are specified in Appendix~\ref{app:numeric}.
From \eqref{eq:current-weight-to-count}, the actual lower coefficient realization,
finite induction, and evaluation of the paired lower functional,
\begin{equation}
 \liminf_{\substack{N\to\infty\\2\mid N}}
         \frac{D_{1,2}(N)}{H_N}\ge R_{\rm low}-2B_{\rm low}.
 \label{eq:current-lower-derived}
\end{equation}
The two allowances leave
\begin{equation}
                  R_{\rm low}-2B_{\rm low}>c+4.9061\cdot10^{-9},
                  \label{eq:current-after-budgets}
\end{equation}
so in particular
\begin{equation}
                  R_{\rm low}-2B_{\rm low}-c>4.9061\cdot10^{-9}>0.
                  \label{eq:current-lower-slack}
\end{equation}
In particular $D_{1,2}(N)\ge cH_N$ eventually, with $c$ as in~\eqref{eq:current-safe-constants}.

Fix the finite set of estimates and their dependencies, the tail
and compact approximation tolerances, the cutoffs defining fixed compact
parameter domains, the prefix caps, and the support reserves. Next choose the common beta and rough
parameters, fine labels, fixed divisor and moment orders, admissible
$\kappa$, and local logarithmic accuracy. Only then let $N$ tend to
infinity through the even integers. Taking the maximum of the finitely
many thresholds makes all estimates simultaneous. Since $C_N\ge C_2>0$,
fixed power savings and $O_A(N/\log^AN)$ with $A>2$ are uniformly
$o(H_N)$. The finite induction therefore gives the conclusion for every
sufficiently large even target.

\section{Finite numerical estimates}
\label{sec:current-numeric}

The preceding upper and lower constructions reduce the numerical
part of the proof to the following finite inequalities.
Table~\ref{num:certificate-table} gives the correspondence between
these quantities and the certificate fields used to verify them.
Appendix~\ref{app:numeric} defines their integrals, partitions,
directed enclosures, and error allowances.

\begin{lemma}
\label{lem:numerical-endpoints}
The upper density functional has an upper endpoint $E_{\rm up}$ as
in~\eqref{eq:current-E}. It admits finite density and coefficient
realizations with errors at most $u$ each, where $u$ is
as in~\eqref{eq:current-u}, and $E_{\rm up}+2u<K$.
The eleven-term lower functional with paired deletion has a lower
endpoint $R_{\rm low}=R_{\rm base}+\Delta$ as in~\eqref{eq:current-R}.
The corresponding two errors are at most $B_{\rm low}$ each, where $B_{\rm low}$ is
as in~\eqref{eq:current-B}, and $R_{\rm low}-2B_{\rm low}>c$.
The generation-$430$ envelopes used in $U_3,U_4$ are generated by
$430$ simultaneous applications of the positive finite action map
to the stated seed. The generation-$248$ envelopes used in the other
specified terms have the separate enclosures described in
Appendix~\ref{app:numeric}.
\end{lemma}

\begin{proof}
For the upper estimate, evaluate the four terms
\[
 E_{\rm up}=\frac4\Theta+I_1+I_{3,\mathrm{compact}}+I_{3,\mathrm{tail}}
       -\sum_{j=1}^{22}G_j,
\]
using upper endpoints for the four costs and lower endpoints for the
nonnegative savings $G_j$. Each successive change is measured from its
parent; overlapping alternatives are maximized before integration.
The noncompact tail is bounded over its entire range.

For the lower estimate, apply lower endpoints to $L_1,L_2,L_5,L_6$
and upper endpoints to $U_3,U_4,U_7,\ldots,U_{11}$
in~\eqref{eq:current-root-functional}. The changes of lower starts,
restored exit regions, and finite-table corrections are inserted on
their prescribed disjoint regions or as complete parent replacements.
The paired-domain terms are then evaluated by
\eqref{eq:current-crop-integral}, \eqref{eq:current-crop-loss},
\eqref{eq:current-extra-gross}, and~\eqref{eq:current-refinement-delta}.
This gives the rational endpoint $R_{\rm low}$.

All scalar operations are rational. Logarithms are enclosed by series
with explicit remainders; the sieve functions are enclosed by
oriented interpolation estimates. Affine cells are triangulated using
exact vertices and exact polynomial moments. Positive savings may be
bounded on a subset, with the remaining saving discarded; an upper
bound for a negative cost covers the full retained region. The finite
recurrence reads only the preceding generation and rounds every
nonnegative contribution downwards. The detailed enclosure and coverage
statements are given in Appendix~\ref{app:numeric}. Rational subtraction
gives~\eqref{eq:current-upper-slack} and~\eqref{eq:current-lower-slack}.
\end{proof}

\paragraph{Numerical data.}
The finite partitions, recurrence, and directed integral bounds of
Appendix~\ref{app:numeric} are supplied in
\emph{Theorem (1+1.75) on the Goldbach Conjecture --- Numerical
Supplement}, version 20261106~\cite{ChanSupplement}. The appendix describes the data and their certificate
correspondence with Lemma~\ref{lem:numerical-endpoints}.

\section{Ordered bad semiprimes and the final subtraction}
\label{sec:current-transfer}

We now estimate $B_b(N)$. Its dependence on $b$ is contained in
the unsifted mass of the ordered semiprimes; the positive product
estimate then imposes the prime complement.

Fix $1/2<b<1$ and set $\beta=b/(1+b)$.  Products below
$2N^{1-\kappa_0}$ have a negligible count.  For a remaining bad
product $m=rq$, the inequality $r>q^b$ implies $r>m^\beta$;
both primes therefore have a fixed positive exponent.  Use the
coordinates
\[
                       t=\frac{\log r}{\log N},\qquad y=\frac{rq}{N}.
\]
On a fixed radial strip $\eta\le y\le C$, PNT in $q=Ny/r$
followed by reciprocal-prime partial summation in $r$ gives the
unsifted mass element
\begin{equation}
                         \frac{N}{\log N}\,
                              \frac{dt\,dy}{t(1-t)}.
 \label{eq:current-pair-measure}
\end{equation}
The PNT errors are uniform on positive exponent compacta.  A prime
upper bound controls $y<\eta$ by $O_b(\eta N/\log N)$; $N$ tends
to infinity first, then $\eta$ tends to zero.

The exact ordered bad inequalities are
\[
 0<y<1,\qquad
 t\le\frac12+\frac{\log y}{2\log N},\qquad
 t>\beta+\frac{b\log y}{(1+b)\log N}.
\]
Their limiting unsifted mass is
\begin{equation}
 J_b=\int_0^1dy\int_\beta^{1/2}\frac{dt}{t(1-t)}
     =\left[\log\frac{t}{1-t}\right]_\beta^{1/2}
     =\log(1/b).
 \label{eq:current-bad-mass}
\end{equation}
The ordering $r\le q$ counts each composite once.  There is no
factor $2$ and no further singular factor in $J_b$.  The
equal-prime diagonal has size $O(\sqrt N)$ and can be retained or
restored after a squarefree argument.

Smooth the complete-product wall in $y$ and the remaining walls
in $t$.  For any fixed $\epsilon>0$ a positive outer approximation
has limiting mass at most $J_b+\epsilon$ and stays in the admissible
finite family.  A partition of unity repeats no mass.
Proposition~\ref{prop:current-positive-product}, followed by
$N\to\infty$ and then $\epsilon\downarrow0$, gives
\begin{equation}
 \limsup_{\substack{N\to\infty\\2\mid N}}
             \frac{B_b(N)}{H_N}\le K\log(1/b).
 \label{eq:current-bad-bound}
\end{equation}

Unique factorization gives the exact identity
\begin{equation}
                         D_b(N)=D_{1,2}(N)-B_b(N).
 \label{eq:current-exact-subtraction}
\end{equation}
The unit was excluded in \eqref{eq:current-d12}; a convention
including it would change the count by at most one.  Zero is
excluded by $p<N$.  A square is a bad remainder when $b<1$, and
a prime remainder is a good one.  Combining
\eqref{eq:current-lower-derived}, \eqref{eq:current-lower-slack}
and \eqref{eq:current-bad-bound} proves
Theorem~\ref{thm:current-main}.

It remains to verify the strict scalar inequalities in
Corollaries~\ref{cor:three-quarters} and~\ref{cor:current-explicit}. For rational $b$ put $z=(1-b)/(1+b)$.
The positive atanh series and its geometric tail give
\begin{equation}
 2\sum_{j=0}^{95}\frac{z^{2j+1}}{2j+1}
 <\log(1/b)
 <L_{96}(b):=
 2\sum_{j=0}^{95}\frac{z^{2j+1}}{2j+1}
       +\frac{2z^{193}}{193(1-z^2)}.
 \label{eq:current-log-certificate}
\end{equation}
Every endpoint is rational.  Substituting
\eqref{eq:current-safe-constants} gives
\begin{align}
 c-KL_{96}(3/4)&>\frac{16532064}{10^9}>0.0165,\notag\\
 c-KL_{96}(74816964/10^8)
     &>\frac{1001149712517311}{10^{20}}>10^{-5},\notag\\
 c-KL_{96}(74816865/10^8)
     &>\frac{106416285358807}{10^{20}}>10^{-6},
                 \label{eq:current-rational-margins}\\
 c-KL_{96}(9/10)
     &>\frac{124934244400949969098}{10^{20}}>1.24889.\notag
\end{align}
The respective surpluses exceed $3.2064\cdot10^{-5}$,
$1.1497\cdot10^{-8}$, $6.4162\cdot10^{-8}$ and
$4.5244\cdot10^{-4}$.  First choose the fixed realization and
positive-smoothing errors inside these strict margins; then let
the remaining uniform asymptotic errors tend to zero.  Taking the
maximum of the four resulting thresholds proves both corollaries.

\paragraph{Range supplied by the fixed constants.}
For $1/2<b<1$, the right-hand side of \eqref{eq:current-general-theorem}
is positive precisely when $b>\exp(-c/K)$. This value describes the range
supplied by the present pair of coefficient bounds. Extending that range
requires stronger applicable lower or upper estimates, with the
moving-residue and coefficient-class conditions preserved.


\appendix
\section{The restricted medium convolution estimate}
\label{app:medium}

We prove Proposition~\ref{prop:current-medium}. The convolution lengths
are $U,V$, with $x=UV$, and $\mathfrak{a}$ denotes the nonzero residue.
The even target is denoted by $N$ in the main text.
Put $\vartheta=7/32$ and
\[
 z=x^{1/(\log\log x)^3},\qquad y=x^{1/\log\log x}.
\]
The coefficients $\alpha_n$ are fixed throughout the modulus sums and
satisfy the auxiliary-coprimality Siegel--Walfisz condition
\eqref{eq:current-aux-sw}.
All constants in this condition are uniform over the coefficient family.
We use $e(t)=\exp(2\pi it)$; an inverse in an exponential is always taken
modulo its displayed denominator.

\subsection{A smooth majorant and divisor estimates}

\begin{lemma}\label{ana:majorant}
There is a fixed function $\psi\in C_c^\infty((0,\infty))$ such that
$0\le\psi\le1$, $\psi=1$ on $[1,2]$, $\supp\psi\subset[1/2,5/2]$, and
\[
 \|\psi^{(j)}\|_\infty\le C^j(j!)^2,
 \qquad |\widehat\psi(t)|\ll \exp(-c\sqrt{|t|}).
\]
Consequently, for polynomially bounded positive $\ell,V$ and
$H\ge C_1\ell(\log x)^5/V$,
\begin{equation}
 \frac V\ell\sum_{|h|>H}
       \left|\widehat\psi\left(\frac{hV}{\ell}\right)\right|
       \ll_D x^{-D}
 \label{ana:poisson-tail}
\end{equation}
for every fixed $D>0$.
\end{lemma}

\begin{proof}
Choose $\eta_j=c_0/j^2$ with $\sum \eta_j=1/2$, let $u_j$ be the uniform
probability measure on $[-\eta_j/2,\eta_j/2]$, and put
$\phi=u_1*u_2*\cdots$. Its support is $[-1/4,1/4]$. Differentiating
the first $j$ factors as signed measures gives total variation at most
$\prod_{i\le j}(2/\eta_i)=C^j(j!)^2$. The remaining convolution is a
probability measure; retaining further undifferentiated factors also
gives a continuous density at each fixed derivative order. The function
$\psi=\one_{[3/4,9/4]}*\phi$ has the asserted support and derivative
bounds. Integration by parts $j$ times, with $j$ a sufficiently small
multiple of $\sqrt{|t|}$, gives the Fourier estimate. A decreasing-integral
comparison bounds the tail in \eqref{ana:poisson-tail} by a polynomial
factor times $\exp(-c_1(\log x)^{5/2})$. If $H<1$, every omitted integer
already has $|h|V/\ell\gg(\log x)^5$, and the same argument applies.
\end{proof}

Write $n=\log_xU$, $q=\log_xQ_1$, $u=\log_xQ_2$ and $v=\log_xQ_3$.
The hypotheses of Proposition~\ref{prop:current-medium} imply
\begin{equation}
 n>\sigma,\qquad n\le\frac12-4\sigma,\qquad
 q+u+v<\frac7{10}-\frac{13\sigma}2.
 \label{ana:lengths}
\end{equation}
In particular $Q_1Q_2Q_3\le x^{4/5}$ for large $x$.
Split $q_2=u_0s$ and $q_3=v_0t$ into their $z$-rough and $z$-smooth
parts. The terms $s>y$ or $t>y$ have arbitrarily large logarithmic
saving. To see this, first collapse the modulus coefficient with its
fixed divisor bound. The primitive progression bound for divisor-bounded
functions~\cite{Shiu}, followed by Rankin's method with exponent $1/\log z$, gives
\begin{equation}
 \sum_{\substack{\ell\le x^{4/5}\\ \operatorname{sm}(\ell;z)\ge y}}
       \frac{\tau(\ell)^C}{\ell}
 \ll \exp\bigl(- (\log\log x)^2+O_C(\log\log x)\bigr).
 \label{ana:smooth-tail}
\end{equation}
Here $\operatorname{sm}(\ell;z)$ is the largest divisor of $\ell$ all of
whose prime factors are at most $z$. The bound is used with divisor
means, rather than with a pointwise factor $x^\epsilon$.

We also use the following translated progression estimate. For fixed
$C,K,\delta>0$, an interval $I\subset[-x^K,x^K]$ of length at most $L$
and $f\le Lx^{-\delta}$ satisfy
\begin{equation}
 \sum_{\substack{v\in I\\v\equiv r\pmod f\\v\ne0}}
        \tau(|v|)^C
       \ll \frac Lf\tau(f)^C(\log x)^{O(1)}.
 \label{ana:shifted-divisor}
\end{equation}
Indeed, put $g=(r,f)$ and $v=gw$. The resulting progression modulo
$f/g$ is primitive and its length-to-modulus ratio is $L/f\ge x^\delta$.
The range $|w|\le(f/g)x^{\delta/2}$ is bounded by a fixed divisor
estimate. Split the rest by sign and into dyadic intervals, subdividing
them into lengths comparable to the smaller of their scale and $L/g$.
The primitive divisor bound applies on each interval with fixed positive
length and modulus margins; summation proves \eqref{ana:shifted-divisor}.
Cauchy's inequality gives its mixed version for
$\tau(v)^B\tau(|v-\mathfrak{a}|)^C$, with $v=\mathfrak{a}$ excluded.

On a dyadic block put $d=st\asymp D$, $w=u_0v_0\asymp\mathcal R$ and
define
\begin{equation}
 \Lambda_d(w)=\sum_{\substack{st=d\\u_0v_0=w}}
       \gamma_2(u_0s)\gamma_3(v_0t)
       \one_{\text{fixed individual intervals and supports}}.
 \label{ana:Lambda}
\end{equation}
Thus $\Lambda_d$ is independent of $q_1$,
$|\Lambda_d(w)|\le\tau(dw)$, $D\le y^2=x^{o(1)}$, and
$D\mathcal R\asymp Q_2Q_3$. With $Q=Q_1$, $Y=QD$ and
$\delta=\sigma/3$, \eqref{ana:lengths} implies
\begin{equation}
 U,V\ge x^\delta,\qquad Y\le Ux^{-\delta},\qquad
 Y\mathcal R\le x^{1-\delta}.
 \label{ana:dispersion-lengths}
\end{equation}
We write $\mathcal S$ for the original factorized discrepancy sum
restricted to this block. All values of $d\asymp D$ remain in this
single sum before Cauchy's inequality is applied.

We reduce to one-bounded $\beta$ before applying dispersion. For
divisor-bounded $\beta$, retain $\tau(m)\le(\log x)^C$ and normalize
by $(\log x)^{BC}$. The discarded coefficient at $v=nm$ is at most
$(\log x)^{-C}\tau(v)^{B+2}$. For $v\ne\mathfrak a$, summing the
moduli costs at most $\tau(|v-\mathfrak a|)^3$; fixed shifted divisor
moments therefore give $O(x(\log x)^{E-C})$, with $E$ independent
of $C$. The exact product $v=\mathfrak a$ contributes separately at
most $x^{4/5+o(1)}$. Ordinary divisor means and the reciprocal
totient sum control the mean term. Choose $C$ first, then request
sufficient logarithmic accuracy for the normalized estimate.

\subsection{Dispersion and the principal variance}

The dispersion step separates a common principal term from the
nonzero frequencies. A remaining nonnegative principal variance is
estimated by the auxiliary Siegel--Walfisz hypothesis before the
exponential sums are treated.

Set $\ell=qd$, $\eta_r=\Lambda_d(r)$ and $h_0=\int\psi$. For fixed
$(q,d)$ define
\begin{align*}
 B_r&=\sum_{(n,\ell r)=1}\alpha_n,
 &A_r(b)&=\sum_{\substack{n\equiv b\pmod\ell\\(n,r)=1}}\alpha_n,\\
 C_r(m)&=\frac{\one_{(m,\ell r)=1}}{\varphi(\ell r)}B_r,
 &X_r(m)&=\sum_{(n,\ell r)=1}\alpha_n
       \bigl(\one_{mn\equiv \mathfrak{a}\pmod{\ell r}}-\one_{mn=\mathfrak{a}}\bigr).
\end{align*}
All $n$ sums retain the original interval. Remove $mn=\mathfrak{a}$ before applying
Cauchy's inequality. The first-moment cost is
$Y\mathcal R x^{o(1)}$, which is power-small by
\eqref{ana:dispersion-lengths}. Nonnegative majorants for $q$ and $m$
then give a Cauchy multiplier $VY(\log x)^C$ and the exact variance
\begin{align}
 \mathcal V'&=S_1-2\Re S_2'+S_3',\notag\\
 S_1&=\sum_{q,d}\psi(q/Q)\sum_m\psi(m/V)
                 \left|\sum_r\eta_rC_r(m)\right|^2,\notag\\
 S_2'&=\sum_{q,d}\psi(q/Q)\sum_m\psi(m/V)
       \left(\sum_r\eta_rX_r(m)\right)
       \overline{\left(\sum_s\eta_sC_s(m)\right)},\notag\\
 S_3'&=\sum_{q,d}\psi(q/Q)\sum_m\psi(m/V)
                 \left|\sum_r\eta_rX_r(m)\right|^2.
 \label{ana:variance}
\end{align}

The common principal term is governed by the kernel
\[
 c_\ell(r,s)=\frac{\varphi(\ell rs)}{
       \ell rs\varphi(\ell r)\varphi(\ell s)}.
\]
If $L$ is any common multiple of $\ell$ and the finitely many supported
$r$, then
\begin{equation}
 c_\ell(r,s)=\frac1L\sum_{v\bmod L}
 \frac{\one_{(v,\ell r)=1}}{\varphi(\ell r)}
 \frac{\one_{(v,\ell s)=1}}{\varphi(\ell s)}.
 \label{ana:gram}
\end{equation}
This is a positive Gram kernel. The size of $L$ is irrelevant because it
is used only in the finite identity. Smooth coprime counting yields
\begin{equation}
 S_1=\mathcal M+O\left(\frac{U^2(\log x)^C}{Y}\right),\qquad
 \mathcal M=Vh_0\sum_{q,d}\psi(q/Q)
        \sum_{r,s}\eta_r\overline{\eta_s}
               c_\ell(r,s)B_r\overline{B_s}.
 \label{ana:mean}
\end{equation}

For $S_2'$, remove $(m,s)>1$ while $mn\ne \mathfrak{a}$. A rough common factor
$f>z$ gives factors $\mathcal R/f$ and $x/f$; applying
\eqref{ana:shifted-divisor} and summing
$\sum_{f>z}\tau(f)^C/f^2$ bounds this error by
$VU^2(\log x)^C/(Yz)$. Only then restore $mn=\mathfrak{a}$ for Poisson summation;
the cost is $\mathcal RUx^{o(1)}$. The zero frequency has coefficient
$1/(\ell r\varphi(\ell s))$, whose ratio to $c_\ell(r,s)$ is
$1+O((\log x)/z)$. Consequently
\begin{equation}
 S_2'=\mathcal M+VE_1+
 O_A\left(\frac{VU^2}{Y(\log x)^A}+\mathcal RUx^{o(1)}\right),
 \label{ana:cross}
\end{equation}
where, with the original coprimalities retained,
\begin{align}
 E_1={}&\sum_{q,d}\psi(q/Q)\sum_{r,s}
     \frac{\eta_r\overline{\eta_s}}{\ell r\varphi(\ell s)}
 \sum_{\substack{n_1,n_2\\(n_1,\ell r)=(n_2,\ell s)=1}}
       \alpha_{n_1}\overline{\alpha_{n_2}}\notag\\
 &\hspace{12mm}\times
 \sum_{0<|h|\le H_1}\widehat\psi\left(\frac{hV}{\ell r}\right)
               e\left(\frac{\mathfrak{a}h\overline{n_1}}{\ell r}\right),
 \qquad H_1\asymp\frac{Y\mathcal R}{V}(\log x)^5.
 \label{ana:E1-definition}
\end{align}
Lemma~\ref{ana:majorant} permits this Fourier truncation.

For $S_3'$, retain both exact-product exclusions during all divisor
cutoffs and common-factor removals. The clean set has
\begin{equation}
 \begin{gathered}
 (r,s)=(s,\ell)=1,\quad (n_1,\ell r)=(n_2,\ell s)=1,
 \quad(n_1,n_2)=1,\\
 n_1\equiv n_2\pmod\ell,\qquad
 |n_1-n_2|\ge U(\log x)^{-C_{\rm sep}}.
 \end{gathered}
 \label{ana:clean-set}
\end{equation}
The removed rough common factors cost
$VU^2(\log x)^C/(Yz)+x(\log x)^C$; the close pairs cost
$VU^2(\log x)^{C-C_{\rm sep}}/Y+x(\log x)^C$.
These estimates follow respectively by counting two multiples or a
multiple and a primitive progression; their additional unit terms are
absorbed using $Y/U\le x^{-\delta}$. Choose the divisor truncation and
then $C_{\rm sep}$ sufficiently large.

Now $n_1\ne n_2$. Restoring one equality $mn_i=\mathfrak{a}$ costs only
$\mathcal R(U+Y)x^{o(1)}=\mathcal RUx^{o(1)}$: fixing its
factorization leaves $O(U/Y+1)$ choices for the other $n$, and the other
rough modulus divides a nonzero shifted product. Both equalities cannot
hold. Apply CRT and Poisson only after this restoration. On
\eqref{ana:clean-set},
\[
 \varphi(\ell)c_\ell(r,s)=\frac1{\ell rs}.
\]
Put $d_r(b)=A_r(b)-B_r/\varphi(\ell)$. Since $\sum_b^*A_r(b)=B_r$,
the zero frequency is $\mathcal M+\mathcal P$, where
\begin{equation}
 \mathcal P=Vh_0\sum_{q,d}\psi(q/Q)\varphi(\ell)
   \sum_{r,s}\eta_r\overline{\eta_s}c_\ell(r,s)
                     \sum_b^*d_r(b)\overline{d_s(b)}\ge0.
 \label{ana:principal-variance}
\end{equation}
Positivity follows from \eqref{ana:gram}. In particular this principal
variance must be bounded; it does not cancel identically.

The auxiliary Siegel--Walfisz hypothesis and the primitive character
large sieve give, for every fixed $A,w>0$,
\begin{equation}
 \sum_{\ell\le8Y}\tau(\ell)^w\sum_b^*|d_r(b)|^2
       \ll_A\frac{\tau(r)^{2C_0}U^2}{(\log x)^A}.
 \label{ana:filtered-BDH}
\end{equation}
For clarity, normalize $\alpha_n\one_{(n,r)=1}$ by $\tau(r)^{C_0}$.
Its auxiliary Siegel--Walfisz constants remain uniform because
$\tau(kr)\le\tau(k)\tau(r)$. Character orthogonality expresses the
unweighted variance through nonprincipal characters. Write each modulus
as $fk$, where $f>1$ is the primitive conductor, and keep $(n,k)=1$
in the coefficient before applying the large sieve. On a dyadic
conductor interval $f\asymp R_0$, its bound at fixed $k$ is
$O(U(U/R_0+R_0)/\varphi(k))$. Summing conductors
$f>(\log x)^b$ and inducing factors gives
$O(U^2(\log x)^{2-b}+UY)$. Small conductors are bounded by the
auxiliary Siegel--Walfisz condition with its $\tau(k)^{C_0}$ cost;
their total is $O(U^2(\log x)^{3b+C_1-2B})$.
Choose $b$, and then $B$, to obtain any prescribed logarithmic saving.
The weight $\tau(\ell)^w$ follows by Cauchy's inequality, ordinary
divisor means and the elementary bound for the variance
$O(U^2/\ell)$.

Using \eqref{ana:gram}, Cauchy over the $r$ variables and
$\varphi(\ell r)^{-1}\ll(\log x)/(Y\mathcal R)$ gives, for a fixed
divisor order $B_0$,
\begin{align*}
 \mathcal P\ll{}&\frac{V(\log x)^2}{Y\mathcal R}
  \sum_{r\asymp\mathcal R}\tau(r)^{2B_0}
  \sum_{\ell\le8Y}\tau(\ell)^{2B_0+1}\sum_b^*|d_r(b)|^2\\
 \ll_A{}&\frac{VU^2}{Y(\log x)^A}.
\end{align*}
The factor $\tau(r)^{2C_0}$ in \eqref{ana:filtered-BDH} is summed by a
divisor mean; it is not included in an unspecified uniform constant.
Combining these estimates with \eqref{ana:variance} gives
\begin{equation}
 |\mathcal S|\ll_A\frac{x}{(\log x)^A}
  +V\sqrt Y(\log x)^C\bigl(\sqrt{|E_1|}+\sqrt{|E_2|}\bigr),
 \label{ana:dispersion-bound}
\end{equation}
where $E_2$ is the nonzero CRT frequency sum on
\eqref{ana:clean-set}. The restored single-zero term contributes at
most $\sqrt{xY\mathcal R}\,x^{o(1)}$, which is power-small.

\subsection{The first exponential sum}

For a fixed smooth $w$, Poisson summation and the complete Kloosterman
Weil bound give, for every integer $A$,
\begin{equation}
 \left|\sum_{(n,\ell)=1}w(n/U)e(A\overline n/\ell)\right|
 \ll_w \tau(\ell)^2\sqrt{\ell}+\frac U{\ell}(A,\ell).
 \label{ana:inverse-completion}
\end{equation}
The zero Fourier frequency is a Ramanujan sum and gives the second
term. For the other frequencies use
\[
 \sqrt{(k,A,\ell)}\le\sum_{\substack{d\mid(A,\ell)\\d\mid k}}\sqrt d,
 \qquad
 \sum_{\substack{k\ne0\\d\mid k}}(1+|k|/V)^{-2}\le\frac{2V}{d}.
\]
This proves \eqref{ana:inverse-completion} also when the frequency
length is below one.

Fix $d$ in \eqref{ana:E1-definition}, put $Q'=dQ$ and restrict $q'$
to multiples of $d$. This restriction is permitted inside the bounded
$q'$-dependent coefficient. The following calculation applies with
$Q'$ denoted temporarily by $Q$ and $\mathcal R$ by $R$.
Absorb the bounded smooth factor into the $r,h$ coefficient and apply
Cauchy while retaining the original unrestricted $r,s$ support. Expansion
of the square gives
\begin{equation}
 |E_1(d)|^2\ll x^\epsilon\frac{U^3}{Q^3R^2}\,T,
 \label{ana:E1-cauchy}
\end{equation}
where
\[
 T=\sum_{\substack{q\asymp Q, r,r'\asymp R\\(\mathfrak{a},qrr')=1}}
       \sum_{0<|h|,|h'|\le H_1}
 \left|\sum_{(n,qrr')=1}w(n/U)
 e\left(\frac{\mathfrak{a}(hr'-h'r)\overline n}{qrr'}\right)\right|.
\]
Put $j=hr'-h'r$. The $j=0$ tuples number $RH_1x^\epsilon$, since
$hr'\ne0$ and $h'r=hr'$ is a divisor equation. Their contribution is
$\ll x^\epsilon QRH_1U$. For $j\ne0$, the square-root term in
\eqref{ana:inverse-completion} contributes
$\ll x^\epsilon Q^{3/2}R^3H_1^2$.
For its Ramanujan term write $r=gu$, $r'=gv$, $(u,v)=1$. Then
\[
 (j,rr')\le g^2(h,r)(h',r'),\qquad
 \sum_{q\asymp Q}(j,q)\le3Q\tau(|j|),
\]
and
\[
 \sum_{0<|h|\le H_1}(h,r)\le2H_1\tau(r),\qquad
 \sum_{r,r'\asymp R}(r,r')^2\ll R^3.
\]
Because $(\mathfrak{a},qrr')=1$, $\mathfrak{a}$ disappears from these gcds. The total
Ramanujan contribution to $T$ is $\ll x^\epsilon UH_1^2R$.
As $H_1\ll QRU(\log x)^5/x$, \eqref{ana:E1-cauchy} becomes
\begin{equation}
 |E_1(d)|^2\ll x^\epsilon\frac{U^4}{Q^2}
  \left(\frac{UQ}{x}+\frac{UQ^{5/2}R^3}{x^2}
                         +\frac{U^2QR}{x^2}\right).
 \label{ana:E1-polynomial}
\end{equation}

Return to $Q=Q_1$ and the logarithmic variables in
\eqref{ana:lengths}. The required gaps follow from
\begin{align*}
 n+q&<1-9\sigma,\\
 2n+q+u+v&<\frac{17}{10}-\frac{29\sigma}2,\\
 n+\frac52q+3u+3v
 &=\frac{35}{16}(q-n)+\frac{23}{16}(2n+u+2v)
   +\frac5{16}(n+q+5u+2v)-\frac v2\\
 &<\frac{1021}{512}-\frac{147\sigma}{8}+O(1/\log x).
\end{align*}
The last inequality uses $|\mathfrak{a}|\ge c_0x$ and $\vartheta=7/32$.
The divisor bound for $\Lambda_d$ and $D=x^{o(1)}$ can now be paid
inside a power saving. Sum $U^2/(dQ)$ over $d\asymp D$ using
$\sum_{d\asymp D}1/d\ll1$. With room to absorb one further $D$,
\begin{equation}
 |E_1|\ll\frac{U^2}{DQ}x^{-\sigma/2}.
 \label{ana:E1-final}
\end{equation}

\subsection{The second exponential sum and the reciprocity phase}

Write $E_2=\sum_{d\asymp D}E_3(d)/d$. For fixed $d$, $E_3(d)$ is the
sum over $q$, $w_1,w_2\asymp\mathcal R$, $n_1,n_2\asymp U$ and
$0<|h|\le H$, where $H\asymp QD\mathcal R^2(\log x)^5/V$, of
\begin{align}
 &\frac{\Lambda_d(w_1)\overline{\Lambda_d(w_2)}
               \alpha_{n_1}\overline{\alpha_{n_2}}}{qw_1w_2}
 \psi(q/Q)\widehat\psi\left(\frac{hV}{qdw_1w_2}\right)\notag\\
 &\hspace{8mm}\times
 e\left(\frac{\mathfrak{a}h\overline{n_1w_2}}{qdw_1}
          +\frac{\mathfrak{a}h\overline{n_2qdw_1}}{w_2}\right).
 \label{ana:E3-definition}
\end{align}
The restrictions are
\[
 \begin{gathered}
 (q,\mathfrak{a})=1,\quad(w_1,\mathfrak{a}w_2)=1,\quad(w_2,\mathfrak{a}qdw_1)=1,\\
 n_1\equiv n_2\pmod{qd},\quad
 (n_1,n_2qdw_1)=(n_2,n_1qdw_2)=1,\quad
 |n_1-n_2|\ge U(\log x)^{-C_{\rm sep}}.
 \end{gathered}
\]
There is no independent bounded $q$ coefficient after the outer Cauchy
step. The denominator in \eqref{ana:E3-definition} is $qw_1w_2$;
the factor $1/d$ occurs outside $E_3(d)$.

Exact reciprocity gives
\begin{align}
 &e\left(\frac{\mathfrak{a}h\overline{n_1w_2}}{qdw_1}
          +\frac{\mathfrak{a}h\overline{n_2qdw_1}}{w_2}\right)\notag\\
 &\qquad=
 e\left(\frac{\mathfrak{a}h(n_1-n_2)\overline{n_2qdw_1}}{n_1w_2}\right)
 e\left(\frac{\mathfrak{a}h}{qdw_1w_2n_1}\right).
 \label{eq:current-retained-phase}
\end{align}
All inverses are justified by the displayed coprimalities. Split $h$
by sign and the $q$ support into a bounded number of intervals on scales
$q_0\asymp Q$. Normalize the five variables
$(q,w_1,w_2,n_1,h)$ by $(q_0,\mathcal R,\mathcal R,U,H)$.
The last phase has amplitude
\[
 \frac{\mathfrak{a}}{x}\frac Q{q_0}\frac Dd(\log x)^5.
\]
Mixed differentiation at most once in each coordinate therefore bounds
the total five-dimensional variation by
\begin{equation}
 \bigl(1+(\log x)^5+(|\mathfrak{a}|/x)(\log x)^5\bigr)^5
       =x^{5\kappa+o(1)}.
 \label{ana:phase-variation}
\end{equation}
Repeated partial summation leaves anchored rectangular sums. Their
endpoints are fixed before estimation, and their $h$ coefficient is a
single fixed signed-interval indicator $\omega_h$, $|\omega_h|\le1$.
It has no dependence on any remaining arithmetic variable.

Partition the rough coprime pairs $(n_1w_2,n_2)$ by the cross-coprime
partition of \cite[Lemma 13.2]{MaynardI}; its number of parts is
$\exp(O((\log\log x)^4))$. Expand the two copies of
\eqref{ana:Lambda} independently, including both choices of the
factorization $st=d$. Put $R\asymp v_0$, $S\asymp u_0$.
A total-product interval is retained as
\[
 \one_{A_i<r_is_i\le B_i}
       =\one_{r_is_i\le B_i}-\one_{r_is_i\le A_i}.
\]
Thus it becomes a difference of the permitted product upper cutoffs,
without inserting a coupled condition into a one-variable coefficient.
For a bilinear form, complex polarization replaces two different
coefficient sequences by diagonal sequences with at most four times
the diagonal bound. Apply this successively to the two $R$, two $S$
and two $n$ sequences, at total constant cost at most $4^3$.

We give the resulting unweighted sum explicitly. Fix $d$, a
cross-coprime class $\mathcal C$, endpoints $q_0\le q_1\le2q_0$,
two product upper cutoffs $B_1,B_2$, and one-bounded sequences
$\gamma_r,\lambda_s,\alpha_n,\omega_h$. All these data are fixed
before the remaining summations. Let $\mathscr C$ be the set of
tuples $(q,r_1,r_2,s_1,s_2,n_1,n_2)$ satisfying
\begin{equation}\label{ana:bare-domain}
\begin{gathered}
 q_0\le q\le q_1,\quad (q,\mathfrak{a})=1,\quad
 r_i\asymp R,\quad s_i\asymp S,\quad n_i\asymp U,\\
 (r_1s_1,\mathfrak{a}r_2s_2)=(r_2s_2,\mathfrak{a}qdr_1s_1)=1,\quad
 r_1s_1\le B_1,\quad r_2s_2\le B_2,\\
 n_1\equiv n_2\pmod{qd},\quad
 (n_1,n_2qdr_1s_1)=(n_2,n_1qdr_2s_2)=1,\\
 (n_1r_2s_2,n_2)\in\mathcal C,\qquad
 |n_1-n_2|\ge U(\log x)^{-C_{\rm sep}}.
\end{gathered}
\end{equation}
Define
\begin{equation}\label{ana:bare-definition}
\mathcal W_{\rm bare}=
 \sum_{\mathscr C}
 \frac{\gamma_{r_1}\lambda_{s_1}
           \overline{\gamma_{r_2}\lambda_{s_2}}
           \alpha_{n_1}\overline{\alpha_{n_2}}}
      {q r_1r_2s_1s_2}
 \sum_{0<|h|\le H}\omega_h
 e\left(\frac{\mathfrak{a}h(n_1-n_2)\overline{n_2r_1s_1dq}}
               {n_1r_2s_2}\right).
\end{equation}
There is no smooth reciprocity phase in this definition.
The preceding endpoint separation and polarization give
\begin{equation}\label{ana:E3-bare-reduction}
 |E_3(d)|\ll x^{5\kappa+o(1)}
       \sup|\mathcal W_{\rm bare}|,
\end{equation}
where the supremum is over the choices specified above. The
$x^{o(1)}$ includes the cross-coprime classes and the two independent
factorizations of $d$, but the five-variable phase loss occurs only
once. It remains to bound \eqref{ana:bare-definition}.

\paragraph{The first Cauchy inequality.}
Put $f=(n_1-n_2)/(dq)$, which is a signed nonzero integer. The new
variable satisfies $0<|f|\le2U/(dq_0)$, and
\begin{equation}\label{ana:f-phase}
 e\left(\frac{\mathfrak{a}h(n_1-n_2)\overline{n_2r_1s_1dq}}
                    {n_1r_2s_2}\right)
 =e\left(\frac{\mathfrak{a}hf\overline{n_2r_1s_1}}
                    {n_1r_2s_2}\right).
\end{equation}
Keep $q_f=(n_1-n_2)/(df)$ integral and in $[q_0,q_1]$ at this stage.
The product cutoff in $s_1$ has the exact finite Fourier expansion
\[
 \one_{r_1s_1\le B_1}=
 \int_0^1D_{r_1}(\xi)e(s_1\xi)\,d\xi,
 \qquad
 D_{r_1}(\xi)=\sum_{1\le j\le B_1/r_1}e(-j\xi).
\]
There is a common nonnegative majorant
$L(\xi)\ll\min(1+B_1/R,1+\|\xi\|^{-1})$ with
$\int_0^1L(\xi)\,d\xi\ll\log x$ and
$|D_{r_1}(\xi)|\le L(\xi)$; set
$c_{r_1}(\xi)=D_{r_1}(\xi)/L(\xi)$.
For fixed $\xi$ write
$\lambda_s'=S\lambda_s e(s\xi)/s$.
These coefficients are one-bounded on $s\asymp S$.
If $\mathscr D$ denotes the conditions in \eqref{ana:bare-domain}
which are independent of $s_1$, after putting $q=q_f$, the exact
Fourier integrand is
\begin{align}\label{ana:W2-definition}
 \mathcal W_2(\xi)={}&\frac1S\sum_{\mathscr D}
 \frac{\gamma_{r_1}c_{r_1}(\xi)
              \overline{\gamma_{r_2}\lambda_{s_2}}
              \alpha_{n_1}\overline{\alpha_{n_2}}}
      {q_f r_1r_2s_2}\notag\\
 &\quad\times
 \sum_{\substack{s_1\asymp S\\(s_1,\mathfrak{a}n_1r_2s_2)=1}}
 \lambda_{s_1}'\sum_{0<|h|\le H}\omega_h
 e\left(\frac{\mathfrak{a}hf\overline{n_2r_1s_1}}{n_1r_2s_2}\right).
\end{align}
In particular
$|\mathcal W_{\rm bare}|\ll(\log x)\sup_\xi|\mathcal W_2(\xi)|$.
For each $f$ the number of pairs $n_1,n_2\asymp U$ with
$n_1\equiv n_2\pmod{df}$ is
$O(U^2/(d|f|)+U)$. Summing over $f$ and then $r_1,r_2,s_2$
gives $O(U^2R^2S\log x)$ outer tuples. Since the outside
coefficient in \eqref{ana:W2-definition} has absolute value at most
$O(1/(q_0R^2S^2))$, Cauchy's inequality gives
\begin{equation}\label{ana:first-Cauchy}
 |\mathcal W_2(\xi)|^2
 \ll\frac{U^2\log x}{q_0^2R^2S^3}\,\mathcal W_3,
\end{equation}
where the following nonnegative sum retains the congruence but
discards other outer restrictions:
\begin{align}\label{ana:W3-definition}
 \mathcal W_3={}&
 \sum_{0<|f|\le2U/(dq_0)}
 \sum_{\substack{n_1,n_2\asymp U\\n_1\equiv n_2\pmod{df}}}
 \sum_{r_1,r_2\asymp R}
 \sum_{\substack{s_2\asymp S\\(n_2r_1,n_1r_2s_2)=1}}\notag\\
 &\quad\times
 \left|\sum_{\substack{s_1\asymp S\\(s_1,\mathfrak{a}n_1r_2s_2)=1}}
 \lambda_{s_1}'\sum_{0<|h|\le H}\omega_h
 e\left(\frac{\mathfrak{a}hf\overline{n_2r_1s_1}}{n_1r_2s_2}\right)
 \right|^2.
\end{align}
Every discarded restriction is outside the squared inner sum. The
cross-coprime property of $\mathcal C$ justifies the displayed
coprimality before the original class restriction is removed.

\paragraph{Completion of the positive square.}
We may now remove the congruence in \eqref{ana:W3-definition}
and enlarge $|f|\le2U/q_0$, because all its summands are nonnegative.
Put $b=n_2r_1$ and $c=n_1r_2s_2$. Their multiplicities are at most
$\tau(b)\tau_3(c)$ and can be bounded before summation.
Split $b,c$ and $|f|$ into dyadic intervals. On each such block
$B\ll UR$, $C\ll URS$, $F\ll U/q_0$, and a fixed smooth
nonnegative majorant $g(b,c)=\psi(b/B)\psi(c/C)$ gives
\begin{equation}\label{ana:second-Cauchy}
 \mathcal W_3\ll x^\epsilon\sum_{B,C,F,\pm}\mathcal W_4,
\end{equation}
where the sign of $f$ is absorbed into $\mathfrak{a}$ and
\begin{equation}\label{ana:W4-definition}
 \mathcal W_4=
 \sum_{\substack{b,c\\(b,c)=1}}g(b,c)
 \sum_{f\asymp F}
 \left|\sum_{\substack{s\asymp S\\(s,\mathfrak{a}c)=1}}
 \lambda_s'\sum_{0<|h|\le H}\omega_h
 e\left(\frac{\mathfrak{a}hf\overline{bs}}c\right)\right|^2.
\end{equation}
There are only logarithmically many dyadic blocks. Combining
\eqref{ana:first-Cauchy} and \eqref{ana:second-Cauchy}, and using
$q_0\asymp Q$, gives the outside factor in the form
\begin{equation}\label{ana:bare-Cauchy-factor}
 |\mathcal W_{\rm bare}|
 \ll x^\epsilon\frac{U}{QRS^{3/2}}
       \sup_{B,C,F}\mathcal W_4^{1/2}.
\end{equation}

Expanding \eqref{ana:W4-definition} and interchanging the two
$s$ labels when necessary yields
\begin{equation}\label{ana:W4-expansion}
 \mathcal W_4=
 \sum_{\substack{s_1,s_2\asymp S\\(s_1s_2,\mathfrak{a})=1}}
 \sum_{0<|h_1|,|h_2|\le H}\lambda_{s_2}'\overline{\lambda_{s_1}'}
       \omega_{h_1}\overline{\omega_{h_2}}
 \sum_{f\asymp F}\sum_{\substack{b,c\\(c,bs_1s_2)=1}}
 g(b,c)e\left(\frac{\mathfrak{a}f(h_1s_1-h_2s_2)\overline{bs_1s_2}}c\right).
\end{equation}
The terms with $h_1s_1=h_2s_2$ have total absolute value at most
$x^\epsilon HSBFC$. For the others set
$z=s_1s_2$, $y=f(h_1s_1-h_2s_2)\ne0$, and define the coefficient
\begin{equation}\label{ana:beta-definition}
 \beta(y,z)=
 \sum_{\substack{s_1s_2=z,\ s_i\asymp S\\
       f\asymp F,\ 0<|h_i|\le H\\
       f(h_1s_1-h_2s_2)=y}}
 \lambda_{s_2}'\overline{\lambda_{s_1}'}
       \omega_{h_1}\overline{\omega_{h_2}}.
\end{equation}
We retain $(z,\mathfrak{a})=1$. The variables $b,c$ do not occur in
\eqref{ana:beta-definition}; the restrictions involving them remain
in the kernel. Consequently the nonzero contribution is exactly
\begin{equation}\label{ana:off-diagonal-kernel}
 \sum_{\substack{y\ne0,z\\(\mathfrak{a},z)=1}}\beta(y,z)
 \sum_{\substack{b,c\\(c,bz)=1}}g(b,c)
                         e(\mathfrak{a}y\overline{zb}/c).
\end{equation}
In particular $|\beta(y,z)|\le b(z,y)$ for the positive tuple count
in the following lemma. This domination also holds after any dyadic
restriction on $y,z$. It is independent of the future Poisson
frequency, which has not yet been introduced.

\begin{lemma}\label{ana:multiplicity}
Let $b(z,y)$ count tuples
$s_i\asymp S$, $f\asymp F$, $0<|h_i|\le H$ with
$z=s_1s_2$ and $y=f(h_1s_1-h_2s_2)\ne0$. Then
\begin{equation}
 \sum_z\sum_{u\ne0}b(z,\nu u)^2
 \ll x^\epsilon(1+H/S)\frac{(HS)^2F}{\nu}
 \label{ana:multiplicity-bound}
\end{equation}
for every $\nu\ge1$ on polynomially bounded scales. The left side
vanishes when $\nu\ge8FHS$.
\end{lemma}

\begin{proof}
At fixed $(z,y)$ there are at most $\tau(z)\tau(|y|)$ decompositions
$(s_1,s_2,f,k)$, with $k=h_1s_1-h_2s_2\ne0$.
Cauchy therefore reduces the square to collisions of $h$ pairs with
the same $(s_1,s_2,k)$, followed by a single $f$ count.
Group according to $c\mid\nu$, $\nu/c\mid k$, with the count of $f$
bounded by $2F/c$. Equal $h$ pairs contribute
$\ll x^\epsilon(HS)^2F/\nu$.
For unequal pairs put $g=(s_1,s_2)$, $s_1=gu$, $s_2=gv$, $(u,v)=1$.
Their differences have the form $(h_1-h_1',h_2-h_2')=(jv,ju)$,
with $0<|j|\le2Hg/S$. If $b=(g,\nu/c)$, the condition on
$k/g$ is divisibility by $\nu/(cb)$. Counting nonzero multiples from
zero, without an additional unit term, gives at most
$O(HScb/(g\nu))$ such values. Factoring the remaining nonzero
integer $h_2v$ costs only a divisor factor. The unequal contribution
is consequently
\[
 \ll x^\epsilon\frac{H^3SF}{\nu}
             \sum_{g\le2S}\frac{(g,\nu)}g
 \ll x^{2\epsilon}\frac{H^3SF}{\nu}.
\]
The last estimate follows from
$\sum_{g\le2S}(g,\nu)/g\le\tau(\nu)(1+\log(2S))$.
Renaming $2\epsilon$ proves the result. The support bound follows from
$0<|y|<8FHS$.
\end{proof}

In the present application
$H/S\le x^{o(1)}(Q/U)D_0<1$, where
$D_0=U^2R^2S/x$. Thus the factor $1+H/S$ is bounded.

\subsection{The reciprocal kernel and the spectral estimate}

We estimate the reciprocal kernel with the coefficient independent
of the future Poisson frequency. This separation makes the same
frequency interval available at every spectral level, as required
by the level-average exceptional estimate.

We use $Y$ for the coefficient length, $Z$ for the level length,
and $P$ for the length of a nonzero Poisson-frequency block in this
subsection. The original convolution lengths remain $U,V$.
For $\beta(u,z)$ supported on $u\asymp Y$, $z\asymp Z$ and
$(\mathfrak{a},z)=1$, put
\begin{equation}
 \mathcal K=\sum_{u,z}\beta(u,z)
   \sum_{\substack{c,b\\(c,bz)=1}}
   \psi(c/C)\psi(b/D)e(\mathfrak{a}u\overline{zb}/c).
 \label{ana:kernel}
\end{equation}
Either sign of $u$ is treated separately. Define
$B^2=\sum_{u,z}|\beta(u,z)|^2$ and
$B_\nu^2=\sum_{u,z}|\beta(\nu u,z)|^2$; further coprimality
indicators or unit phases do not increase these norms.

Poisson summation in $b$ is the exact identity
\[
 \sum_{(b,c)=1}\psi(b/D)e(\mathfrak{a}u\overline z\,\overline b/c)
 =\frac Dc\sum_{m\in\mathbb Z}\widehat\psi(mD/c)
                   S(m\overline z,\mathfrak{a}u;c).
\]
The zero frequency contributes at most
$x^\epsilon D\sqrt{YZ}\,B$, since
$|S(0,\mathfrak{a}u;c)|\le(\mathfrak{a}u,c)$ and
$\sum_{c\asymp C}(\mathfrak{a}u,c)\ll C\tau(|\mathfrak{a}u|)$.
Its square is $x^\epsilon D^2YZB^2$.
For nonzero $m$ truncate at the common value $(C/D)x^\epsilon$ and
split into common signed dyadic intervals. Ordinary smooth Fourier
decay, with a fixed sufficiently large derivative order, controls
the discarded tail.

On an actual interval $m\asymp P$, put
\[
 Y_0=\frac{4\pi\sqrt{|\mathfrak{a}|PY}}{C\sqrt Z},\quad
 v=m/P,\quad w=u/Y,\quad \zeta=z/Z,\quad
 t=\frac{4\pi\sqrt{|\mathfrak{a}|mu}}{c\sqrt z\,Y_0}.
\]
Then $c/C=\sqrt{vw/\zeta}/t$. This change of variables removes $\mathfrak{a}$
from the normalized amplitude. Fourier expansion on a fixed compact
box separates the amplitude into a Bessel test function, a common
factor $\one_I(m)e(m\omega)$ and an arbitrary complex $u,z$
coefficient. Choosing the Fourier decay order above the fixed Bessel
derivative order makes the coefficient norm summable. Thus the
$m$ interval $I$ and phase $\omega$ are common at every level $z$.

The cusp conversion and packet operation are given precisely in
Appendix~\ref{app:classical}. They transform the $u$ coefficient to
\begin{equation}
 \beta_\nu(n,z)=\one_{(n,\mathfrak{a})=1}\beta(\nu n,z)e(-n/z),
 \qquad \nu\mid \mathfrak{a}^\infty,
 \label{ana:packet-coefficient}
\end{equation}
at packet cost $(|\mathfrak{a}|\nu)^{\vartheta/2}$, while leaving the common
$m$ functional unchanged. Let
\[
 P_0=\frac{PC^4Z}{|\mathfrak{a}|PY+C^2Z},\qquad T_0=\frac{|\mathfrak{a}|PY}{C^2}.
\]
The regular spectral contribution and the exceptional estimates yield
\begin{align}
 |\mathcal L_P|^2\ll x^\epsilon P_0\bigg\{
 &\sum_{\nu\mid \mathfrak{a}^\infty}(|\mathfrak{a}|\nu)^\vartheta
       (Y/\nu+Z+T_0)(P+Z+T_0)B_\nu^2\notag\\
 &+(C^2Z)^\vartheta(Y+Z)^{1-\vartheta}
                    (P+Z)^{1-2\vartheta}B^2\bigg\}.
 \label{ana:spectral-polynomial}
\end{align}
Here $\mathcal L_P$ is the separated Kuznetsov sum, before restoring
the Poisson factor $D/C$.

We give the spectral summation leading to this polynomial. Throughout
this step, a regular spectral norm is the positive direct-sum norm
\eqref{ana:regular-norm}: its Maass atoms have weight
$1/\cosh(\pi t_j)$, its holomorphic atoms weight $1/(2\pi)$ after
the stated Fourier normalization, and its continuous part measure
$dt/\pi$. Cauchy's inequality is applied within each component of this
same measure, before the three components are added.
At fixed $z,\nu$, write $A(t),B(t)$ for the $n$- and $m$-side
coefficient arrays on $d\mu_z$, with entries
$\sum_n\beta_\nu(n,z)R_{t,1}(n)$ and
$\sum_{m\in I}e(m\omega)R_{t,\infty}(m)$, respectively.
Here $t$ labels an entry of the regular direct sum, including its
component and basis labels, and $R_{t,\mathfrak c}$ denotes the
corresponding Fourier coefficient in~\eqref{ana:regular-norm}.
Let $A_0,B_0\ge0$ denote the norm-bound amplitudes, and let
$u_0,v_0\ge0$. If the
two cumulative regular squared norms are bounded by
$(T^2+u_0)A_0^2$ and $(T^2+v_0)B_0^2$, the Bessel bounds
\[
 |\mathcal H(t)|\ll\frac{1+|\log Y_0|}{1+Y_0},\qquad
 |\mathcal H(t)|\ll |t|^{-5/2}+Y_0|t|^{-3}
 \quad(|t|\ge\max(1,2Y_0))
\]
give
\[
 \int|\mathcal H(t)A(t)B(t)|\,d\mu_z(t)
 \ll \frac{1+|\log Y_0|}{1+Y_0}
 \sqrt{(1+Y_0^2+u_0)(1+Y_0^2+v_0)}A_0B_0.
\]
Indeed, above $\max(1,2Y_0)$ the cumulative norm product grows on
the $j$th shell by at most $2^{2j}$, leaving the summable factors
$2^{-j/2}+2^{-j}$. The level Cauchy factor is $\sqrt Z$.
With $L=Y/\nu$, this gives the first term of
\eqref{ana:spectral-polynomial}, including the packet cost.

For exceptional forms put $\tau_j=2\sqrt{1/4-\lambda_j}\le\vartheta$,
$X_0=1+Z/L$, and $X=1/Y_0$. Use
\[
 1+X^{\tau_j}\le
 2(1+X/\sqrt{X_0})^{\tau_j}(1+X_0)^{\tau_j/2}.
\]
The fixed-level $n$ norm is
$\ll x^\epsilon(1+L/Z)B_\nu^2$. The common-interval $m$ norm is
\[
 \ll x^\epsilon P(P+Z)
 \left[1+\left\{
  \frac{P(1+X/\sqrt{X_0})^2}{(P+Z)^2}\right\}^{\vartheta}\right].
\]
Multiplying these bounds and the squared packet cost gives a term
absorbed in the regular expression and the second term of
\eqref{ana:spectral-polynomial}. The factor $|\mathfrak{a}|\nu$ cancels only
in this latter term after the two norms are multiplied.
The same argument applies to the holomorphic component with height
$k-1$ and to the continuous component with real height $t$ in
\eqref{ana:regular-norm}. Their measures remain unchanged throughout
the shell summation. On the exceptional spectrum,
$1/\cosh(\pi t_j)=1/\cos(\pi|\operatorname{Im}t_j|)\asymp1$,
since $|\operatorname{Im}t_j|\le7/64$; hence the unweighted
exceptional estimates used above differ only by an absolute constant.

Restoring $D/C$, the elementary inequality
\[
 \frac{Z(L+Z+T_0)(P+Z+T_0)}{Z+T_0}
       \le(L+Z)(P+Z)+ZT_0
\]
and $P\le C/D$ give
\[
 |\mathcal K|^2\ll x^\epsilon\mathcal J^2,
\]
where the complete nonnegative polynomial is
\begin{align}
 \mathcal J^2={}&\sum_{\nu\mid \mathfrak{a}^\infty}(|\mathfrak{a}|\nu)^\vartheta
 \bigl\{C(Y/\nu+Z)(C+DZ)+|\mathfrak{a}|YZ\bigr\}B_\nu^2
 \notag\\
 &+\bigl\{
 C(CD\sqrt Z)^{2\vartheta}(Y+Z)^{1-\vartheta}
                         (C+DZ)^{1-2\vartheta}+D^2YZ\bigr\}B^2.
 \label{ana:returned-polynomial}
\end{align}
The actual cutoff $P\le x^\epsilon C/D$ is accommodated by reducing
the initially chosen $\epsilon$. If $C/D<1/2$, estimates are applied
only on nonempty actual intervals; the real comparison value $C/D$
is not used as the length of a nonexistent integer interval.

\paragraph{Return to the convolution scales.}
Apply \eqref{ana:returned-polynomial} first at the actual dyadic
lengths in \eqref{ana:off-diagonal-kernel}. All its coefficients are
nonnegative and increasing in the four lengths, since
$0\le\vartheta\le1/2$. We may then replace these lengths by
\begin{equation}\label{ana:scale-substitution}
 D=B_*=UR,\quad C=C_*=URS,\quad Z=S^2,\quad
 Y=\frac{UR^2S^3}{V}.
\end{equation}
The accompanying counting parameters satisfy
$F\le U/Q$ and $H\le x^\epsilon H_0$, where
$H_0=QR^2S^2/V$. These selectable powers are included in the final
$x^\epsilon$ after choosing the input exponents smaller. Put
\[
 D_0=\frac{U^2R^2S}{x}=\frac YZ<1.
\]
Then $B_*Z\ge C_*$, $B_*^2<x$, and $C_*B_*Z=xY$.
Lemma~\ref{ana:multiplicity} and
$H/S\le x^\epsilon(Q/U)D_0<1$ give
\begin{equation}\label{ana:norm-insertion}
 B_\nu^2\le x^\epsilon\frac{(HS)^2F}{\nu},\qquad
 B^2\le x^\epsilon(HS)^2F.
\end{equation}
The auxiliary divisor sum satisfies
\[
 \sum_{\nu\mid\mathfrak a^\infty}\nu^{\vartheta-1}
 =\prod_{p\mid\mathfrak a}(1-p^{-(1-\vartheta)})^{-1}
 =x^{o(1)}.
\]
Its logarithm is $O((\log x)^\vartheta)$, by splitting the prime
divisors at $\log x$; this is uniform for polynomially bounded
$\mathfrak a$. Consequently \eqref{ana:returned-polynomial} gives
\begin{equation}\label{ana:J-reduced}
 \mathcal J^2\ll x^\epsilon(HS)^2F
 \left\{
 |\mathfrak a|^\vartheta(C_*B_*Z+|\mathfrak a|Y)Z
 +C_*^{1+2\vartheta}B_*Z^{2-2\vartheta}
 +B_*^2YZ\right\}.
\end{equation}
For clarity, the four monomials on the right yield the following
ratios after multiplication by
\[
 \frac{(H_0S)^2(U/Q)}{U^4R^4S^6}
                  =\frac{Q}{U^3V^2}:
\]
\begin{equation}\label{ana:ratio-table}
\begin{array}{c|c}
 \text{monomial}&\text{normalized contribution}\\ \hline
 |\mathfrak a|^\vartheta C_*B_*Z^2
     &D_1=|\mathfrak a|^\vartheta UQR^2S^5/x^2\\
 |\mathfrak a|^{1+\vartheta}YZ
     &(|\mathfrak a|/x)D_1\\
 C_*^{1+2\vartheta}B_*Z^{2-2\vartheta}
     &D_2=U^{1+2\vartheta}QR^{2+2\vartheta}S^{5-2\vartheta}/x^2\\
 B_*^2YZ&QR^4S^5/V^3\le D_1.
\end{array}
\end{equation}
The last comparison follows from $B_*^2<x$ and $|\mathfrak a|\ge1$.
The zero-difference contribution in \eqref{ana:W4-expansion}
satisfies
\[
 \left(\frac{U}{QRS^{3/2}}\right)^2H_0SB_*\frac UQ C_*
       =\left(\frac{U^2}{Q}\right)^2D_0.
\]
For the nonzero part, \eqref{ana:off-diagonal-kernel} gives
$|\mathcal W_{4,\ne0}|\ll x^\epsilon\mathcal J$.
Combining this square root with that in
\eqref{ana:bare-Cauchy-factor} gives explicitly
\[
 \left(\frac{U}{QRS^{3/2}}\right)^4
 \frac{\mathcal J^2}{(U^2/Q)^4}
 =\frac{\mathcal J^2}{U^4R^4S^6}
 \ll x^\epsilon\bigl\{(1+|\mathfrak a|/x)D_1+D_2\bigr\}.
\]
The two contributions prove
\begin{align}
 |\mathcal W_{\rm bare}|&\ll\frac{U^2}{Q}x^\epsilon
   \left\{D_0^{1/2}
    +\bigl[(1+|\mathfrak{a}|/x)D_1+D_2\bigr]^{1/4}\right\},
 \label{eq:current-bare-bound}\\
 D_0&=\frac{U^2R^2S}{x},\qquad
 D_1=\frac{|\mathfrak{a}|^\vartheta UQR^2S^5}{x^2},\qquad
 D_2=\frac{U^{1+2\vartheta}QR^{2+2\vartheta}S^{5-2\vartheta}}{x^2}.
 \notag
\end{align}

\subsection{Completion of the convolution estimate}

The scale hypotheses give $D_0\ll x^{-8\sigma}$ and
$D_1,D_2\ll x^{-15\sigma}$ on every small-factor block.
Together with $|\mathfrak{a}|/x\le x^\kappa$ and the single variation cost
\eqref{ana:phase-variation}, \eqref{eq:current-bare-bound} gives
\[
 |E_3(d)|\ll\frac{U^2}{Q}
 \left(x^{-4\sigma+5\kappa+\epsilon}
       +x^{-15\sigma/4+21\kappa/4+\epsilon}\right).
\]
All selectable epsilon losses, the vanishing subpowers from the
partitions and one additional normalizing factor $D$ fit inside
$x^{\sigma/10}$. Since $\kappa\le\sigma/100$, both remaining
savings exceed $\sigma/2$. Harmonic summation in $d$ therefore gives
\begin{equation}
 |E_2|\ll\frac{U^2}{DQ}x^{-\sigma/2}.
 \label{ana:E2-final}
\end{equation}
Insert \eqref{ana:E1-final} and \eqref{ana:E2-final} into
\eqref{ana:dispersion-bound}. The exponential-sum contribution is
$O(x^{1-\sigma/4}(\log x)^C)$. The earlier single-zero restoration
is $O(x^{1-\sigma/6+o(1)})$, still power-small.
Choose the auxiliary Siegel--Walfisz and logarithmic precisions to
pay for the finitely many dyadic choices, the fixed logarithmic number
of blocks and the normalization of $\beta$. Equation
\eqref{ana:smooth-tail} deals with the omitted smooth parts.
This proves \eqref{eq:current-medium-conclusion}.

\section{Normalization of the spectral inputs}
\label{app:classical}

We record the normalizations to identify the positive measures used
in the spectral Cauchy inequalities. The regular spectrum and the
exceptional spectrum are treated in their respective ranges.

These inputs are used in \eqref{ana:spectral-polynomial}; the packet coefficient and the common frequency interval are specified in \eqref{ana:packet-coefficient} and the discussion preceding that estimate.

We specify the classical statements used in the preceding proof.
They are applied at trivial character, level $z$, cusps $\infty$ and
$1$, with $(\mathfrak{a},z)=1$. Write
$R(n)=\sqrt{|n|}\rho(n)$ for the Deshouillers--Iwaniec normalization
of the Fourier coefficients. The cusp scaling matrices are
\[
 \sigma_\infty=I,\qquad
 \sigma_1=\begin{pmatrix}\sqrt z&0\\\sqrt z&1/\sqrt z\end{pmatrix}.
\]
The cusp conversion in \cite[Lemma 4.3]{Drappeau} introduces
$e(\mathfrak{a}y/z)$; hence the ordinary Kloosterman coefficient carries
$e(-\mathfrak{a}y/z)$.

For $y=\nu n$, $\nu\mid \mathfrak{a}^\infty$, $(n,\mathfrak{a})=1$, the coprime packet
identities in \cite[Lemmas 3.1 and 3.2]{ABL} give, on every normalized
oldform dilation,
\[
 R_v(\mathfrak{a}\nu n)=\lambda_v(\mathfrak{a}\nu)
                   e\bigl(n(\mathfrak{a}\nu-1)/z\bigr)R_v(n).
\]
Multiplication by $e(-\mathfrak{a}\nu n/z)$ leaves exactly $e(-n/z)$.
The packet coefficients in \cite[(3.3)]{ABL} have Frobenius norm
$z^\epsilon$ after divisor losses. Cauchy's inequality in that packet
and the Kim--Sarnak exponent bound~\cite{KimSarnak} therefore cost
$(|\mathfrak{a}|\nu)^{7/64+\epsilon}$ and give the common coefficient
\eqref{ana:packet-coefficient}. No eigenvalue is assumed common to
all dilations. The holomorphic and continuous quadratic forms in
\cite[Lemma 3.2]{ABL} use the corresponding identities with
$\vartheta=0$. The continuous term is integrated with the formal
Plancherel measure, rather than treated as the norm of an individual
square-integrable Eisenstein series; see also
\cite[Proposition 8.2]{Young}.

For completeness we specify the positive norm used for every regular
spectral Cauchy inequality. Let $\mathfrak c$ be either cusp. In the
notation of \cite[Proposition E]{Pascadi}, write
\begin{align*}
 R^M_{j,\mathfrak c}(n)&=
       \rho^{\rm DI}_{j,\mathfrak c}(n)
       =\sqrt{|n|}\rho^{\rm ABL}_{j,\mathfrak c}(n),\\
 R^H_{jk,\mathfrak c}(n)&=
       \left(\frac{(k-1)!}{(4\pi)^{k-1}}\right)^{1/2}
       n^{-(k-1)/2}\psi_{jk,\mathfrak c}(n),\\
 R^E_{\mathfrak d,\mathfrak c}(n,t)&=
       n^{it}\varphi_{\mathfrak d\mathfrak c}
                       (n;\tfrac12+it).
\end{align*}
The symbol $R_j$ in the packet formulas denotes $R_j^M$.
For $T\ge1$ and $b$ supported on $n\asymp L$, $L\ge1/2$, define
\begin{align}\label{ana:regular-norm}
 \mathfrak Q_{z,\mathfrak c}(T;b)
 ={}&\sum_{\substack{t_j\in\mathbb R\\|t_j|\le T}}
 \frac1{\cosh(\pi t_j)}
 \left|\sum_{n\asymp L}b_nR^M_{j,\mathfrak c}(n)\right|^2\notag\\
 &+\frac1{2\pi}\sum_{\substack{k\in2\mathbb Z_{>0}\\k-1\le T}}
      \sum_j\left|\sum_{n\asymp L}b_nR^H_{jk,\mathfrak c}(n)\right|^2
       \notag\\
 &+\frac1\pi\sum_{\mathfrak d}\int_{-T}^T
      \left|\sum_{n\asymp L}b_n
                 R^E_{\mathfrak d,\mathfrak c}(n,t)\right|^2dt.
\end{align}
This is the direct-sum measure $d\mu_z$ used above. The constants
$1/(2\pi)$ and $1/\pi$ are the constants in the Kuznetsov formula;
the factor $1/\cosh(\pi t_j)$ is part of the Maass measure and is
not absorbed by the change from $\rho^{\rm ABL}$ to $R^M$.
By \cite[Proposition 4.7]{Drappeau} and
\cite[Lemma H]{Pascadi},
\begin{equation}\label{ana:regular-large-sieve}
 \mathfrak Q_{z,\mathfrak c}(T;b)
 \ll (zLT)^\epsilon\left(T^2+\frac Lz\right)\sum_n|b_n|^2.
\end{equation}
The harmless change from $k\le T$ to $k-1\le T$ costs only a fixed
factor. The continuous term in this convention has measure $dt/\pi$,
as in \cite[(3.27)]{Pascadi}; it has no additional hyperbolic factor.
The test-function bounds used with this norm are
\cite[Lemma D]{Pascadi}. For opposite-sign Kuznetsov sums the
holomorphic term is absent, and the same positive norm is still an
upper bound.

Put $\tau_j=2\sqrt{1/4-\lambda_j}\le7/32$ on the exceptional
spectrum. The fixed-level estimate \cite[Theorem A]{Pascadi}, at
cusp $1$, supplies
\begin{equation}
 \sum_{\rm exc}X^{\tau_j}
       \left|\sum_{n\asymp L}b_nR_{j,1}(n)\right|^2
 \ll(zL)^\epsilon(1+L/z)\sum_n|b_n|^2,
 \qquad X\ll\max(1,z/L).
 \label{ana:exceptional-fixed}
\end{equation}
The coefficients may depend on $z$ but are independent of spectral
height. The level-average estimate \cite[Theorem J]{Pascadi} gives
\begin{equation}
 \sum_{z\asymp Z}\sum_{\rm exc}X^{\tau_j}
       \left|\sum_{m\in I}e(m\omega)R_{j,\infty}(m)\right|^2
 \ll(ZV)^\epsilon(Z+V)V,
 \quad X\ll\max(V,Z^2/V),
 \label{ana:exceptional-average}
\end{equation}
where $I\subset(V,2V]$ and $\omega$ are common to every level.
For a partial interval $I$, choose an integer $P>4V$ and expand
$\one_I(m)=\sum_{k=0}^{P-1}c_ke(km/P)$ on residues modulo $P$.
The coefficient norm is $\sum|c_k|\ll\log(2+V)$.
Minkowski in the Hilbert norm indexed by $(z,j)$ and
Theorem J with phase $\omega+k/P$ prove
\eqref{ana:exceptional-average}, with only a logarithmic squared
loss absorbed in $(ZV)^\epsilon$. This argument requires the same
interval at every level. For larger $X$, the bound
$\tau_j\le\vartheta$ gives the continuation factor
$1+(VX/(V+Z)^2)^\vartheta$.

The transformed estimates apply on each nonempty interval of length
at least $1/2$, including an endpoint interval containing $n=1$.
The $m$ interval is common to all levels, and the $n$ coefficient
is independent of spectral height. The condition
$(\mathfrak a,z)=1$ is retained in every packet operation.
These are precisely the interval and coefficient restrictions used
in the preceding spectral estimate.
\section{Factorization of the finite sieve coefficients}
\label{app:coefficients}

We prove that the actual upper and lower sieve coefficients have
the factorization required by the distribution estimate. Throughout
this appendix $N$ is the even target. Finite partitions, levels,
and approximation tolerances are chosen before $N$ tends to
infinity; the arithmetic coefficients themselves may depend on $N$.
We write $\omega_{\mathrm{pf}}(n)$ for the number of distinct prime
factors of $n$.

\subsection{Finite policies and residual square conditions}

Write
\[
 \rho_0=\frac{19}{96},\qquad \mu_0=\frac{25}{128},\qquad
 \Lambda_{\rm cap}=\frac{153}{256},\qquad \Theta=\frac{299}{500}.
\]
The number $\Theta$ bounds the complete carrier; $\Lambda_{\rm cap}$ bounds the
high-rank levels.  A finite policy specifies a level $N^\lambda$ at each
of finitely many prime prefixes.  An ancestor decision is fixed before
the children are subdivided.  For decreasing exponents $t_1,\ldots,t_j$,
put $S_j=t_1+\cdots+t_j$.  The full square conditions at an auxiliary
level $N^\ell$ are
\begin{equation}
 S_{j-1}+2t_j\le\ell.
 \label{coef:full-squares}
\end{equation}
If $I$ is a selected subset of the marked slots, put
$b_I=\sum_{i\in I}t_i$.  In addition to \eqref{coef:full-squares}, retain
all the conditions
\begin{equation}
 b_I+\sum_{\substack{i<j\\i\notin I}}t_i+2t_j\le\ell
 \qquad(j\notin I).
 \label{coef:selected-squares}
\end{equation}
For $j$ after the last selected index, the left side of
\eqref{coef:selected-squares} is exactly $S_{j-1}+2t_j$.
Thus adjoining smaller primes preserves the selected conditions whenever
the full conditions are preserved.  An admitted cube also implies the
next square, since
\[
 S_j+2t_{j+1}\le S_{j-1}+3t_j.
\]
This observation applies to both parities of the ordinary sieve tree:
upper first exits have odd rank, and lower first exits have even rank.

The physical level $\lambda$ and the auxiliary covering level $\ell$
have different uses.  The coefficient and its density use $\lambda$;
factorization may use a slightly larger $\ell$.  No change in $\lambda$
is made in obtaining this cover.  Here is the needed choice at the
three-slot alternative.  Put
\[
 A(t)=\frac{1+t}{2},\qquad W(z)=\Lambda_{\rm cap}-\frac38z.
\]
On a coarse box let $\Gamma=\inf\max(A,W)$ and
$\lambda=\Gamma-2e$, where the variation allowance satisfies
$\chi<e/10$.  If $\inf A\ge\lambda+e/2$, choose
$\ell=\lambda+e/4$.  Otherwise $\sup A<\lambda+11e/20$,
so $W\ge\Gamma$ throughout the box; choose $\ell=\lambda+e$.
This leaves a reserve of at least $9e/20$ above the lower bound and $e$
below the upper bound.  The remaining initialization reserves include
\[
 1-\Lambda_{\rm cap}-2\mu_0=\frac3{256},\qquad
 \frac{525}{8704}+\frac{49z}{272}>0.
\]
The higher-rank changes are restricted to compact pieces on which their
support conditions have positive reserve.  Fine endpoint enlargement is then taken
smaller than this reserve.

At an ordinary lower start $t_1\le {h_{\rm s}}=19/125<\mu_0$, the endpoint
admissibility condition satisfies
\[
 \max\left(\frac{1+t_1}{2},\Lambda_{\rm cap}-\frac38t_1\right)
 \ge\frac{249}{448}.
\]
A level with a fixed buffer below $249/448$ therefore has a strict
covering interval.  A restored fourth exit is admitted only after its
own cube and all the full and selected squares have been checked.
For upper prefixes above $\mu_0$, the function
\[
 v(x)=\min\left(\frac{1+x}{2},\frac{57}{64}-\frac32x\right)
\]
satisfies $v(x)>3x$ for $x<\rho_0$.  Selecting $y<x$ gives
$2x+y<3x$; selecting $y+z<x$ gives $2x+y+z<3x$.
These are the initial square bounds for the two additional alternatives.

In a marked state write $S=\log_N B$, $L=\lambda-S$ and
$\alpha=(1-S)/L\ge2$.  A singleton of exponent $Lu$ changes the
residual length to $L(1-u)$ and a pair of exponents $Lu,Lv$ changes it
to $L(1-u-v)$.  In both cases $\alpha$ increases.  The conditions
$2u<1$, or $2v<1$ and $v+2u<1$, are precisely the new full square
conditions.  The selected conditions then follow from
\eqref{coef:selected-squares}; the parent prime ceiling is retained.
The finite construction has at most $440$ elementary transitions.
Its residual exponents are at least $L_{\min}=\ell_0 101^{-440}$
for some $\ell_0>0$, and its appended prime exponents are at least
$L_{\min}/6$.  In particular all switched leaves have a common positive
prime cutoff.

\subsection{Finite approximation and strict margins}

The finite partition is chosen jointly for all the marked queries.
At the third upper exit, the omitted last-prime tail is
\[
 O\!\left(\eta\left(1+\log^2\frac{\rho_0}{\eta}\right)\right),
 \qquad\eta\downarrow0.
\]
At lower exits of rank $r=3,4,5$, the corresponding bound is
$O(\eta(1+\log^{r-1}({h_{\rm s}}/\eta)))$.  The higher changes have positive
last-prime cutoffs; in particular the eighth-rank domains have
$t_8>53/13824$.  On the compact remainder, proper affine walls can be
trimmed with arbitrarily small mass.  A region on which the chosen
label is identically inherited is retained rather than treated as a
wall.  Freezing the coarse decision before refining the children gives
one common finite policy for the count and the density.

For the third envelope, the coarse parameters $e,\chi$ give
$0\le L_{\mathrm{end}}-\lambda\le4e+6\chi$.
Fix the finite policy and its positive reserves. The admissibility
functions are continuous on finitely many compact boxes, so their
strict reserves have a positive minimum. Hence some $\kappa_*>0$
preserves all the inequalities when
$\log_xN\in[1,1+\kappa_*]$ and the endpoint enlargement is
sufficiently small. The margin is obtained on these fixed compact
pieces, after the limiting boundary regions have been removed.

\subsection{Three-factor representations}

We describe the endpoint construction, since the factorization must hold
for the coefficient as a whole.  Let $x$ be the complete local scale,
$q_0=\log_x N$, $\vartheta=7/32$, and $\alpha_0\ge q_0$ the residue
parameter in the local admissibility condition.  For a small fixed $\zeta>0$ put
\[
 \theta_{\rm cov}=q_0\ell+2\zeta,\qquad D=x^{\theta_{\rm cov}-2\zeta},
 \qquad A=Ux^{-\zeta},
\]
where $U$ is the chosen medium length.  Define
\begin{align*}
 v_0&=2\theta_{\rm cov}-1,\\
 u_0&=\min\left\{
 \frac{2-(2+2\vartheta)\theta_{\rm cov}}{3-4\vartheta},
 \frac{2-2\theta_{\rm cov}-\alpha_0\vartheta}{3}\right\}.
\end{align*}
In this subsection the slot exponents and $b_I$ are measured to base
$x$: they are $q_0$ times their preceding base-$N$ values.  The covering
exponent in the full and selected square conditions is therefore
$q_0\ell$.
The selected subset admissibility condition is $v_0\le b_I\le u_0$,
with \eqref{coef:full-squares}--\eqref{coef:selected-squares} imposed
on ordered upper endpoints.  All lines of this admissibility condition are retained
after rescaling, including a line that is dominated when $\alpha_0=1$.

Suppose first that $C=D/(Ax^{b_I})\ge1$.  Put the selected slots in
group two and insert the remaining slots greedily into capacities $A,C$.
At every insertion the selected square bounds the square of the new
slot by the product of the remaining capacities, so at least one accepts
it.  The resulting group endpoints satisfy
\[
 B_1\le A,\qquad B_2=x^{b_I},\qquad B_3\le C
\]
and
\begin{align}
 A^2B_2B_3^2&\le x^{1-4\zeta},\notag\\
 N^\vartheta AB_1B_2^5B_3^2&\le x^{2-4\zeta},
 \label{coef:three-monomials}\\
 A^{1+2\vartheta}B_1B_2^{5-2\vartheta}B_3^{2+2\vartheta}
 &\le x^{2-(4+4\vartheta)\zeta}.\notag
\end{align}
If $C<1$, insert all slots into $A,D/A$ using the full squares.
Here $1<D/A<x^{b_I}$.  Set $B_3=1$ and replace $x^{b_I}$ by $D/A$
in the last two bounds.  The first bound follows from
$A^2B_2\le AD\le x^{23/24-5\zeta/2}$.  Thus no factor with an
endpoint below one is introduced.

For the three-slot $W$ alternative initialize capacities
$A,x^{v_0},D/(Ax^{v_0})$ with the first two slots.  One fits $x^{v_0}$;
the square of the other is at most the product
$x^{1-\theta_{\rm cov}-2\zeta}$ of the two remaining capacities.  Continue
inserting whole slots.  If a first slot fits none of the capacities,
adjoin it to group two.  Its new product $B$ lies in
$(x^{v_0},x^{u_0}]$, since that slot is at most the third slot and
$p_3\le x^{u_0-v_0}$.  Failure of the first capacity together with the
full square gives $ABc<D$, so the two residual capacities are legal.
The full-square greedy argument finishes the partition.  If no failure
occurs, the capacity product gives \eqref{coef:three-monomials} directly.
These algorithms use only endpoint data of a whole atom and its medium
length.  They never use the individual arithmetic tuple or the sign of
a progression discrepancy.

\subsection{The coefficient identity and squarefree restriction}

Refine the interval labels to a common disjoint half-open alphabet.
If $m_I$ slots have label $I$, set $M=\prod_I m_I!$.
The unrestricted prime convolution $N_J$ of the slot list has value
$M/\prod_p\nu_p(n)!$ on its support.  Its sorted squarefree coefficient
is therefore $\one_{\mathrm{sf}}N_J/M$.  For a fixed three-group
partition put $M_j=\prod_I m_{I,j}!$ and let $\gamma_j$ be the group
convolution divided by $M_j$.  Then $0\le\gamma_j\le1$ on every integer
and the exact identity is
\begin{equation}
 \frac{N_J}{M}
 =\frac{M_1M_2M_3}{M}\,
       (\gamma_1*\gamma_2*\gamma_3).
 \label{coef:tensor-identity}
\end{equation}
An empty group is the unit mass at one.  A marked convolution already
contains its marked multiplicities.  Only its ordered residual slots
are divided by their factorials: regrouping a signed atom
$\epsilon\one_{\mathrm{sf}}N_{\mathrm{all}}/M_{\mathrm{res}}$
gives the external scalar $\epsilon M_1M_2M_3/M_{\mathrm{res}}$.
This scalar is fixed, although it need not be bounded by one.

Choose $w=N^{\eta_0}$ below every marked prime exponent.  For $p<w$,
$p\mid BN$ if and only if $p\mid N$. Hence one pair of small-prime
beta weights, with support below the smallest residual reserve, serves
every state in the finite marked family.  The sign of the atom determines
which member is attached.  The small and rough prime universes are
disjoint, so decomposition into these two parts is unique.  Attaching
the small-prime beta weight to one of the groups in
\eqref{coef:tensor-identity} preserves one-boundedness.

Global squarefreeness remains in the physical coefficient and its density.
It is removed only in the analytic tensor expansion.  If the support is
$D_*\le N^{\theta_*}$ with $\theta_*<1$, the change in the discrepancy
is at most a fixed divisor factor times
\begin{align}
 &\sum_{\substack{q\le D_*\\p^2\mid q\text{ for some }p\ge N^{\eta_0}}}
 \left(\frac{N}{q}+1+\frac{N}{\varphi(q)}\right)
 \ll (N+D_*)N^{-\eta_0+o(1)}.
 \label{coef:collision}
\end{align}
The term $1$ includes an exact complete product equal to $N$.
No divisor function is evaluated at zero.

\subsection{Parameter order and the medium estimate}

Before small-prime attachment, endpoint enlargement and local rescaling,
the complete supports at rough rank at most two have exponents
\[
 \frac{19}{48},\qquad \frac{38}{125},\qquad
 \frac{1527}{3584}
\]
for the upper, ordinary lower and marked families, respectively.
Their gaps below $1/2$ are $5/48$, $49/250$ and $265/3584$.
For the last bound use $x<57/224$, $x+2y\le\Lambda_{\rm cap}$, and hence
$x+y\le(\Lambda_{\rm cap}+57/224)/2=1527/3584$.
The fixed coefficient norm permits the below-half-level
Bombieri--Vinogradov estimate for these pieces.  A high-rank coefficient
requires the factorization above.

First fix the finite collection of estimates and their dependencies, coarse policy, factorization conditions and
positive reserves.  Choose $\zeta$ below these reserves and the prime
cutoffs, and set $\sigma=\zeta/2000$.  Next choose a common beta support
small enough for the density and support tolerances.  If the rough-rank
bound is $J$, take the endpoint mesh $\tau$ so that $(J+2)\tau$ is below
every remaining reserve.  Only after the fixed divisor and moment orders
are known choose
\begin{equation}
 0<\kappa\le\min\left\{\kappa_*,\frac{\sigma}{200},
                         \frac1{2K_{\mathrm{mom}}}\right\},
 \qquad
 0<\kappa_0<\frac{\kappa}{4(1+\kappa)}.
 \label{coef:parameter-order}
\end{equation}
Here $K_{\mathrm{mom}}$ bounds the fixed moment orders in the Fourier
tail and complementary-coefficient truncation.  Make $\kappa_0$ smaller
if necessary so that the total carrier, safely bounded by
$\theta_*=599/1000$, satisfies $\theta_* /(1-\kappa_0)<3/5$ and all
three low-rank supports remain below one half.

Replacing $A$ by $U$ in \eqref{coef:three-monomials} leaves savings
$\zeta,2\zeta,9\zeta/16$.  After two reserves of $\zeta/16$, the
smallest saving is $7\zeta/16>15\sigma$.  Small-prime attachment enlarges
only group three; its costs $2\ell_\beta,2\ell_\beta,
(2+2\vartheta)\ell_\beta$ are included in the reserves.  Dyadic subdivision
therefore produces $Q_1,Q_2,Q_3$ satisfying all four size conditions of
Proposition~\ref{prop:current-medium}.  The strict margins absorb fixed
endpoint constants.  There are finitely many atoms, with fixed total
absolute scalar norm, before $N$ grows.  This establishes a uniform
three-factor representation for the actual finite coefficient families.

\hypertarget{coef:medium-checks}{}
\paragraph{Checks for the medium branch.}
The following checks connect Proposition~\ref{prop:current-medium} with
its application in Proposition~\ref{coef:prime-product-insertion}. They are
made after fixing the common finite policy and its strict support reserves.
\begingroup
\renewcommand{\LTcaptype}{}
\stepcounter{LT@tables}
\small
\setlength{\tabcolsep}{4pt}
\begin{longtable}{@{}>{\raggedright\arraybackslash}p{.21\textwidth}>{\raggedright\arraybackslash}p{.45\textwidth}>{\raggedright\arraybackslash}p{.28\textwidth}@{}}
\toprule
 Condition or input & Realization in the application & Verification \\
\midrule
\endhead
\bottomrule
\endfoot
Complete scale and residue & $x=UV$ is the scale of the complete product and $\mathfrak a=N$. The trim and local rescaling give $c_0x\le N\le x^{1+\kappa}$. & \eqref{coef:complete-scales}, \eqref{coef:parameter-order}\\[3pt]
Rough medium sequence & An individual interval carries the rough $1$ or M\"obius sequence, its bounded weight and permitted Mellin twist; the auxiliary-coprimality estimate is retained. & \eqref{coef:rough-SW} and its proof\\[3pt]
Complementary sequence & Unique factorization gives a fixed divisor bound. Any logarithmic truncation is charged by the mixed moment estimate before applying the local bound. & Lemma~\ref{coef:mixed-moment}; proof of Proposition~\ref{coef:prime-product-insertion}\\[3pt]
Three modulus factors & The normalized group convolutions are separate and one-bounded, with the residue-coprimality support retained. Marked multiplicities remain external scalars. & \eqref{coef:tensor-identity}; full and selected squares \eqref{coef:full-squares}--\eqref{coef:selected-squares}\\[3pt]
Four strict size tests & The three monomial bounds, the medium capacity, small-prime attachment and local rescaling yield all the tests in \eqref{eq:current-medium-ranges}. & \eqref{coef:three-monomials}; the reserves following \eqref{coef:parameter-order}\\[3pt]
Moving-residue phase and spectral input & The exact reciprocal phase is retained. Its variation and the packet and spectral norms use the stated local margins. & \eqref{eq:current-retained-phase}, \eqref{ana:phase-variation}; Appendix~\ref{app:classical}\\[3pt]
Uniform summation errors & Fixed atom scalars, Fourier norms and logarithmically many charts are covered by stronger local logarithmic accuracy. Prime powers, squarefree collisions and exact products equal to $N$ are counted separately. & \eqref{coef:collision}; proof of Proposition~\ref{coef:prime-product-insertion}\\[3pt]
\end{longtable}
\endgroup

For the four size tests, deducting the stated support reserves from the
savings $\zeta,2\zeta,9\zeta/16$ leaves strict margins that absorb fixed
dyadic endpoint multiples. This check concerns the medium branch; the
below-half estimate and Appendix~\ref{app:brun} cover the other branches of
Proposition~\ref{coef:prime-product-insertion}. Upper coefficients, lower
coefficients and ordinary marked descendants enter through that proposition
with their actual signed coefficient $\rho_N$.

\section{Distribution of the smooth prime products}
\label{app:prime-products}

This appendix proves \eqref{eq:current-PP} for the finite coefficient families of Appendix~\ref{app:coefficients}. The medium branches use Proposition~\ref{prop:current-medium}; the remaining branches use the below-half estimate or Appendix~\ref{app:brun}.

For a smooth weight $W_N$ on at most $k_*$ labelled prime slots, with
$p_i\ge cN^{a_*}$ and product $O(N)$, define
\[
 E_q(N;W)=\sum_{p_1,\ldots,p_k}W_N(p_1,\ldots,p_k)
 \left(\one_{p_1\cdots p_k\equiv N\ (q)}
       -\frac{\one_{(p_1\cdots p_k,q)=1}}{\varphi(q)}\right).
\]
The geometric weights have fixed polylogarithmic mixed Euler derivative
bounds.  They carry no additional arithmetic filters.

\subsection{Two auxiliary mean estimates}

Let $d_j$ denote the $j$-fold divisor function.
We record the elementary moment estimate with its dependence on the
moving shift, since logarithmic truncations require a fixed moment order.

\begin{lemma}\label{coef:mixed-moment}
Fix nonnegative integers $R,S$ and a constant $C\ge1$.  Put
$K_{\rm mom}=2^{2\max(R,S,1)}$ and $0\le\kappa\le1/(2K_{\rm mom})$.
Uniformly for $|a|\le Cx^{1+\kappa}$ and an interval
$J\subseteq[1,Cx]$,
\[
 \sum_{\substack{b\in J\\b\ne a}}
       \tau(b)^R\tau(|b-a|)^S
 \ll_{C,R,S} x(1+\log x)^{K_{\rm mom}-1}.
\]
\end{lemma}

\begin{proof}
If an interval $I\subseteq[1,Z]$ has length $Y$, choosing a largest
coordinate of an ordered $j$-tuple whose product belongs to $I$ gives
\begin{align*}
 \sum_{n\in I}d_j(n)
 &\le j\sum_{d\le Z^{1-1/j}}d_{j-1}(d)(Y/d+1)\\
 &\le j\{Y(1+\log Z)^{j-1}
              +Z^{1-1/j}(1+\log Z)^{j-2}\}.
\end{align*}
Indeed the product of the other $j-1$ coordinates is at most
$Z^{1-1/j}$; dropping the largest-coordinate restriction only increases
the sum.  The two final bounds follow by expanding the divisor function
as an ordered product and summing the last factor.
For every integer $e\ge0$,
$\tau(n)^e\le d_{2^e}(n)$.  On a prime power this follows by mapping
an $e$-tuple in $\{0,\ldots,v\}^e$ to the histogram of the subsets
of coordinates reaching each of the positions $1,\ldots,v$; the
histogram is a weak composition of $v$ into $2^e$ parts and determines
the tuple.

The nonzero values $|b-a|$ lie, with multiplicity at most two, in at most
two intervals of length $O_C(x)$ and upper endpoint $O_C(x^{1+\kappa})$.
Apply the preceding interval bound to $d_{K_{\rm mom}}$ to obtain
\[
 \sum_{\substack{b\in J\\b\ne a}}\tau(|b-a|)^{2S}
 \ll x(1+\log x)^{K_{\rm mom}-1}
       +x^{(1+\kappa)(1-1/K_{\rm mom})}(1+\log x)^{K_{\rm mom}-2}.
\]
The second exponent is at most $1-1/(2K_{\rm mom})-1/(2K_{\rm mom}^2)<1$.
The unshifted $2R$ moment has the same bound.
Cauchy--Schwarz proves the assertion.  The term $b=a$ is removed
before applying the divisor function.
\end{proof}

\begin{lemma}[A convolution estimate below one half]
\label{coef:low-rank-BV}
Fix $\eta>0$ and $c\ge1$. Let $x=UV$,
$x^\eta\le U,V\le x^{1-\eta}$, and let $\alpha,\beta$ be supported
on $[U,cU]$ and $[V,cV]$ with
\[
 \sum_m|\alpha(m)|^2\ll U(\log x)^C,\qquad
 \sum_n|\beta(n)|^2\ll V(\log x)^C.
\]
They are independent of the summation modulus.  Suppose that for every
fixed $A_1,B>0$, every $k\ge1$, $r\le(\log x)^B$ and $(b,r)=1$,
\begin{equation}
 \left|\sum_{\substack{m\equiv b\ (r)\\(m,k)=1}}\alpha(m)
 -\frac1{\varphi(r)}\sum_{(m,rk)=1}\alpha(m)\right|
 \ll_{A_1,B}\tau(k)^{C_0}U(\log x)^{-A_1},
 \label{coef:small-aux-SW}
\end{equation}
with one fixed exponent $C_0$.  For every $A>0$ there is a $D>0$ such
that $Q\le\sqrt x(\log x)^{-D}$ implies
\[
 \sum_{q\le Q}\max_{(a,q)=1}
 \left|\sum_{mn\equiv a\ (q)}\alpha(m)\beta(n)
 -\frac1{\varphi(q)}\sum_{(mn,q)=1}\alpha(m)\beta(n)\right|
 \ll_A x(\log x)^{-A}.
\]
\end{lemma}

\begin{proof}
Put $L=\log x$.  The primitive multiplicative large sieve
\cite[Theorem~18.2]{Kedlaya} is
\begin{equation}
 \sum_{r\le R}\frac r{\varphi(r)}\sum_{\chi\ (r)}^{*}
 \left|\sum_nz(n)\chi(n)\right|^2
 \le(H-1+R^2)\sum_n|z(n)|^2,
 \label{coef:multiplicative-LS}
\end{equation}
when $z$ is supported on an interval of at most $H$ consecutive integers.
Character orthogonality removes the principal character exactly.
A character modulo $q=rk$ induced by a primitive character $\chi^*$
satisfies $\chi(m)=\chi^*(m)\one_{(m,k)=1}$, including when $(r,k)>1$.
Put $\alpha_k=\alpha\one_{(m,k)=1}$ and
$\beta_k=\beta\one_{(n,k)=1}$ and write $A_k(\chi),B_k(\chi)$ for
their character sums.  Since $\varphi(rk)\ge\varphi(r)\varphi(k)$,
the required sum is bounded by
\begin{equation}
 \sum_{k\le Q}\frac1{\varphi(k)}
 \sum_{1<r\le Q/k}\frac1{\varphi(r)}
              \sum_{\chi\ (r)}^{*}|A_k(\chi)B_k(\chi)|.
 \label{coef:induced-characters}
\end{equation}
Both coprimality restrictions are essential.  The restrictions do not
increase the $L^2$ norms and, for fixed $k$, do not depend on $r$.

Let $R_0=L^{B_0}$.  For $R<r\le2R$, $R\ge R_0$, Cauchy--Schwarz
and \eqref{coef:multiplicative-LS} bound the inner sum by
\[
 \ll L^{C'}\left\{\frac xR+
 \sqrt x(\sqrt U+\sqrt V)+R\sqrt x\right\}.
\]
Summing dyadically, then over $k$ in
\eqref{coef:induced-characters}, gives
\begin{equation}
 \ll L^{C'}\left\{\frac{xL}{R_0}
       +L^2\sqrt x(\sqrt U+\sqrt V)+Q\sqrt x\right\}.
 \label{coef:large-conductors}
\end{equation}
Here $\sum_{k\le Y}1/\varphi(k)\ll1+\log Y$ and
$\sum_{k\ge1}1/(k\varphi(k))<\infty$; the latter sum occurs in
the last term.  Both follow from
$n/\varphi(n)=\sum_{d\mid n}\mu(d)^2/\varphi(d)$ and
$1/(d\varphi(d))\le\tau(d)/d^2$.
The middle term in \eqref{coef:large-conductors} has a fixed power
saving because $\sqrt x(\sqrt U+\sqrt V)\le2x^{1-\eta/2}$.
Taking $B_0,D$ sufficiently large controls the other two terms.

For $r\le R_0$, sum \eqref{coef:small-aux-SW} against the primitive
nonprincipal character.  Its mean vanishes, yielding
$|A_k(\chi)|\ll\varphi(r)\tau(k)^{C_0}UL^{-A_1}$.
The $L^1$ bound for $\beta_k$ is $O(VL^{C/2})$, so the small-conductor
part of \eqref{coef:induced-characters} is at most
\[
 \ll xL^{C/2-A_1}
       \left(\sum_{r\le R_0}\varphi(r)\right)
       \left(\sum_{k\le Q}\frac{\tau(k)^{C_0}}{\varphi(k)}\right).
\]
Choose an integer $J\ge C_0+1$.  Since
$n/\varphi(n)\le\tau(n)$ and $\tau(n)^J\le d_{2^J}(n)$,
the last sum is at most $(1+\log Q)^{2^J}$.
After $B_0$ is fixed, taking $A_1$ sufficiently large finishes the proof.
All divisor orders are fixed before the requested logarithmic accuracy.
\end{proof}

For an original prime slot, the classical prime Siegel--Walfisz theorem
\cite[Exercise~64]{TaoSW}, followed by partial summation, gives
\eqref{coef:small-aux-SW} at $k=1$ with any prescribed saving.
For arbitrary $k$ the change deletes only primes dividing $k$ and is at
most $2\|\alpha\|_\infty\omega_{\mathrm{pf}}(k)$; its fixed logarithmic norm is
absorbed by the positive-power scale $U$.  Thus $C_0=1$ works for all
the required exponents.  Fixed smooth profiles, interval endpoints and
Mellin twists with fixed polylogarithmic variation preserve the estimate.
For several original primes, put one such slot in $\alpha$ and the
remaining factors in $\beta$.  Unique factorization supplies a fixed
divisor bound and hence the required $L^2$ norm for $\beta$.

\subsection{Proof of the prime-product estimate}

\begin{proposition}\label{coef:prime-product-insertion}
For the finite upper, lower and marked coefficients constructed in
Appendix~\ref{app:coefficients}, the restricted medium estimate of
Proposition~\ref{prop:current-medium} and the below-half-level
Bombieri--Vinogradov estimates imply
\[
 \sum_q\rho_N(q)E_q(N;W)=O_{A,W}\bigl(N(\log N)^{-A}\bigr)
 \qquad(A>0).
\]
The coefficient in the count and exact coprime mean is the same.
The estimate is uniform for the fixed finite family and all sufficiently
large even $N$.
\end{proposition}

\begin{proof}
Choose a smooth $\eta$ equal to zero on $[0,1]$ and one on $[2,\infty)$,
and multiply $W_N$ by $\eta(p_1\cdots p_k/N^{1-\kappa_0})$.
The discarded products have size less than $2N^{1-\kappa_0}$, and are
strictly below $N$ for large $N$.  Their complete modulus-divisor sum
and exact mean have error $O(N^{1-\kappa_0+o(1)})$.
All Euler derivatives of this trim are bounded.

Insert a smooth nonnegative dyadic partition in each original prime.
On an original chart put $R_i=2^{j_i}$, $X=\prod_iR_i$ and $X^+=2^kX$.
With $f(v)=v^{1/(\log\log v)^3}$ fix the common roughness threshold
\[
 z=f(2^{6k_*+2}X^+)
\]
before applying a convolution identity.  Every original prime exceeds
$z$ eventually.  Replacing its prime indicator by $\Lambda(n)/\log n$
and deleting proper prime powers costs a fixed power saving: their
reciprocal mass above $cN^{a_*}$ is $N^{-a_*/2+o(1)}$.
Nonzero complete differences have divisor cost $N^{o(1)}$, and an exact
product $N$ contributes at most $N^{\theta_*+o(1)}$ separately.

For a group at scale $R$ let $Z=(2R)^{1/3}$ and
$\mu_Z=\mu\one_{[1,Z]}$.  If $\delta$ denotes the convolution identity,
$\delta-1*\mu_Z$ is supported above $Z$; consequently
$\Lambda*(\delta-1*\mu_Z)^{*3}=0$ below $2R$.  Expanding gives
\begin{equation}
 \Lambda=\sum_{j=1}^3(-1)^{j-1}\binom3j
             \mu_Z^{*j}*1^{*(j-1)}*\log.
 \label{coef:HB}
\end{equation}
The original smooth chart is retained.  The roughness projector factors
over all the variables in \eqref{coef:HB}; use $P^-(1)=\infty$.
There is no additional sharp coupled group boundary.

Each original group has at most six variables.  Individual rectangles
meeting its smooth support give complete local scales satisfying
\begin{equation}
 2^{-7k-1}X\le x\le2^{6k-1}X^+,
 \qquad f(x)\le z\le x^{2/(\log\log x)^3}.
 \label{coef:complete-scales}
\end{equation}
Thus $cN^{1-\kappa_0}\le x\le CN$ and
$c_0x\le N\le x^{1+\kappa}$ by \eqref{coef:parameter-order}.
The scale $x$ is always that of the complete product.

To separate the geometric profile, use coordinates
$u_j=\log(n_j/U_j)$ and a smooth compact extension $w$ of the profile.
Writing
$D_w=\sum_{S\subseteq\{1,\ldots,m\}}\|\partial_{2S}w\|_1$,
integration by parts gives
\[
 |\widehat w(\xi)|\le D_w\prod_j(1+\xi_j^2)^{-1},
 \quad \frac1{(2\pi)^m}\int|\widehat w(\xi)|\,d\xi
 \le\frac{D_w}{2^m}.
\]
The tail outside $|\xi|_\infty\le H$ is $O(mD_w/H)$.
Here $D_w$ is a fixed logarithmic power.  The arithmetic sums can be
interchanged with the absolutely integrable Fourier representation.
After the complete modulus sum, the tail at a nonzero difference
$b-N$ is bounded by fixed mixed divisor moments
\[
 \sum_{b\asymp x}\tau(b)^R\tau(|b-N|)^S
 \ll x(\log x)^C.
\]
Apply Lemma~\ref{coef:mixed-moment}; its fixed order is included in
$K_{\mathrm{mom}}$.  The exact
difference zero is estimated separately by $N^{\theta_*+o(1)}$.
Choose a sufficiently large logarithmic $H$.  For the retained modes
one pays the Fourier integral norm, rather than the number of modes.

Choose an individual medium slot, if present, in
$X^{3\zeta}\le U\le X^{1/3+\zeta/4}$.
Its local length lies in
$x^{2\zeta}\le U\le x^{1/3+\zeta/2}$ by
\eqref{coef:complete-scales}.  Its sequence has the form
\[
 \alpha_n=e_n\one_I(n)\one_{P^-(n)>z}h(n)n^{it},
 \qquad e_n=1\text{ or }\mu(n),\quad |h|\le1,
\]
with fixed polylogarithmic variation and frequency.  The auxiliary
Siegel--Walfisz property required in the medium estimate is
\begin{equation}
 \left|\sum_{\substack{n\equiv b\ (q)\\(n,k)=1}}\alpha_n
 -\frac1{\varphi(q)}\sum_{(n,kq)=1}\alpha_n\right|
 \ll_A U\tau(k)(\log x)^{-A}.
 \label{coef:rough-SW}
\end{equation}
For completeness, roughness can be inserted at even Brun order
$h=2\lfloor(\log\log x)^2\rfloor$.  Its next divisor support is
$z^{h+1}<U^{1/4}$ and its gap density has arbitrary logarithmic decay.
For squarefree $\ell$,
$\mu(\ell v)=\mu(\ell)\mu(v)\one_{(v,\ell)=1}$.
Put $V_\ell=U/\ell$ for the quotient scale. The identity
\[
 \mu(v)\one_{(v,\ell)=1}
 =\sum_{\substack{w\mid v\\p\mid w\Rightarrow p\mid\ell}}\mu(v/w)
\]
reduces the small-modulus calculation to the unrough
Siegel--Walfisz estimate for $\mu$.  Truncation at $w\le V_\ell^{1/4}$
leaves $O(V_\ell^{7/8}\tau(\ell)^2)$; the smallest inner scale is $U^{9/16}$.
After summing $\ell\le U^{1/4}$ this is
$O(U^{29/32}\log^4U)$.  Large moduli are bounded by the one-bounded
interval estimate.  Inserting $(n,k)=1$ deletes only primes $p\mid k$
above $z$ and costs at most
$2(U/z+1)\omega_{\mathrm{pf}}(k)$.  Partial summation absorbs the allowed variation
and twists by choosing a stronger initial logarithmic accuracy.  This proves \eqref{coef:rough-SW} for the actual slots.

The complementary convolution is bounded by a fixed divisor power.
If necessary, truncate only this complementary coefficient at a fixed
logarithmic threshold, using the same mixed divisor moment bound for
nonzero differences and $N^{\theta_*+o(1)}$ at zero.  The threshold is
chosen after the fixed orders, and the stronger local logarithmic
estimate restores its normalization.  The medium slot is not truncated.
Now apply Proposition~\ref{prop:current-medium} to the exact tensor
representation \eqref{coef:tensor-identity}.

If no medium slot occurs, every original group contains a long
slot: since $\zeta<a_*/100$, six slots below $X^{3\zeta}$ cannot
have product comparable to $R_i$. There are at most two long slots,
because three would exceed the complete scale. A truncated M\"obius
slot is $O(X^{1/3})$ and cannot be long. Thus one or two long
unweighted integer variables remain, with their logarithms, smooth
twists, individual intervals, and common roughness projector.
Their complementary short products are bounded by $x^{20\zeta}$
and $x^{40\zeta}$, respectively. Appendix~\ref{app:brun} estimates
both cases with the exact coprime mean.

Finally, the low-rank coefficient pieces have complete support with a
fixed gap below $x^{1/2}$.  For one original prime use maximal prime
Bombieri--Vinogradov.  For two or more original primes, apply
Lemma~\ref{coef:low-rank-BV} using one original prime slot as explained
above.  The bounded total coefficient norm permits its absolute
form.  All fixed atom scalars, logarithmically many charts and the
Fourier integral norm are paid by requesting stronger local logarithmic
accuracy.  The removed prime powers, collisions and trim have fixed
power savings.  The finite family therefore has one eventual threshold
for the asserted estimate.
\end{proof}

\section{One and two long rough variables}
\label{app:brun}

The estimates here supply the branches without a medium slot in the proof of Proposition~\ref{coef:prime-product-insertion}.

For the one- and two-long-variable cases we use direct progression
estimates and the classical complete Kloosterman bound.
The argument keeps the exact coprime mean and is uniform in the
size of the primitive residue.

\subsection{Brun polynomials with a uniform density gap}

Let $L=\log\log x$, $h=2\lfloor L^2\rfloor$, $r=h+1$ and
$\log z\le2\log x/L^3$.  Put $P(z)=\prod_{p\le z}p$ and
$y=z^r=x^{o(1)}$.  Define
\[
 \lambda^+(d)=\mu(d)\one_{d\mid P(z),\,\omega_{\mathrm{pf}}(d)\le h},
 \qquad
 \lambda^-(d)=\mu(d)\one_{d\mid P(z),\,\omega_{\mathrm{pf}}(d)\le h+1},
\]
and
\[
 U_z(n)=\sum_{d\mid n}\lambda^+(d),\quad
 L_z(n)=\sum_{d\mid n}\lambda^-(d),\quad
 G_z=U_z-L_z,\quad r_z(n)=\one_{(n,P(z))=1}.
\]
Both divisor sums have coefficients of absolute value at most one and
support below $y$.  The alternating binomial identity gives
\begin{equation}
 L_z\le r_z\le U_z,\qquad U_z\ge0,\qquad
 G_z(n)=\binom{\omega_{\mathrm{pf}}((n,P(z)))}{r}\ge0.
 \label{coef:Brun-order}
\end{equation}
This also covers $(n,P(z))=1$, when both polynomials are one.
The lower polynomial need not be nonnegative.

For every $q$ set
\[
 A_q=\sum_{(d,q)=1}\frac{\lambda^+(d)}d,\quad
 g_q=\sum_{(d,q)=1}\frac{\lambda^+(d)-\lambda^-(d)}d,\quad
 V_q=\prod_{\substack{p\le z\\p\nmid q}}\left(1-\frac1p\right).
\]
Finite inclusion-exclusion, applied to independent Bernoulli variables
of means $1/p$, gives
\begin{equation}
 0\le A_q-V_q\le g_q
 \le\frac{(\sum_{p\le z}1/p)^r}{r!}=:\epsilon_x.
 \label{coef:Brun-gap}
\end{equation}
The bound is uniform in the set of primes dividing $q$.  The elementary
prime harmonic bound $\sum_{p\le z}1/p\ll L$ and $r!\ge(r/e)^r$ yield
$\epsilon_x\le\exp(-cL^2\log L)$ eventually.  Thus $\epsilon_x$
is smaller than every fixed negative power of $\log x$, whereas
$y=x^{o(1)}$ will occur only beside a fixed power saving.

\subsection{Smooth progression cores}

For a compactly supported smooth profile define
$N_1(f)=\|f\|_1+\|f'\|_1$ and
$N_2(f)=\|f\|_1+\|f''\|_1$.
Progression summation and M\"obius inversion give
\begin{align}
 \sum_{n\equiv b\ (q)}f(n/U)&=\frac Uq\int f+O(N_1(f)),\notag\\
 \sum_{(n,q)=1}f(n/U)&=U\frac{\varphi(q)}q\int f
                        +O(\tau(q)N_1(f)).
 \label{coef:one-core}
\end{align}
These estimates are uniform in $U>0$ and in the primitive class $b$.
For a fixed $m$ coprime to $q$, the discrepancy of $mn\equiv a\pmod q$
from its exact coprime mean is therefore
$O((1+\tau(q)/\varphi(q))N_1(f))$.  If $(m,q)>1$ and $(a,q)=1$,
both terms vanish.

For two variables, use $e(t)=e^{2\pi it}$ and
$\widehat f(t)=\int f(s)e(-ts)\,ds$.
Let $c_q$ be the Ramanujan sum and
$S(j,k;q)=\sum_{u\bmod q}^{*}e((ju+k\bar u)/q)$.
For $b=a\bar m\pmod q$, double periodization gives the exact
discrepancy
\begin{equation}
 \frac{UV}{q^2}\sum_{j,k\in\mathbb Z}
 \widehat f(jU/q)\widehat g(kV/q)
 \left(S(bj,k;q)-\frac{c_q(j)c_q(k)}{\varphi(q)}\right).
 \label{coef:two-Poisson}
\end{equation}
Both axes vanish identically, because
$S(0,k;q)=c_q(k)$, $S(bj,0;q)=c_q(j)$ and $c_q(0)=\varphi(q)$.
For $jk\ne0$, use
\[
 |S(bj,k;q)|\le\tau(q)\sqrt q\,(j,k,q)^{1/2},\qquad
 |\widehat f(t)|\ll N_2(f)(1+|t|)^{-2}.
\]
For every $H>0$ and $d\ge1$,
$\sum_{j\ne0,\,d\mid j}(1+|j|/H)^{-2}\le2H/d$.
Expanding common divisors and using
$\sum_{d\mid q}d^{-3/2}\le\zeta(3/2)$ therefore bounds
\eqref{coef:two-Poisson} by
\begin{equation}
 N_2(f)N_2(g)
 \left(\sqrt q\,\tau(q)+\frac{\tau(q)^2}{\varphi(q)}\right).
 \label{coef:two-core-bound}
\end{equation}
For the Ramanujan term one uses
$|c_q(j)|\le(j,q)$ and
$\sum_{d\mid q}\varphi(d)/d\le\tau(q)$.
When $q=1$ the discrepancy is zero.

Let a short coefficient $\alpha_m$ have
$\sum|\alpha_m|\ll M\log^C x$, with $x\asymp MU$ or $MUV$.
If $q\le D$ and $|\rho(q)|\le\tau(q)^J\log^B x$ with fixed $J,B$,
the sums of the one- and two-variable core errors are
\begin{equation}
 O(MD\log^{C'}x),\qquad O(MD^{3/2}\log^{C'}x),
 \label{coef:core-summed}
\end{equation}
respectively.  On inserting Brun divisors $d_1,d_2\le y$, the profiles
are at scales $U/d_1,V/d_2$ and their norms in
\eqref{coef:two-core-bound} are unchanged.  Replace $m$ by $md_1d_2$
only in the primitive residue.  The fixed-product bound is uniform in
that residue, so summing the divisor pairs costs at most $y^2$.

\subsection{Restoring roughness with the exact mean}

For a nonnegative smooth $F$, expansion of the divisor coefficients in
\eqref{coef:one-core} gives
\begin{align}
 \sum_{(n,q)=1}F(n/U)U_z(n)
 &=U\frac{\varphi(q)}q A_q\int F+O(y\tau(q)N_1(F)),\notag\\
 \sum_{(n,q)=1}F(n/U)G_z(n)
 &=U\frac{\varphi(q)}q g_q\int F+O(y\tau(q)N_1(F)).
 \label{coef:exact-Brun-means}
\end{align}
Both left sides are nonnegative.  They are exact coprime sums;
the factor $\varphi(q)/q$ is retained throughout.

For two variables write $U_i=U_z(n_i)$, $r_i=r_z(n_i)$ and
$G_i=G_z(n_i)$.  The valid product comparison is
\begin{equation}
 0\le U_1U_2-r_1r_2\le G_1U_2+U_1G_2.
 \label{coef:rough-product-gap}
\end{equation}
Indeed $U_1U_2-r_1r_2=(U_1-r_1)U_2+r_1(U_2-r_2)$.
This argument does not multiply two possibly negative lower polynomials.
For complex profiles choose nonnegative smooth majorants $F\ge|f|$,
$G\ge|g|$ and use $|\alpha|$ only in estimating the absolute gap.
If $C_q$ denotes the congruence count and $M_q$ its exact coprime mean,
and $B=G_1U_2+U_1G_2$, then
\begin{align*}
 |(C_q-M_q)(\alpha fg(r_1r_2-U_1U_2))|
 &\le C_q(|\alpha|FGB)+M_q(|\alpha|FGB)\\
 &\le |(C_q-M_q)(|\alpha|FGB)|+2M_q(|\alpha|FGB).
\end{align*}
The first term is covered by expanding the Brun divisors and
\eqref{coef:core-summed}; the second is estimated by
\eqref{coef:exact-Brun-means}.  Fixed divisor/totient mean bounds give
the total two-variable error
\begin{equation}
 \ll\log^{C'}x\left\{
 MUV\epsilon_x+M(U+V)y+My^2+My^2D^{3/2}\right\}.
 \label{coef:two-rough-error}
\end{equation}
The one-variable version is
\begin{equation}
 \ll\log^{C'}x\{MU\epsilon_x+My+MyD\}.
 \label{coef:one-rough-error}
\end{equation}
Only fixed logarithmic factors multiply $\epsilon_x$.
For the long alternatives, $M\le x^{20\zeta}$ or $x^{40\zeta}$,
$D\le x^{\Theta_{\mathrm{loc}}}$ with $\Theta_{\mathrm{loc}}<3/5$,
and $\min(U,V)\ge x^\nu$ for a fixed $\nu>0$.
Choose $\zeta$ so that
\[
 20\zeta+\Theta_{\mathrm{loc}}<1,\qquad
 40\zeta+\frac32\Theta_{\mathrm{loc}}<1.
\]
Every remaining term in \eqref{coef:two-rough-error} and
\eqref{coef:one-rough-error} then has a fixed power saving.
The errors are consequently $O_A(x\log^{-A}x)$ for every fixed $A$.

\subsection{Individual sharp intervals}

Let $\delta=(\log x)^{-B}$.  An interval in a fixed positive scale has
a smooth approximation whose error is bounded by nonnegative smooth
strips of normalized width $O(\delta)$ at its endpoints.  If the interval
is shorter than $\delta$, a single strip covers it.  Their integrals
are $O(\delta)$, their first-derivative norms are $O(1)$, and their
second-derivative norms are bounded by a fixed power of $\delta^{-1}$.
Multiply by the fixed envelopes for the other profiles.

Apply the proved rough smooth estimates to these positive strips.
Their exact means are bounded by
\[
 \ll\log^{C'}x\{x\delta+M(U+V)y+My^2\}
\]
in the two-variable case, and by
$\log^{C'}x\{x\delta+My\}$ in the one-variable case.
The exponent in the first term uses the strip integral, not its
derivatives, and is independent of $B$.  The $B$-dependent derivative
losses affect only the fixed-power errors.  Choose $B$ after the desired
logarithmic accuracy.  This controls both the count and the exact mean,
uniformly in every primitive residue, including a residue equal to an
intermediate complete product.  It proves the required extension to
sharp individual intervals.  No assertion about arbitrary sharp coupled
domains is used.

\section{Causal upper sieve policies}\label{app:upper}

This appendix gives the upper-policy density and coefficient realization used in Section~\ref{sec:current-upper}; its finite savings are evaluated in Appendix~\ref{app:numeric}.

We construct the divisor coefficients from finitely many decisions
on prime prefixes. Prefixes are ordered by decreasing primes:
if $p_i=N^{t_i}$, then $t_1>\cdots>t_r>0$.
Write $S_j=t_1+\cdots+t_j$ and $S_0=0$.
A parent is obtained by deleting the last prime of a prefix.
Every decision at that parent is retained when smaller primes
are appended.

\subsection{Admissible caps and the parity of first exits}

For a nonempty subset $I\subseteq\{1,\ldots,r\}$ write
\begin{equation}\label{sv:selected-gate}
 b_I=\sum_{i\in I}t_i,\qquad
 G_I(t)=\max_{\substack{j\le r\\j\notin I}}
 \left(b_I+\sum_{\substack{i<j\\i\notin I}}t_i+2t_j\right),
 \qquad \max\varnothing=0.
\end{equation}
The full-square support condition is
\begin{equation}\label{sv:full-gate}
 G_{\mathrm{full}}(t)=\max_{1\le j\le r}(S_{j-1}+2t_j).
\end{equation}
The candidate selected-product endpoint is
\begin{equation}\label{sv:tent}
 v(b)=\min\left\{\frac{1+b}{2},\frac{57}{64}-\frac{3b}{2}\right\}.
\end{equation}
A selected branch on a compact prefix region is admitted only when
its constant cap $\lambda$ satisfies, with a fixed positive margin,
\begin{equation}\label{sv:branch-licence}
 \max\{1/2,G_{\mathrm{full}}(t),G_I(t)\}
       <\lambda<v(b_I)
\end{equation}
throughout that region.  The separate three-slot branch with endpoint
$\Lambda_{\rm cap}-3t_3/8$ is subject to the whole-box factorization conditions
of Appendix~\ref{app:coefficients}.  A maximum of candidate endpoints
is always taken over branches satisfying their support conditions.

An ordinary upper residual tree at total exponent $\lambda$, after
a prefix of rank $r$, applies the cube tests
\begin{equation}\label{sv:upper-cubes}
 S_{j-1}+3t_j<\lambda\qquad(j>r,\ j\text{ odd}).
\end{equation}
If $r$ is even, these are the ordinary upper residual tests.  If $r$
is odd, they are the ordinary lower residual tests.  The earlier
membership decisions are retained.  At an even prefix a change of cap
is made after admitting the prefix; it introduces no exclusion at that
rank.  At an odd prefix a cap may be chosen before its own odd test,
provided that the decision depends only on that prefix and the
previously fixed policy.  Either operation leaves every first
exclusion at odd full rank.

\begin{lemma}[First-exit identity]\label{sv:first-exit}
Let $\mathcal T$ be a prefix-closed collection of squarefree products
from a finite prime set $\mathcal P$, containing the empty product.
Suppose that all first excluded prefixes have odd rank.  For every
multiplicative $\xi$ with $0\le \xi(p)\le1$,
\begin{equation}\label{sv:upper-exit}
 \sum_{d\in\mathcal T}\mu(d)\xi(d)
 =\prod_{p\in\mathcal P}(1-\xi(p))
 +\sum_{e\text{ first exit}}\xi(e)
           \prod_{\substack{p\in\mathcal P\\p<p_{\mathrm{last}}(e)}}
                            (1-\xi(p)).
\end{equation}
In particular $\sum_{d\mid n,\ d\in\mathcal T}\mu(d)$ majorizes
$\one_{(n,\prod_{p\in\mathcal P}p)=1}$ for every integer $n$.
\end{lemma}

\begin{proof}
Each excluded divisor has a unique first excluded prefix $e$; its
remaining prime factors are smaller than $p_{\mathrm{last}}(e)$.
Summing the complete continuation gives
$\mu(e)\xi(e)\prod_{p<p_{\mathrm{last}}(e)}(1-\xi(p))$.
Subtract these disjoint continuations from the complete Euler product.
Since $\mu(e)=-1$, the sign in \eqref{sv:upper-exit} is positive.
Take $\xi(p)=\one_{p\mid n}$ for the last assertion.  The same
finite identity applies at $n=0$ and at negative integers.
\end{proof}

The square tests persist through descendants.  A retained odd cube
implies its own square.  The following even square follows from
ordering:
\begin{equation}\label{sv:descendant-square}
 S_{j-1}+2t_j\le S_{j-2}+3t_{j-1}<\lambda.
\end{equation}
After all selected slots have appeared, the expression in
\eqref{sv:selected-gate} for any appended slot equals its full square.
Consequently \eqref{sv:branch-licence}, checked at the selected prefix,
continues to control the complete descendant vector.

\subsection{The cap changes and their density differences}

Set
\[
 \rho_0=\frac{19}{96},\qquad \mu_0=\frac{25}{128},\qquad
 \Theta=\frac{299}{500},\qquad \Lambda_{\rm cap}=\frac{153}{256}.
\]
The first prime is less than $N^{\rho_0}$.  The three-slot endpoint
below $t_1=\mu_0$ is the envelope of admissible branches of
\begin{equation}\label{sv:third-envelope}
 v(t_1),\ v(t_1+t_2),\ v(t_1+t_3),\ v(t_2+t_3),\ v(S_3),\
                         \Lambda_{\rm cap}-\frac{3t_3}{8};
\end{equation}
above $\mu_0$ the initial three-slot endpoint is $v(t_1)$.
Every branch in \eqref{sv:third-envelope} retains its individual
factorization conditions.  In particular the displayed list is not
an unconditional permission to select any nonleading pair.

At rank four the selected candidates are $M=v(b_I)$ with
$4\in I$, satisfying \eqref{sv:branch-licence} and exceeding the
inherited third cap.  Such a change is made only on prefixes that
already pass the third cube.  Membership through rank four is
unchanged and subsequent odd cubes use the new cap.
The limiting density difference on a region $\mathcal C$ is
\begin{equation}\label{sv:even-gain}
 4\int_{\mathcal C}
 \frac{\overline F((L-S_r)/t_r)
                   -\overline F((M-S_r)/t_r)}
      {t_1\cdots t_{r-1}t_r^2}\,dt_1\cdots dt_r,
 \qquad \overline F=\frac{F}{2e^\gamma},
\end{equation}
where $r=4$, $L$ is the inherited cap and $M$ the selected endpoint.
Indeed an even prefix $P_r$ contributes
$g_N(P_r)V_N^+(N^\lambda/P_r,p_r)$ to its complete subtree.
Subtracting the two subtrees gives \eqref{sv:even-gain}; it is an
identity of raw densities before taking an asymptotic limit.

At rank five a candidate $v(b_I)$, $5\in I$, is compared with the
inherited cap and the fifth cube $S_4+3t_5$.  Its cap is fixed before
that cube.  Thus a fifth prefix rejected at the smaller cap may be
retained at the larger one; all earlier decisions remain fixed.
The negative fifth subtree has the ordinary lower residual, and its
saving is
\begin{equation}\label{sv:odd-gain}
 4\int_{\mathcal C}
 \frac{\overline f((M-S_5)/t_5)
                   -\overline f((L-S_5)/t_5)}
      {t_1t_2t_3t_4t_5^2}\,dt_1\cdots dt_5,
 \qquad \overline f=\frac{f}{2e^\gamma},
\end{equation}
with the zero extension of $\overline f$ below $2$.
The inherited first-five cap is bounded above by
\begin{equation}\label{sv:fifth-envelope}
 R_5(t)=\max\left\{\Lambda_{\rm cap}-\frac{3t_3}{8},\
         v(b_I):\varnothing\ne I\subseteq\{1,\ldots,5\}\right\}.
\end{equation}
Using an upper bound for the inherited cap in a positive saving is
conservative.  It does not change the physical cap or its density.

The rank-six changes are even-prefix changes of the form
\eqref{sv:even-gain}.  They are applied after preserving the third
and fifth admission tests.  In particular
$L_3-S_6>2t_5-t_6>t_6$ on a region whose fifth cube is controlled by
$L_3$.  A convex piecewise-affine majorant $U_5\ge R_5$ can be obtained by
choosing one defining affine upper plane for each tent in
\eqref{sv:fifth-envelope}.  For a new concave endpoint $M$, put
$d=M-U_5$ and restrict the region to $M-S_6\le3t_6$.
Write $L$ for the inherited cap.  Where $d>0$, the retained fifth
cube gives
\[
 1<\frac{L-S_6}{t_6}\le\frac{U_5-S_6}{t_6}
       <\frac{M-S_6}{t_6}\le3.
\]
Thus $\overline F(s)=1/s$ throughout the interval of comparison, and
\[
 \overline F((L-S_6)/t_6)-\overline F((M-S_6)/t_6)
 \ge \frac{t_6(M-U_5)}{(U_5-S_6)(M-S_6)}.
\]
The cap is changed only where $d>0$, giving the saving lower bound
\begin{equation}\label{sv:sixth-lower}
 \frac{4d_+}{R_oR_n\prod_{i=1}^6t_i},
\end{equation}
where $R_n\ge M-S_6$ and
$R_o\ge\min\{U_5-S_6,M-S_6\}>0$ are whole-region upper bounds.
The positive part is taken before a lower bound for the positive
prime-density measure is extracted.  For a rectangular region,
concavity gives
$\int d_+/\operatorname{Vol}(\mathcal C)
\ge 2^{-6}\sum_{v\text{ vertex}}d(v)$.

Two additional changes occupy disjoint parts of $\mu_0<x<\rho_0$.
For the rank-two change, put $a(x)=2v(x)-1=25/32-3x$ and use
\[
 a(x)<y<x,\qquad L=v(x),\quad M=v(y).
\]
The support conditions follow from
$2x,x+2y,2x+y< v(x)<v(y)$; the selected slot is $2$.
The saving is the $r=2$ instance of \eqref{sv:even-gain}.
For the rank-three change, put $t=y+z$ and impose
\[
 a(x)<t<x,\quad z=tw/2,\quad 5/9\le w<1,
 \qquad L=v(x),\quad M=v(t).
\]
Here $y<3/16$, whereas the rank-two change has $y\ge3/16$.
Selecting slots $\{2,3\}$ gives the extra support condition $2x+t<M$.
The third node is retained and its lower residual is changed, giving
the $r=3$ analogue of \eqref{sv:odd-gain}.  Its arguments lie in
$[2,4]$, where $\overline f(s)=\log(s-1)/s$.

At rank eight a selected endpoint $M$ is compared with a valid
upper bound $R_6$ for the inherited sixth cap.  The retained seventh
cube and the additional condition $M-S_8<3t_8$ give
\begin{equation}\label{sv:eighth-range}
 1<\frac{\lambda_6-S_8}{t_8}
 \le\frac{R_6-S_8}{t_8}<\frac{M-S_8}{t_8}<3.
\end{equation}
Consequently \eqref{sv:even-gain} is bounded below by the integral of
\begin{equation}\label{sv:eighth-density}
 \frac{4(M-R_6)}{(R_6-S_8)(M-S_8)\prod_{i=1}^8t_i}.
\end{equation}
Each chosen region must retain every full and selected-complement
square, as well as all ancestor tests.  Regions acted upon at one
stage have disjoint interiors.  Changes at different stages are
combined by the exact telescope
$V_3-V_6=(V_3-V_4)+(V_4-V_5)+(V_5-V_6)$.
For each density difference we use one lower enclosure.

\subsection{Literal ordinary densities and finite realization}

Let $g_N(p)=\one_{p>2,\ p\nmid N}/(p-1)$ and
$V_N(y)=\prod_{p<y}(1-g_N(p))$.  The ordinary cubic Rosser
coefficients satisfy, uniformly on compact subsets of the stated
ranges,
\begin{equation}\label{sv:literal-density}
 \frac{V_N^+(N^\beta,N^t)}{V_N(N^t)}=F(\beta/t)+o(1)
       \quad(\beta/t>1),\qquad
 \frac{V_N^-(N^\beta,N^t)}{V_N(N^t)}=f(\beta/t)+o(1)
       \quad(\beta/t\ge2).
\end{equation}
Here the densities are those of the literal coefficients, not merely
upper and lower bounds for an arbitrary sifted sequence.

To see the uniformity in $N$, compare a normalized fixed-rank exit
with the corresponding exit for $g_0(p)=\one_{p>2}/(p-1)$.
A surviving prefix gains the multiplier
\[
 \prod_{\substack{p_{\rm last}\le p<y\\p\mid N}}(1-g_0(p))\le1;
\]
other prefixes disappear.  A rank-$j$ exit has last prime at least
$N^{\beta_{\min}/(3\cdot2^{j-1})}$ on a positive compact parameter
range.  Deleting prime divisors of $N$ therefore changes any
fixed-rank integral by $o(1)$ uniformly.  The positive Rosser exit
series converges uniformly on compact parameter ranges, and the
raw one-sided bounds give the opposite side of a finite-rank
sandwich; see \cite[Proposition~11.8 and Lemma~14.7]{Suzuki}.
First choose the rank and then let $N$ increase.  This proves
\eqref{sv:literal-density}.  For $1<s\le3$ the identity at the first
cube reduces the upper calculation to cutoff $N^{\beta/3}$, giving
$F(s)=2e^\gamma/s$.

The normalization follows from Mertens' formula.  In particular the
exact ratio to $2C_N\prod_{p<z}(1-1/p)$ is
\begin{equation}\label{sv:euler-ratio}
 \frac{V_N(z)}{2C_N\prod_{p<z}(1-1/p)}
 =\prod_{p\ge z}\left(1-\frac1{(p-1)^2}\right)^{-1}
  \prod_{\substack{p\mid N\\p\ge z}}\left(1-\frac1{p-1}\right).
\end{equation}
Its logarithm is $O((1+\log N/\log z)/(z-2))$.
Thus $V_N(N^t)\sim2e^{-\gamma}C_N/(t\log N)$ uniformly in $N$.
The factor $4$ in \eqref{sv:even-gain} and \eqref{sv:odd-gain}
is exactly $2e^{-\gamma}\cdot2e^\gamma$.

The finite cell data consist of rational regions, their ranks,
selected subsets, endpoints, and parent labels. Every active region
has positive margins in its admission and support inequalities;
an inherited region keeps its parent's decision. The density is
the sum of the exact subtree differences. For a positive saving,
any selected subfamily gives a lower bound. For a negative cost,
the upper bound must include the whole retained domain.

\begin{proposition}[Finite-policy realization]\label{sv:finite-realization}
Fix such a finite admissible policy and a positive error $\epsilon$.
Assume the coefficient factorization conditions of
Appendix~\ref{app:coefficients} on its selected branches.
There is a fixed finite refinement and a bounded squarefree
coefficient $\rho_N^{\rm up}$, supported on $(q,N)=1$, with the following
properties.  Its divisor sum is a nonnegative rough-indicator
majorant.  The signed distribution estimate \eqref{eq:current-PP}
applies to that same coefficient.  Its density differs, in the
required upper direction, from the ideal policy density by at most
$\epsilon C_N/\log N+o(C_N/\log N)$.
\end{proposition}

\begin{proof}
First truncate any small-last-prime tail.  At the third activation
one may use, with $u=\rho_0/2$ and $0<h<u$,
\begin{equation}\label{sv:activation-tail}
 T_3(h)\le\frac{18h}{(u-h)^2}
   \left(\log^2\frac{\rho_0}{h}
                   +2\log\frac{\rho_0}{h}+2\right)=o(1).
\end{equation}
This follows from $\overline F(s)-\overline f(s)\le9/s^2$
for $s\ge3$, the residual lower bound $\lambda-t_1-t_2\ge u$,
and the ordered reciprocal-prime triangle of mass
$\tfrac12\log^2(\rho_0/t_3)$.  Higher selected changes have
positive last-prime cutoffs.  Proper affine walls have arbitrarily
small strip volume.  Identities and open inherit regions are
retained.  On the remaining compact pieces choose coarse caps with
strict margins and refine them in increasing prefix order.  No
ancestor decision may depend on an appended prime.

The third endpoint is piecewise affine and Lipschitz.  For boxes
of width $\chi<e/10$, caps of the form
$\lambda_{\mathscr B}=\Gamma_{\mathscr B}-2e$ approximate it with error at most
$4e+6\chi$.  Near a three-slot branch with endpoint
$W=\Lambda_{\rm cap}-3t_3/8$, an auxiliary factorization cover may exceed
the physical cap.  If $\inf_{\mathscr B} A\ge\lambda_{\mathscr B}+e/2$, where
$A=(1+t_1)/2$, take $\ell_{\rm aux}=\lambda_{\mathscr B}+e/4$.
Otherwise $\sup_{\mathscr B} A<\lambda_{\mathscr B}+11e/20$, and
$\ell_{\rm aux}=\lambda_{\mathscr B}+e$ lies at least $9e/20$ above the
necessary support condition and $e$ below $W$.  The other initialization square
has margin $3/256$.  The auxiliary cover gives a factorization for the coefficient;
it never replaces the physical cap in a density or residual level.

Choose a presieve $z_0=N^{\eta_0}$.  At even rough rank use outer
endpoint tests and at odd rough rank use inner tests, retaining all
ancestor labels.  This gives $a_N(d)\ge\lambda_N(d)$ for the
raw upper coefficients.  If $B^-\le I_0\le B^+$ are the
dimension-one small-prime beta polynomials, attach $B^+$ to
positive $a_N(d)$ and $B^-$ to negative $a_N(d)$.  Term by term,
$a_N(d)B_{\rm sign}\ge a_N(d)I_0$.  Unique small/rough
factorization defines the resulting $\rho_N^{\rm up}$ and preserves
squarefreeness and boundedness.

The density direction needs an additional identity.  Projecting
\eqref{sv:upper-exit} to primes at least $z_0$ gives
\begin{equation}\label{sv:upper-projection}
 V_N(z_0)\sum_{d\text{ rough}}\lambda_N(d)g_N(d)
       \le\sum_{d\text{ full}}\lambda_N(d)g_N(d).
\end{equation}
Only positive exits with last prime at least $z_0$ survive on the
left.  For a fixed coarse policy the absolute beta and fine-box
errors are
$O(e^{-c/\eta}\eta_0^{-2}+C_{\rm policy}\tau\eta_0^{-5})$
relative to the root Euler product.  Choose $\eta,\eta_0$ and
then $\tau$ so that these errors and the truncated pieces fit
within $\epsilon$.

All these choices precede $N$.  The fine rank is bounded by
$J_0>\theta_*/\eta_0+2$; support inflation can be kept below
$1/1000$, giving $\theta_*=599/1000$.
Finally choose the fixed moment order and complete-scale trim so
that $\theta_*/(1-\kappa_0)<3/5$ and all local inequalities have
strict margins.  Appendix~\ref{app:coefficients} then gives
\eqref{eq:current-PP} for this very coefficient.  The pointwise
bound follows from Lemma~\ref{sv:first-exit}; the density bound
follows from \eqref{sv:upper-projection} and the absolute errors.
\end{proof}

\section{Lower policies and finite feedback}\label{app:lower}

This appendix supplies the lower-policy realizations and positive marked transition inequalities used in Section~\ref{sec:current-lower}. Appendix~\ref{app:numeric} evaluates the corresponding finite maps and integrals.

\subsection{The ordinary lower tree and fourth-prefix restoration}

Put ${a_{\rm s}}=417/5000$, ${h_{\rm s}}=19/125$ and $\nu=249/448$.
For either $w={a_{\rm s}}$ or $w={h_{\rm s}}$, begin with the ordinary lower support
at total exponent $\nu_\delta=\nu-\delta$, where
$0<\delta$ is fixed and $2{h_{\rm s}}<\nu_\delta$.
Its first square $2t_1<\nu_\delta$ is automatic below $w$;
its subsequent tests are
\begin{equation}\label{sv:lower-cubes}
 S_{j-1}+3t_j<\nu_\delta\qquad(j\text{ even}).
\end{equation}
At an admitted third prefix, freeze $\lambda_3\ge\nu_\delta$
for future even cubes.  An ordinary fourth change is permitted only
after $S_3+3t_4<\lambda_3$; it freezes
$\lambda_4\ge\lambda_3$.  The fifth prefix is then admitted, and
may freeze $\lambda_5\ge\lambda_4$ for future even tests.
Every active selected cap is below its admissible endpoint and above
its parent and all support conditions \eqref{sv:selected-gate}--\eqref{sv:full-gate}.
For the lower starting families the selected endpoint is
\begin{equation}\label{sv:lower-endpoint}
 v_-(b)=\min\left\{\frac{1+b}{2},
 \frac{32-34b}{39},\frac{57}{64}-\frac{3b}{2}\right\}.
\end{equation}
The selected branch is kept through descendants; the residual size
ratio of a marked state, introduced below, is separate from the
local residue exponent in the distribution estimate.

There is also a fourth-prefix restoration.  On an otherwise first
excluded fourth prefix choose $\lambda_4$ before its own cube,
subject to
\begin{equation}\label{sv:restore-gates}
 \lambda_3\le S_3+3t_4<\lambda_4,
 \qquad \max\{G_{\rm full},G_I\}<\lambda_4<v_-(b_I).
\end{equation}
Attach the ordinary lower residual below $p_4$ at level
$N^{\lambda_4}/P_4$.  Since
$\lambda_4-S_4>2t_4$, its first square cannot fail.
Its later first exits have even residual rank, hence even full
rank.  The restored tree remains a lower tree.  Its density gain is
exactly
\begin{equation}\label{sv:restoration-density}
 4\int_{\mathcal B}
 \frac{\overline f((\lambda_4-S_4)/t_4)}
      {t_1t_2t_3t_4^2}\,dt_1dt_2dt_3dt_4.
\end{equation}
There was no retained residual at such a first exit, so no parent
residual term is subtracted in \eqref{sv:restoration-density}.

Define a finite rational family of restoration regions as follows. Let
$r_0=(\nu-{h_{\rm s}})/9$, $n\in\{400,800\}$, and
$\mathscr B=\prod_{i=1}^4[\ell_i,\ell_i+1/n]$ with
$r_0\le\ell_4<\ell_3<\ell_2<\ell_1<{h_{\rm s}}$ and
$u_1+3u_2<\nu$, where $u_i=\ell_i+1/n$.
Let $U_3$ be the maximum of $\nu$, the upper bounds of all
three-slot selected endpoints, and $\Lambda_{\rm cap}-3\ell_3/8$.
It bounds the inherited third cap from above.
With $D=\operatorname{lcm}(448,256,78n,64n,8n)$, retain a mask
$I$ containing $4$ only if
\[
 c_I=\inf_{\mathscr B} v_-(b_I)-D^{-1}
       >\max\{U_3,\sup_{\mathscr B} G_{\rm full},\sup_{\mathscr B} G_I\}.
\]
Choose the largest retained $c_I$; write it as $c$.
For $A=t_1+t_2+t_3+3t_4$, clip the box to
\[
 \max(A_-,U_3+D^{-1})<A<\min(A_+,c-D^{-1}),
\]
and split this interval into finitely many equal rational layers.
Every layer satisfies \eqref{sv:restore-gates} with positive
margins and lies outside previously retained fourth prefixes.
On the clipped region,
\[
 2<\frac{c-S_4}{t_4}
       <2+\frac{9(\Lambda_{\rm cap}-\nu)}{\nu-{h_{\rm s}}}
       =2+\frac{84375}{90452}<3.
\]
Thus $\overline f(s)=\log(s-1)/s$ is increasing throughout.
If $A<a_j$ on a layer and $u_4$ is the fourth upper endpoint, use
$s_j=2+(c-a_j)/u_4$ in the lower bound
$4\overline f(s_j)\operatorname{Vol}(E_j)/(u_1u_2u_3u_4^2)$.
The exact volume of $\mathscr B\cap\{A<t\}$ is
\begin{equation}\label{sv:slab-volume}
 \frac1{72}\sum_{e\in\{0,1\}^4}(-1)^{|e|}
 \left(t-A_--\frac{e_1+e_2+e_3+3e_4}{n}\right)_+^4.
\end{equation}
This is the inclusion--exclusion formula for weighted box widths
$(1,1,1,3)/n$.  Differences of \eqref{sv:slab-volume} give
the layer volumes exactly.

\subsection{The lower coefficient and its original prime sequence}

For a prefix-closed tree with even first exits the counterpart of
\eqref{sv:upper-exit} is
\begin{equation}\label{sv:lower-exit}
 \sum_{d\in\mathcal T}\mu(d)\xi(d)
 =\prod_p(1-\xi(p))
 -\sum_{e\text{ first exit}}\xi(e)
               \prod_{p<p_{\rm last}(e)}(1-\xi(p)).
\end{equation}
Its proof is identical, using $\mu(e)=1$.
Taking $\xi(p)=\one_{p\mid n}$ gives a rough-indicator
minorant.  With $\xi=g_N$ and presieve $z_0$ it gives
\begin{equation}\label{sv:lower-projection}
 V_N(z_0)R_{\rm rough}-V_{\rm full}
 =\sum_{\substack{e\text{ full even exit}\\p_{\rm last}(e)<z_0}}
                      g_N(e)V_N(p_{\rm last}(e))\ge0.
\end{equation}
This is the density direction required for lower coefficient realization.

Round positive even coefficients inward and negative odd coefficients
outward, preserving the prefix decisions.  Attach beta-minus to a positive
rounded coefficient and beta-plus to a negative one.  Term by term
$a_{\rm rnd}(d)B_{\rm sign}\le a_{\rm rnd}(d)I_0$.
Consequently the physical coefficient satisfies
\begin{equation}\label{sv:physical-lower}
 \sum_{q\mid n}\rho_N^{\rm low}(q)\le
       \one_{(n,P_N(N^w))=1},\qquad
 \sum_q\frac{\rho_N^{\rm low}(q)}{\varphi(q)}
       \ge V_{\rm full}-\text{absolute density error}.
\end{equation}
The second inequality uses \eqref{sv:lower-projection} and the
absolute beta and box estimates.  It does not follow merely from
coefficientwise order, and beta-minus need not be nonnegative.
The sharp-prime consequence in Appendix~\ref{app:cross-review}, applied
to this same $\rho_N^{\rm low}$, turns \eqref{sv:physical-lower} into the
ordinary lower estimate for $\{N-p:p<N,\ p\text{ prime}\}$.

For an admissible lower prefix, the third-, fourth-, and fifth-stage
density changes are, respectively, upper parent minus upper child,
lower child minus lower parent, and upper parent minus upper child.
Evaluate each by~\eqref{sv:literal-density}.
At ranks $r=3,4,5$, the difference of the ordinary upper and lower
densities cancels the apparent inverse-square singularity in the last
coordinate, leaving a tail
$O(\epsilon(1+\log^{r-1}({h_{\rm s}}/\epsilon)))$.
The remaining compact regions admit one common finite refinement
that preserves every inherited decision.

\subsection{Finite feedback on positive marked families}\label{sv:finite-feedback}

The normalized states are estimates averaged over a nonnegative finite
marked family $\sigma(B)$ with a common total level $Q=N^\lambda$.
Put
\[
 L=\lambda-\log_NB>0,\qquad
 \Theta_Q=4\li(N)\sum_B
       \frac{\sigma(B)C_{BN}}{\varphi(B)\log(Q/B)},\qquad
 \alpha=\frac{1-\log_NB}{L}\ge2.
\]
For a singleton child of exponent $Lu$, or a pair of exponents $Lu,Lv$,
the residual level and parameter are, respectively,
\begin{equation}\label{sv:child-state}
 L'=L(1-u),\quad
 \alpha'=1+\frac{\alpha-1}{1-u};\qquad
 L'=L(1-u-v),\quad
 \alpha'=1+\frac{\alpha-1}{1-u-v}.
\end{equation}
The total $Q$, selected slots and prime ceiling are retained.
In particular $\alpha$ increases; the child is not reset to $\alpha=2$.

An admitted marked family here means a finite nonnegative combination
of labelled prime-box indicators, with fixed positive lower bounds on
all marked prime exponents and on $L$. In the classical case its
complete modulus level is below $N^{1/2}$ by a fixed power. In the
enhanced case $1/2<\lambda<153/256$ with fixed margins; it retains all full and selected-complement square
conditions and one factorization branch from
Appendix~\ref{app:coefficients}. It also retains a number $\chi>0$
such that adjoining primes of exponent below $\chi$ preserves that
branch and its support conditions. The number $\chi$ is no larger than the least
specified exponent. If a branch uses a bound $c$ for a third slot
not yet specified, it also satisfies $\chi\le c$.
Every enhanced ordinary node at argument $s$ has physical cutoff
$N^{L/s}$ below $N^\chi$. All these restrictions are imposed with
fixed margins on each box before endpoint limits are taken. The
estimates concern the marked average, rather than an individual fixed
prime prefix.

The child-to-parent measures are
$du/[u(1-u)]$ and $du\,dv/[uv(1-u-v)]$.
They are positive.  To specify the one-step switched cost, define
\begin{equation}\label{sv:one-step-mass}
 I(s,s_0,\alpha)=
 \int_{1/s_0<u_1<u_2<u_3<1/s}
 \omega\!\left(\frac{\alpha-u_1-u_2-u_3}{u_2}\right)
                    \frac{du_1\,du_2\,du_3}{u_1u_2^2u_3}.
\end{equation}
The one-step inequality gives the candidate
\begin{align}\label{sv:one-step-operator}
 \mathcal W_1[H,h](s;s_0,\alpha)
 ={}&A_+(s)-A_+(s_0)+H(s_0)-\frac K8 I(s,s_0,\alpha)\notag\\
 &+\frac12\int_{1/s_0}^{1/s}
 \frac{A_-(s_0(1-u))+h(s_0(1-u))}{u(1-u)}\,du,
 \qquad 2\le s<s_0\le5.
\end{align}
Here $A_+(s)=sF(s)/(2e^\gamma)$ and
$A_-(s)=sf(s)/(2e^\gamma)$, with the zero extension below
$2$.  The physical singleton child cutoff is $N^{L/s_0}$ and its
argument is $s_0(1-u)$, not the argument at a different total level.
The positive switched family is
$Bq_1q_2q_3n\le N$ with
$N^{L/s_0}\le q_1<q_2<q_3<N^{L/s}$ and $P^-(n)\ge q_2$.
Its unsifted mass, in the parent normalization, is $I/4$; the
factor $1/2$ in the pointwise weighted inequality gives the cost
$KI/8$ in \eqref{sv:one-step-operator}.  All its prime factors
have fixed positive exponents.  Its unit quotient has a fixed power
saving because $\alpha-3/s\ge1/2$.

A uniform bound usable in \eqref{sv:one-step-mass} is obtained by
putting $c_\omega=70893/125000$ and
\begin{equation}\label{sv:omega-envelope}
 \omega_{\rm env}(u)=
 \begin{cases}
 1,&u<1,\\
 \max\{1/u,c_\omega\},&u\ge1.
 \end{cases}
\end{equation}
For $v\ge u\ge1$ one has $\omega(v)\le\omega_{\rm env}(u)$.
On $[1,2]$ this follows from $\omega(v)=1/v$; on $[2,3]$
maximize $1+\log(v-1)-c_\omega v$.  The maximum is
$-\log c_\omega-c_\omega<0$, since the exact rational inequality
$c_\omega\sum_{j=0}^8c_\omega^j/j!>1$ implies
$c_\omega e^{c_\omega}>1$.  The delay equation propagates the
bound for $v\ge3$.  Thus replacing the integrand by
$\omega_{\rm env}((2-u_1-u_2-u_3)/u_2)$ bounds the cost for every
$\alpha\ge2$.  The branch below $1$ is a conservative bound for
the switched estimates that use it; unit quotient terms are counted
separately when their domains have nonnegligible mass.

For a two-step action write $a_0=1/s_0$, $a_i=1/\phi_i$ and $b=1/s$.
The parameter conditions are
\begin{gather}\label{sv:two-step-gates}
 2<s<\phi_3<\phi_2<\phi_1<s_0\le5,\qquad
 a_3+2a_1<1,\quad3a_1<1,\quad a_2+2a_1<1,\notag\\
 a_0<(1-a_2)/2<1/3<a_2<1/2,\qquad
 s_0(1-b)>2,\quad s_0-1\le4,\quad s_0(2-2a_2)>2.
\end{gather}
Its three singleton regions are
\begin{equation}\label{sv:singleton-domains}
 (a_0,b),\qquad(a_0,a_2),\qquad(a_0,a_3),
 \quad \rho(u)=s_0(1-u),
 \quad\text{physical cutoff }N^{La_0}.
\end{equation}
The regions on which the ordinary pair estimate and its feedback are used are
\[
\begin{array}{c|c|c}
 &\text{region}&\text{child argument}\\ \hline
 D_5^+&a_0<u<v<a_2,\ v+2u<1&s_0(1-u-v)\\
 D_6&(a_0,a_1)\times(a_2,a_3)&s_0(1-u-v)\\
 D_7&a_0<u<v<a_1&(1-u-v)/u\\
 D_8&(a_0,a_1)\times(a_1,a_2)&(1-u-v)/u,
\end{array}
\]
For the pointwise weighted inequality the fifth term has the complete
mother region
\begin{equation}\label{sv:complete-fifth-region}
 D_5=\{a_0<u<v<a_2\}
     =D_5^+\sqcup D_5^-,\qquad
 D_5^-=D_5\cap\{v+2u\ge1\}.
\end{equation}
Thus $\Gamma_5=\Gamma_5^++\Gamma_5^-$; it is not restricted to
$D_5^+$ in the combinatorial inequality. The ordinary estimate and
child defect are used only on $D_5^+$. On $D_5^-$ the complete count
is retained and bounded by the positive product estimate.
The test $2v<1$ holds throughout $D_5$, and the other three pair
regions satisfy both $2v<1$ and $v+2u<1$ by
\eqref{sv:two-step-gates}.
If $J_2^0$ denotes the ordinary and positive switched base for the
chosen admissible parameter tuple, the feedback added to that base is
\begin{equation}\label{sv:two-step-feedback}
 \frac15\left\{4H(s_0)+H(\phi_1)
 +\sum_{c\in\{b,a_2,a_3\}}\int_{a_0}^{c}
       \frac{h(s_0(1-u))}{u(1-u)}\,du
 +\sum_{i=5^+,6,7,8}\iint_{D_i}
       \frac{H(\rho_i(u,v))}{uv(1-u-v)}\,du\,dv\right\}.
\end{equation}
We now define this base and prove the sign of every contribution.

\paragraph{The twenty-one-term inequality.}
For a marked state let $R=Q/B$ and, in this paragraph, put
$\mathcal A_{N,B}=\{(N-p)/B:p<N,\ p\text{ prime},\ B\mid N-p\}$.
Let $\Phi_s$ be its positively marked average sifted at $R^{1/s}$.
Write $[l,r]^k_\uparrow$ for $k$ increasingly ordered exponents in
$[l,r]$.  For a positive integer $n$, let
\[
 Q_d(y;n)=\one_{d\mid n}\,
                  \one_{P^-(n/d)\ge y}.
\]
All additional primes in the following expressions are distinct;
collisions with the marked prefix, and square factors, are removed
with their fixed power-saving errors.  Every expression is then
averaged with the same nonnegative marked weights.
The first four terms are
\[
 \Gamma_1=4\Phi_{s_0}+\Phi_{\phi_1},\qquad
 \Gamma_2,\Gamma_3,\Gamma_4
   =\sum_{R^{a_0}\le q<R^c}Q_q(R^{a_0};n),
 \quad c=b,a_2,a_3,
\]
where the last expressions include the average over $n\in\mathcal A_{N,B}$.
For $\Gamma_5$ use the full region $D_5$ in
\eqref{sv:complete-fifth-region}; for $\Gamma_6,\Gamma_7,\Gamma_8$
use the other three pair regions above. The summand is
$Q_{q_1q_2}(R^{a_0};n)$ for $i=5,6$, and
$Q_{q_1q_2}(q_1;n)$ for $i=7,8$, with
$q_1=R^u<q_2=R^v$.
For $9\le i\le21$ the ordered exponent domains are
\begin{equation}\label{sv:switch-domains}
\begin{array}{c|l@{\qquad}c|l}
9 &[a_1,a_3]^3_\uparrow
 &10 &[a_1,a_2]^2_\uparrow[a_2,b]\\
11 &[a_1,a_2][a_2,a_3]^2_\uparrow
 &12 &[a_0,a_1]^2_\uparrow[a_3,b]\\
13 &[a_0,a_1][a_1,a_2][a_2,b]
 &14 &[a_0,a_1][a_2,b]^2_\uparrow\\
15 &[a_1,a_2][a_2,a_3][a_3,b]
 &16 &[a_2,a_3]^4_\uparrow\\
17 &[a_2,a_3]^3_\uparrow[a_3,b]
 &18 &[a_2,a_3]^2_\uparrow[a_3,b]^2_\uparrow\\
19 &[a_1,a_2][a_3,b]^3_\uparrow
 &20 &[a_2,a_3][a_3,b]^4_\uparrow\\
21 &[a_3,b]^6_\uparrow.&&
\end{array}
\end{equation}
For a domain with $r$ primes, the summand is
$Q_{q_1\cdots q_r}(q_{r-1};n)$.  Juxtaposition in
\eqref{sv:switch-domains} denotes independent successive bands,
with ordering only within each indicated ordered block.
The inequality to be used is
\begin{equation}\label{sv:twenty-one-weight}
 5\Phi_s\le\Gamma_1-\Gamma_2-\Gamma_3-\Gamma_4
                       +\sum_{i=5}^{21}\Gamma_i+o(\Theta_Q).
\end{equation}
It is the second weighted inequality of \cite[Lemma~4.2]{Wu2004};
the following verification fixes its complete ranges and multiplicities.

For a squarefree $R^{a_0}$-rough quotient let $A,B',C,D$ count
its prime factors in the four bands
$[a_0,a_1),[a_1,a_2),[a_2,a_3),[a_3,b)$, respectively.
(The symbol $B'$ in this verification is a band count, not the
marked prefix $B$.)  Pointwise the first four terms are
\[
 \Gamma_1=4+\one_{A=0},\qquad
 (\Gamma_2,\Gamma_3,\Gamma_4)
       =(A+B'+C+D,A+B',A+B'+C).
\]
The remaining terms have the following values, where
$x_+=\max(x,0)$:
\begin{equation}\label{sv:band-counts}
\begin{array}{c|l@{\quad}c|l}
5&\binom{A+B'}2&6&AC\\
7&(A-1)_+&8&\one_{A\ge1}B'\\
9&\one_{A=0}(B'+C-2)_+
 &10&\one_{A=0,B'\ge2}(C+D)\\
11&\one_{A=0,B'=1}(C-1)_+
 &12&\one_{A\ge2}D\\
13&\one_{A=1,B'\ge1}(C+D)
 &14&\one_{A=1,B'=0}(C+D-1)_+\\
15&\one_{A=0,B'=1,C\ge1}D
 &16&\one_{A=B'=0}(C-3)_+\\
17&\one_{A=B'=0,C\ge3}D
 &18&\one_{A=B'=0,C=2}(D-1)_+\\
19&\one_{A=0,B'=1,C=0}(D-2)_+
 &20&\one_{A=B'=0,C=1}(D-3)_+\\
21&\one_{A=B'=C=0}(D-5)_+.&&
\end{array}
\end{equation}
Indeed a quotient cutoff at the next-to-last chosen prime forces
all preceding chosen primes to be the smallest factors of $n$.
This gives each entry of \eqref{sv:band-counts}.  The signed
right-hand side in \eqref{sv:twenty-one-weight} is, respectively,
\[
\begin{array}{c|l}
A\ge2&(A+B'-2)(A+B'-3)/2+(A-2)C\\
A=1,\ B'\ge1&(B'-1)(B'-2)/2\\
A=1,\ B'=0&(1-C-D)_+\\
A=0,\ B'\ge2&(B'-2)(B'-3)/2\\
A=0,\ B'=1,\ C\ge1&0\\
A=0,\ B'=1,\ C=0&(2-D)_+\\
A=B'=0,\ C\ge3&0\\
A=B'=0,\ C=2,1,0&(1-D)_+,\ (3-D)_+,\ (5-D)_+.
\end{array}
\]
All are nonnegative, and at $A=B'=C=D=0$ the value is five.
The left side is five times the indicator of that last case.
Outside the $R^{a_0}$-rough support every summand vanishes.
This proves the pointwise inequality.  On a fixed positive exponent
box the omitted square and prefix-intersection terms have
$N^{1-\epsilon+o(1)}$ total mass for some fixed $\epsilon>0$,
which gives the stated error after the finite marked averaging.

\paragraph{Termwise estimates and the ordinary base.}
The two components of $\Gamma_1$ have upper multipliers
$4A_+(s_0)+A_+(\phi_1)$.  The three negative terms have lower
multipliers
\begin{equation}\label{sv:negative-singleton-base}
 \mathcal B(c)=\int_{a_0}^{c}
   \frac{A_-(s_0(1-u))}{u(1-u)}\,du,
 \qquad c=b,a_2,a_3.
\end{equation}
The support conditions put $s_0(1-u)$ in $(2,4]$ on these intervals,
so $A_-(s_0(1-u))=\log(s_0(1-u)-1)$.
The merged singleton and pair coefficients satisfy the complete
square and selected-complement tests of the parent state.  In
particular $D_5^+$, $D_6$, $D_7$, $D_8$ have the ordinary upper
kernel $(1-u-v)^{-1}$ against $du\,dv/(uv)$.
For a child argument below $2$, comparison with a smaller sifting
cutoff gives this same kernel, since $A_+=1$ up to $3$.
On
$D_5^-=\{a_0<u<v<a_2:v+2u\ge1\}$ the ordinary pair is not used.
The positive product estimate gives the valid kernel
$Kc_\omega s_0/4$, because its rough argument is at least
$s_0(2-2a_2)>2$.  Define
\begin{align}\label{sv:pair-base}
 \mathcal U_5&=
  \iint_{D_5^+}\frac{du\,dv}{uv(1-u-v)}
   +\frac{Kc_\omega s_0}{4}\iint_{D_5^-}\frac{du\,dv}{uv},\notag\\
 \mathcal U_i&=\iint_{D_i}\frac{du\,dv}{uv(1-u-v)}
                  \qquad(i=6,7,8).
\end{align}
Thus the unsupported corner is replaced by a positive upper
estimate; its ordinary mass or feedback is not retained as well.

For a switched domain $\mathcal D_i$ with $r$ primes, let
$w=u_{r-1}$ and $U=u_1+\cdots+u_r$.  Its normalized unsifted
mass at the actual parent ratio $\alpha$ is
\begin{equation}\label{sv:full-switched-mass}
 \mathcal M_i(\alpha)=
 \int_{\substack{\boldsymbol u\in\mathcal D_i\\U\le\alpha-w}}
 \frac{\omega((\alpha-U)/w)}{w\prod_j u_j}\,d\boldsymbol u
 +\int_{\substack{\boldsymbol u\in\mathcal D_i\\U=\alpha}}
                     \frac{du_1\cdots du_{r-1}}{\prod_j u_j}.
\end{equation}
In the surface integral one coordinate is eliminated with derivative
one.  The surface counts quotient $1$; it is separate from the
ordinary Buchstab term.  Prime summation and the rough-number
formula on fixed positive boxes give \eqref{sv:full-switched-mass}.
The upper count in parent units is $K\mathcal M_i(\alpha)/4$.

The following explicit bounds are uniform in $\alpha\ge2$.
Put $\ell_{ij}=\log(\phi_i/\phi_j)$,
$\ell_{01}=\log(s_0/\phi_1)$ and $\ell_{3b}=\log(\phi_3/s)$.
For $i=9,\ldots,15$ take $C_i=\max\{c_\omega,1/\max(1,m_i)\}$,
where
\[
\begin{array}{c|l@{\quad}c|l}
i&m_i&i&m_i\\ \hline
9&2\phi_3-3&10&2\phi_2-2-\phi_2/s\\
11&2\phi_3-\phi_3/\phi_2-2
 &12&2\phi_1-2-\phi_1/s\\
13&2\phi_2-\phi_2/\phi_1-1-\phi_2/s
 &14&2s-s/\phi_1-2\\
15&2\phi_3-\phi_3/\phi_2-1-\phi_3/s.&&
\end{array}
\]
These are lower bounds for the rough arguments on the respective
domains.  Equation~\eqref{sv:omega-envelope} gives
\begin{align}\label{sv:switch-mass-bounds}
 \widehat M_9&=C_9\{(\phi_1+\phi_3)\ell_{13}-2(\phi_1-\phi_3)\},\notag\\
 \widehat M_{10}&=C_{10}\{\phi_1-\phi_2(1+\ell_{12})\}
                                          \log(\phi_2/s),\notag\\
 \widehat M_{11}&=C_{11}\ell_{12}\{\phi_2(\ell_{23}-1)+\phi_3\},\notag\\
 \widehat M_{12}&=C_{12}\{s_0-\phi_1(1+\ell_{01})\}\ell_{3b},\notag\\
 \widehat M_{13}&=C_{13}\ell_{01}(\phi_1-\phi_2)\log(\phi_2/s),\notag\\
 \widehat M_{14}&=C_{14}\ell_{01}
                       \{\phi_2(\log(\phi_2/s)-1)+s\},\notag\\
 \widehat M_{15}&=C_{15}\ell_{12}(\phi_2-\phi_3)\ell_{3b},\notag\\
 \widehat M_{16}&=\phi_2\ell_{23}^4/24,
 &\widehat M_{17}&=\phi_2\ell_{23}^3\ell_{3b}/6,\notag\\
 \widehat M_{18}&=\phi_3\ell_{23}^2\ell_{3b}^2/4,
 &\widehat M_{19}&=\phi_3\ell_{12}\ell_{3b}^3/6,\notag\\
 \widehat M_{20}&=\phi_3\ell_{23}\ell_{3b}^4/24
                         +\phi_2\ell_{3b}^4/24,\notag\\
 \widehat M_{21}&=\phi_3\ell_{3b}^6/720+\phi_3\ell_{3b}^5/120.
 \end{align}
For $i\le19$ there is no unit surface because the sum of specified
exponents is at most $4/s<2\le\alpha$.
For $i=20$, eliminate the first variable on the surface and use
$1/u_1\le\phi_2$; the other four variables are ordered in
$[a_3,b]$.  This gives the second summand of $\widehat M_{20}$.
For $i=21$, eliminate the last variable, use $1/u_6\le\phi_3$,
and integrate the other five ordered variables.  The other
summands in \eqref{sv:switch-mass-bounds} are obtained by
integrating $1/(w\prod u_j)$ over their displayed bands, with
$\omega\le C_i$ for $9\le i\le15$ and $\omega\le1$
otherwise.  Hence $\mathcal M_i(\alpha)\le\widehat M_i$.

The ordinary base is now the explicit functional
\begin{equation}\label{sv:two-step-base}
 J_2^0=\frac15\left\{
 5A_+(s)-4A_+(s_0)-A_+(\phi_1)
 +\mathcal B(b)+\mathcal B(a_2)+\mathcal B(a_3)
 -\sum_{i=5}^8\mathcal U_i
 -\frac K4\sum_{i=9}^{21}\widehat M_i\right\}.
\end{equation}
Insert the upper estimates for positive terms and lower estimates
for negative terms in \eqref{sv:twenty-one-weight}.
The defects in the two components of $\Gamma_1$ subtract
$4H(s_0)+H(\phi_1)$ from its upper estimate.  The three singleton
lower defects subtract their integrals from the negative terms.
The four good-pair upper defects subtract the pair integrals from
their positive upper estimates.  After division by $5$ these are
exactly \eqref{sv:two-step-feedback}.  Therefore
\begin{equation}\label{sv:two-step-admissibility}
 \Phi_s\le\{A_+(s)-J_2^0-\mathcal P_2[H,h]+o(1)\}\Theta_Q,
\end{equation}
where $\mathcal P_2[H,h]$ is the expression in
\eqref{sv:two-step-feedback}.  This proves admissibility of a
finite two-step action from its mathematical base and positively
weighted preceding-generation defects.

Sharper upper bounds may replace the positive costs
in~\eqref{sv:pair-base} and~\eqref{sv:full-switched-mass}.
Such a bound must cover the same labelled family for every admissible
$\alpha\ge2$, including its unit surfaces. Replace that cost once
in~\eqref{sv:two-step-base}, and bound the resulting explicit
functional in the required direction to obtain the corresponding transition
constant. This preserves the sign argument for the two-step
inequality.

\subsection{A finite seed on the admitted marked families}

Put
\begin{equation}\label{sv:seed-parameters}
 S=\frac{11}{5},\qquad s_0=\frac{227}{50},\qquad
 (\phi_1,\phi_2,\phi_3)
 =\left(\frac{353}{100},\frac{29}{10},\frac{61}{25}\right),
 \qquad K_{\rm seed}=\frac{681}{100}.
\end{equation}
Here $K<K_{\rm seed}$, so Proposition~\ref{prop:current-positive-product}
also supplies the positive upper estimate with $K_{\rm seed}$.
Let $\mathcal J_{\rm seed}$ be the right-hand side of
\eqref{sv:two-step-base} at the parameters
\eqref{sv:seed-parameters}, with $s=S$ and $K$ replaced by
$K_{\rm seed}$. In particular its fifth cost is the full split cost
\eqref{sv:pair-base}, and the two unit-surface terms in
\eqref{sv:switch-mass-bounds} are included.

\begin{lemma}[Finite seed]\label{sv:finite-seed}
On an admitted marked family with $\alpha\ge2$, assume additionally
$L/S<\chi$ in the enhanced case. Then
\[
 \Phi_S\le\left(1-\frac{2327}{100000}+o(1)\right)\Theta_Q.
\]
The same upper defect is available at every $1\le w\le S$ whenever
the state admits the comparison cutoff $N^{L/S}$. The conclusion is
uniform over each fixed finite collection of such families, and follows
from four weighted steps after an initial seed of $1/50$.
\end{lemma}

\begin{proof}
All primes appended in the twenty-one-term inequality have exponent
at most $L/S$. They therefore lie below the retained ceiling, so the
largest specified slots, selected subsets and any third-slot bound
remain valid. The singleton residual square is $2u<1$. The pair squares
are $2v<1$ and $v+2u<1$; they hold on $D_5^+,D_6,D_7,D_8$.
Every new complete square is obtained by multiplying the corresponding
residual square by $B$. Once the selected old slots have occurred,
the new selected-complement square is the same full square. The
ordinary residual atoms are evaluated at their actual level
$Q/(Bq)$ or $Q/(Bq_1q_2)$. The remaining part $D_5^-$ is retained
with its positive cost, using no signed distribution claim there.
In the classical case the complete ordinary modulus is below $Q$,
and the absolute below-half remainder gives the same conservative
functional. The switched terms require only $\alpha\ge2$ and fixed
positive prime cutoffs.

The directed enclosure of the explicit functional is
\begin{equation}\label{sv:seed-base-certificate}
 \mathcal J_{\rm seed}\ge
 \frac{11672890052052913}{576460752303423488}
 >\frac{20249}{10^6}>\frac1{50}.
\end{equation}
The finite certificate for this inequality is specified in the
ancillary data under \texttt{seed\_certificate}. Its ordinary
integrals are those in \eqref{sv:negative-singleton-base} and
\eqref{sv:pair-base}; its switched masses are the explicit upper
bounds in \eqref{sv:switch-mass-bounds}. The ordinary double integral
is first integrated in one coordinate. The resulting one-dimensional
logarithmic integrals are enclosed by rational Darboux intervals and
the logarithm remainder of \eqref{num:log-enclosure}. Upper endpoints
are used for positive costs and lower endpoints for negative costs in
the definition of $\mathcal J_{\rm seed}$. Thus
\eqref{sv:seed-base-certificate} proves the initial defect $J_0=1/50$
before any later array recurrence is used.

Suppose an upper defect $J>0$ on $[1,S]$ has already been proved on
these admitted families. Positivity and cutoff comparison give
\[
 H_J(w)=J\one_{[1,S]}(w).
\]
The lower Buchstab identity at the same total level gives the
previously proved option
\begin{equation}\label{sv:seed-lower-child}
 h_J(w)=J\log\frac{S}{w-1}\,
                    \one_{[2,1+S]}(w).
\end{equation}
The normalization is independent of $w$, and $A_+(w)=1$ for
$1\le w\le S$. To check the comparison inside
\eqref{sv:seed-lower-child}, a Buchstab child indexed by
$v\in[w-1,S]$ has appended exponent $L/(v+1)$ and residual
exponent $Lv/(v+1)$. Its seed comparison exponent satisfies
\[
 \frac{Lv}{S(v+1)}\le\frac{L}{v+1}.
\]
It is therefore below that appended prime and the preceding stable
ceiling. Its new square holds because $v\ge1$. The lower child
need not itself admit argument $S$; the seed is applied to its
upper Buchstab children.

Insert $H_J,h_J$ into \eqref{sv:two-step-feedback}. Since
$s_0,\phi_1>S$, the unchanged-prefix defects vanish. Write
$u_*=67/227$, $b=1/S$, and set
\[
 E_i=\begin{cases}
 D_i\cap\{u+v\ge117/227\},&i=5^+,6,\\
 D_i\cap\{v\ge1-16u/5\},&i=7,8.
 \end{cases}
\]
The feedback is exactly $\lambda_{\rm seed} J$, where
\begin{align}\label{sv:seed-feedback-integral}
 \lambda_{\rm seed}=\frac15\bigg\{&
 \sum_{c\in\{b,a_2,a_3\}}
 \int_{u_*}^{c}
 \frac{\log\{S/[s_0(1-u)-1]\}}{u(1-u)}\,du\notag\\
 &+\sum_{i=5^+,6,7,8}\iint_{E_i}
                     \frac{du\,dv}{uv(1-u-v)}\bigg\}.
\end{align}
For $i=5^+,6$ the support condition is
$s_0(1-u-v)\le S$; for $i=7,8$ it is
$(1-u-v)/u\le S$. All retained pair-child arguments exceed $1$,
and the singleton lower arguments exceed $2$. There is no
feedback from $D_5^-$.

Every boundary in \eqref{sv:seed-feedback-integral} is rational
and affine. For a vertical slice $l(u)<v<r(u)$ the exact pair
integral is
\[
 \frac1{u(1-u)}
 \log\frac{r(u)(1-u-l(u))}{l(u)(1-u-r(u))}.
\]
For $E_{5^+}$ the successive $u$ intervals have endpoints
$50/227,117/454,19/58,1/3$. The domain $E_6$ is its entire
rectangle. The nonempty support in $E_7$ starts at $1265/5648$
and changes its lower boundary at $5/21$; the lower boundary of
$E_8$ changes at $1265/5648$. Rational Darboux enclosures of
these pieces give
\begin{equation}\label{sv:seed-feedback-certificate}
 \lambda_{\rm seed}\ge
 \frac{75057365825327405}{576460752303423488}
 >\frac{13}{100}.
\end{equation}
The exact endpoints, subdivisions and source correspondence are
specified by the same \texttt{seed\_certificate} data as
\eqref{sv:seed-base-certificate}.

Equations \eqref{sv:two-step-admissibility},
\eqref{sv:seed-base-certificate} and
\eqref{sv:seed-feedback-certificate} prove
\[
 J_{n+1}=\frac{20249}{10^6}+\frac{13}{100}J_n.
\]
At each step \eqref{sv:seed-lower-child} is derived from $J_n$
before the weighted inequality is applied. Starting at $J_0=1/50$
the four exact values are
\[
 \frac{22849}{10^6},\quad
 \frac{2321937}{10^8},\quad
 \frac{232675181}{10^{10}},\quad
 \frac{23273777353}{10^{12}}>\frac{2327}{100000}.
\]
There are finitely many appended-prime operations. Their exponents
remain positive after all compact choices have been fixed, so the
positive-product estimate and the distribution remainders apply at
each of them. The finitely many errors are $o(\Theta_Q)$; the strict
final surplus absorbs them in the stated order of limits.
\end{proof}

\paragraph{Use at the roots and descendants.}
For $U_3,U_4$, the specified prime has exponent $t>{a_{\rm s}}$, the residual
exponent is $L=v(t)-t$, and the upper argument is
$\max(2,L/{a_{\rm s}})$. On the enhanced range $t<57/224$, the largest-prime branch has
$2t<v(t)\le(1+t)/2$, so its stable ceiling can be taken below $t$
and arbitrarily close to it. Its physical cutoff exponent
$\min({a_{\rm s}},L/2)\le {a_{\rm s}}<t$ lies below that ceiling. Whenever the seed is used, its comparison exponent
$L/S$ is no larger than this physical cutoff. The classical tail
uses the below-half interface. For $U_{10},U_{11}$ the same argument
uses the least specified exponent $w$ and physical cutoff
$\min(w,L/2)$; the original selected subsets and first three slots
are already fixed. For the lower pair terms $L_5,L_6$, the seed
occurs only inside \eqref{sv:seed-lower-child}. The first appended
prime is at most $N^{a_{\rm s}}$, and every seed comparison prime is below
it by the preceding displayed inequality. These comparisons also
preserve a third-slot bound equal to ${a_{\rm s}}$.

More generally, every finite action in
\eqref{sv:one-step-operator} or \eqref{sv:two-step-feedback} retains
the total $Q$, every old square and the stable ceiling. New
singletons have $2u<1$; new pairs in the ordinary part have
$2v<1$ and $v+2u<1$. Their physical cutoffs are no larger than
the least newly appended prime. Equation~\eqref{sv:child-state}
shows that their ratios remain at least $2$. Induction therefore
keeps the seed and its Buchstab consequences within the admitted
families. No seed defect is inserted directly into an empty-prefix
lower start.

The finite induction uses the preceding generation on every
right-hand side. At each state take the maximum of alternative
estimates; when a child argument crosses a knot, use the minimum
over all adjoining closed cells. Each step therefore gives a bound
for the same positive marked counts. Only finitely many steps
are used.

Here is an explicit common approximation bound.  The positive operator
masses satisfy
\begin{align*}
 \|\mathcal B\|&\le\log5<2,\\
 \|\mathcal W_1\|&\le1+\tfrac12\log4<2,\\
 \|\mathcal W_2\|&\le1+\tfrac35\log4
                 +\tfrac45\,125(3/10)^2<12.
\end{align*}
For the last estimate $1/5\le u,v\le1/2$ and
$1-u-v\ge u\ge1/5$, so the pair kernel is at most $125$ on
each of four regions.  The unchanged-prefix coefficients sum to one.
The switched positive costs belong to the base of the action and
are not additional feedback edges.
If each local oriented approximation error is at most $\delta$,
then
\begin{equation}\label{sv:feedback-error}
 e_{j+1}\le\delta+12e_j,\qquad
 e_j\le\delta\frac{12^{j+1}-1}{11}.
\end{equation}
The absolute coefficients of the eleven-term functional sum to $14$
before its factor $1/4$, giving root norm $7/2$.

At depth at most $440$, every residual level is at least
$\lambda_*101^{-440}>0$, and an appended prime exponent is at
least this number divided by $6$.  Recursively refine appended
exponents on each fixed positive parent, using an error within
$\delta$ times its normalized mass.  Proper cap and knot walls are
null; structural clamps and inherit identities are retained.
The denominators are bounded away from zero.  Equation
\eqref{sv:feedback-error} then gives one finite set of queries with
arbitrarily small total error.  Only after fixing that set does $N$
increase.  A positive-mass vertex has a positive limiting
$\Theta_Q/H_N$, so the absolute distribution errors become
$o(\Theta_Q)$; zero-mass pieces are handled by absolute errors.

\section{Endpoint estimates and the order of limits}\label{app:cross-review}

The sharp-prime lemma (Lemma~\ref{lem:current-sharp-prime}) applies to
the same coefficient constructed above.  We record the absolute bounds
that allow it to be used for signed lower weights.  If
$|\rho_N^{\rm low}(q)|\le C$, its squarefree support is $q\le N^\theta$,
$\theta<1$, and $\delta=(\log N)^{-B}$, then the endpoint contribution
is bounded by
\begin{equation}\label{sv:sharp-boundary}
 \sum_{\substack{p\text{ prime}\\0<|p-N|\le\delta N}}
          \sum_{q\mid N-p}|\rho_N^{\rm low}(q)|
 \le2C\sum_{1\le n\le\lceil\delta N\rceil}\tau(n)
 \ll\delta N\log N+\log N.
\end{equation}
The difference $p-N=0$ is impossible for even $N>2$.
The exact-coprime mean boundary has the same bound because
$\sum_{q\le N^\theta}1/\varphi(q)\ll\log N$.
The smooth lower trim below $2N^{1-\kappa_0}$ costs
$N^{1-\kappa_0+o(1)}$ in the progression sum and
$O(N^{1-\kappa_0}\log N)$ in the mean.  These have fixed power
savings.  Choosing $B>A+3$ makes all boundary terms smaller than the
prescribed logarithmic error.

For the exact-coprime prime mean, let $\omega_{\mathrm{pf}}(q)$ be
the number of distinct prime divisors of $q$.  Every such divisor is
below $N$, and squarefreeness gives
\begin{equation}\label{sv:coprime-prime-error}
 \sum_q\frac{|\rho_N^{\rm low}(q)|\omega_{\mathrm{pf}}(q)}{\varphi(q)}
 \ll\sum_{\ell\le N^\theta}\frac1{\ell-1}
                  \sum_{e\le N^\theta/\ell}\frac1{\varphi(e)}
 \ll\log^2N.
\end{equation}
The sum over $\ell$ is over primes; summing over all integers already
gives the stated crude bound.  Removing the $O(\log N)$ prime divisors
of $N$ from the original sequence contributes $N^{o(1)}$ by the
divisor bound.  These estimates justify using the sharp original
prime count in \eqref{sv:physical-lower}, without obtaining a prime
lower bound from the positive-product upper constant $K$.

The upper and lower count estimates follow the parameter order in
Section~\ref{sec:lower-limits}. The fixed moment orders precede the
choice of $\kappa$ in~\eqref{coef:parameter-order}, and the boundary
exponent satisfies $B>A+3$, as above.

\section{Finite numerical estimates}\label{app:numeric}

We specify the finite quantities that determine the two constants
in Lemma~\ref{lem:numerical-endpoints}. The construction separates
integral enclosures, the finite defect recurrence, and the final
rational subtraction, so that each contribution has an explicit
mathematical definition and direction.

\subsection{Supplementary data}
The finite data are supplied in \emph{Theorem (1+1.75) on the Goldbach
Conjecture --- Numerical Supplement}, version 20261106~\cite{ChanSupplement}.
They comprise $106$ affine transitions, $200$-cell defect arrays at
depths $248$, $312$, and $430$, and $17572$ paired rectangles.
Rational quantities in \filepath{finite-construction-data.json} are
represented as quotients of integers. The fields
\texttt{upper\_saving\_bounds}, \texttt{lower\_root\_bounds}, and
\texttt{lower\_integral\_bounds} give directed rational endpoints.
The map \filepath{certificate-map.json} identifies their source
certificates and the corresponding domains, subdivisions, and interval
witnesses. Table~\ref{num:certificate-table} gives the correspondence
between the mathematical estimates and their certificate fields.

\hypertarget{num:verification-scope}{}
\paragraph{Verification commands and inputs.}
Run the commands below from the supplement root with Python 3.11 or later,
with assertions enabled. The sixth-rank computations also use NumPy.
The rows distinguish reconstruction from saved certificates, exact checks of
the finite construction, and recalculation of specified witness bounds.
All output files requested with \texttt{-{}-output} or \texttt{-{}-output-dir}
are written outside the supplement.

\begingroup
\renewcommand{\LTcaptype}{}
\stepcounter{LT@tables}
\small
\setlength{\tabcolsep}{4pt}
\begin{longtable}{@{}>{\raggedright\arraybackslash}p{.40\textwidth}>{\raggedright\arraybackslash}p{.29\textwidth}>{\raggedright\arraybackslash}p{.25\textwidth}@{}}
\toprule Entry point & Quantities checked or recomputed & Numerical inputs retained \\
\midrule
\endhead
\bottomrule
\endfoot
\filepath{verify-package.py} & File set, sizes and recorded SHA-256 identities. & The versioned \filepath{MANIFEST.json}.\\[3pt]
\filepath{finite-data/reproduce-data.py} \texttt{-{}-check} & Exact extraction of finite data and field correspondence from source certificates. & The supplied source-certificate values.\\[3pt]
\filepath{finite-data/verify-finite-data.py} & Finite defect recurrence; paired-cell geometry, masses and loss proxy; $\Gamma_5$ decompositions; density sums and strict scalar margins. & Action coefficients and integral endpoints identified by \filepath{certificate-map.json}.\\[3pt]
\filepath{replay-records.py} & Ordered-cover keys, stored integer-floor sums and the cost identity. & Stored floors in the ordered witness.\\[3pt]
\filepath{replay-ordered-cover.py} & The selected-method floor on every row of the grid-48 base cover; use \texttt{-{}-grid 32} for its predecessor. & The recorded method choices and generation-248 defect table.\\[3pt]
\filepath{verify-sixth-domains.py} & Affine sieve-argument ranges on the saved polyhedral cells. & Recorded vertices, coordinate maps and caps.\\[3pt]
\filepath{replay-sixth-rank.py} & Grid-400 sixth-rank cube contributions and their sum, together with the preceding polyhedral range checks. & The supplied cube witness and polyhedral records.\\
\end{longtable}
\endgroup

The recurrence and paired-cell checks use the action and integral inputs
listed in Table~\ref{num:certificate-table}. Their primitive enclosure
obligations remain those of the corresponding integrals and policies in
this appendix and Appendices~\ref{app:upper}--\ref{app:lower}.
The \texttt{enclosure\_obligations} entries of \filepath{certificate-map.json}
list the active assertions, complete records and supplied algorithm sources
for each primitive family. Witness generation is described in the package
README; in particular \filepath{witnesses/generator-0041.py} generates the ordered
base cover and \filepath{witnesses/generator-0045.py} the sixth-rank cube witness.
The generation-248 profile and other numerical inputs retained by a replay
are checked through their own entries, so that each verification has an
explicit set of inputs. The package README gives the command options; the
fixed-point arithmetic for the finite data is specified below.

\subsection{Exact scalar endpoints}
The exact values used in the upper density estimate are
\begin{align}
 E_{\rm up}&=\frac{9979702576710911234850619807816847864907359959}
 {1475913550962201227946780619099353215139840000},
 \label{eq:current-E}\\
 u&=\frac{145007961227048312210593427530128000784681}
 {11807308407697609823574244952794825721118720000}.
 \label{eq:current-u}
\end{align}
For the lower construction, before the paired-fibre increment, put
\begin{align}
 R_{\rm base}&=
 \frac{1345575129530003739571681432301188857312754737}
 {685738127494661151198444592354557952000000000},
 \label{eq:base-R}\\
 B_{\rm base}&=
 \frac{2504748318091942808229694219379799307550987}
 {5485905019957289209587556738836463616000000000}.
 \label{eq:base-B}
\end{align}
The two copies of $u$ and of $B_{\rm base}$ pay for distinct finite
approximation and coefficient-realization errors. Quadrature errors
are already contained in the oriented integral endpoints. The ancillary
fields \texttt{scalar\_bounds} contain the resulting exact sums and
strict surpluses.

\subsection{Directed arithmetic}
All rounding used to prove an upper bound is outward, and all rounding
used to prove a lower bound is inward.  More explicitly, at denominator
$\mathsf Q=2^{b_{\rm fp}}$ a rational interval $[x,y]$ is represented by
$[\lfloor \mathsf Q x\rfloor,\lceil \mathsf Q y\rceil]$.  Addition and subtraction
are exact at the integer level.  For multiplication the four endpoint
products are enclosed by their minimum and maximum and rounded in the
indicated directions after division by $\mathsf Q$.  Division is performed
only after verifying that the denominator interval is positive.

For logarithms put $z=(x-1)/(x+1)$.  If $1\le x\le2$, then
\begin{equation}
 2\sum_{j=0}^{n-1}\frac{z^{2j+1}}{2j+1}
 \le\log x\le
 2\sum_{j=0}^{n-1}\frac{z^{2j+1}}{2j+1}
 +\frac{2z^{2n+1}}{(2n+1)(1-z^2)}.
 \label{num:log-enclosure}
\end{equation}
This follows by summing the positive power series for
$2\operatorname{arctanh}z$ and bounding the remaining denominators
below by $2n+1$.  For $n=96$ the remainder is at most
$9/(4\cdot193\cdot3^{193})$.  Range reduction and a separate
enclosure of $\log2$ extend the formula to every positive rational
argument used here.  The paired-cell calculation uses denominator
$2^{224}$ for the series and denominator $2^{160}$ for its returned
intervals.  Thus no floating-point evaluation of a logarithm determines
the direction of an inequality.

\subsection{Enclosing delayed functions and integrals}
It is convenient to propagate the numerators
\[
 A_+(s)=\frac{sF(s)}{2e^\gamma},\qquad
 A_-(s)=\frac{sf(s)}{2e^\gamma}.
\]
Their initial branches are $A_+(s)=1$ for $1\le s\le3$ and
$A_-(s)=0$ for $s\le2$, and subsequent branches satisfy
\[
 A_+'(s)=\frac{A_-(s-1)}{s-1},\qquad
 A_-'(s)=\frac{A_+(s-1)}{s-1}.
\]
The equation for a branch involves only the preceding unit interval.
One therefore obtains rational interval enclosures successively on a
rational mesh, splitting at every integer breakpoint.  Lower and upper
integral bounds follow from the corresponding interval bounds for the
integrand.  Chord bounds and curvature bounds sharpen these enclosures
without changing their direction.  An arbitrary query is enclosed on
its entire support, including both adjoining cells when an endpoint
is a mesh knot.  The analogous construction for $s\omega(s)$ uses
$(s\omega(s))'=\omega(s-1)$ and its initial branch
$\omega(s)=1/s$ on $[1,2]$.

The chord integration can be implemented without numerical
quadrature.  On $[j/n,(j+1)/n]$ the integral of the affine
endpoint interpolant of a numerator $V$, divided by its argument,
is
\[
 \bigl((j+1)\log(1+1/j)-1\bigr)V(j/n)
 +\bigl(1-j\log(1+1/j)\bigr)V((j+1)/n).
\]
Both weights are positive because they are integrals of
nonnegative affine basis functions against $du/u$.
The monotonicities of $F$ and $f$ imply that $A_+$ is convex
and $A_-$ is concave on their nonconstant branches.
Their numerator curvature on a delayed cell is bounded
in absolute value by $1/(j/n-1)$.  Thus the integrated
chord error is at most
\[
 \frac{1}{12nj(j-n)}\qquad(j>n).
\]
This quantity is subtracted from a chord lower estimate
or added to a chord upper estimate as appropriate.
The weights themselves have alternating rational series
\[
 (j+1)\log(1+1/j)-1
  =\sum_{k\ge1}\frac{(-1)^{k+1}}{k(k+1)j^k},\qquad
 1-j\log(1+1/j)
  =\sum_{k\ge1}\frac{(-1)^{k+1}}{(k+1)j^k}.
\]
Their fourth and fifth partial sums enclose the exact
weights.  This supplies rational enclosures even where
direct subtraction of nearly equal logarithms would be
numerically inconvenient.

The finite integrations use rational boxes or rational polytopes.
For a polytope given by affine inequalities, every possible active
basis is solved exactly, and each resulting vertex is tested against
every original inequality.  Boundedness and the face incidences are
verified before a pulling triangulation is formed.  Polynomial
moments are then exact: on a $d$-simplex $T$ with barycentric
coordinates $\lambda_0,\ldots,\lambda_d$,
\begin{equation}
 \int_T\prod_{j=0}^d\lambda_j^{\alpha_j}\,dx
 =d!\operatorname{vol}(T)
   \frac{\prod_{j=0}^d\alpha_j!}{(d+\sum_j\alpha_j)!}.
 \label{num:simplex-moment}
\end{equation}
Rational lower or upper bounds for the remaining kernel factors,
together with these moments and \eqref{num:log-enclosure}, give the
required integral enclosure.  A boundary box in a negative-cost upper
integral is retained with an upper bound for its whole contribution.
In a positive saving, a box may instead be omitted after establishing
that its omitted contribution is nonnegative.  These two rules have
different signs and are used separately.

\subsection{The upper density}
Write
\[
 U_0=\frac4\Theta+I_1+I_{3,c}+I_{3,t},\qquad
 E_{\rm up}=U_0-\sum_{j=1}^{22}G_j.
\]
Here $I_{3,c}$ is the integral over the region with last coordinate
at least $1/100$, and $I_{3,t}$ bounds its entire complementary
tail.  Put $D=(F-f)/(2e^\gamma)$.  Its delay equation gives
\[
 D'(s)=-\frac{D(s)+D(s-1)}s\le-\frac{2D(s)}s
 \quad(s\ge3),
\]
using positivity and monotonicity of $D$.  Therefore
$D(s)\le D(10)(10/s)^2$ for $s\ge10$.  If
$\delta=1/100$ and $\rho_0=19/96$, the upper-tree cap
inequalities give a residual exponent at least $\rho_0$.
For $0<z<\delta$ its delayed argument is at least
$(\rho_0-z)/z\ge10$.  Enlarging the two remaining ordered
coordinates to $z\le y\le x\le\rho_0$ yields the complete
tail bound
\begin{equation}
 I_{3,t}\le
 \frac{200D(10)\delta}{(\rho_0-\delta)^2}
 \left\{\log^2\frac{\rho_0}{\delta}
        +2\log\frac{\rho_0}{\delta}+2\right\}.
 \label{num:complete-tail}
\end{equation}
Here the ordered $x,y$ integral is
$\tfrac12\log^2(\rho_0/z)$, and the remaining integral
over $0<z<\delta$ is elementary.  In particular the
logarithmic singularity at $z=0$ is integrated rather
than discarded.  The first integral is
enclosed by a positive series with its geometric remainder.
The compact third integral is covered by a disjoint partition,
including conservative bounds on its boundary boxes.

Table~\ref{num:upper-table} gives readable lower bounds for the
nonnegative savings $G_j$: its integer entries satisfy
$G_j\ge m_j10^{-14}$. The full rational bounds are in the
supplementary data. Successive changes on one region are measured
from their respective parents. In particular, the last two rows
give only the curvature and mixed increments over the preceding
bound for that integral.

\begin{longtable}{@{}r p{.62\textwidth} r@{}}
\caption{Rational lower bounds $10^{-14}m_j$ for the upper-density savings.}
\label{num:upper-table}\\
\toprule
$j$ & Saving & $m_j$\\
\midrule
\endfirsthead
\toprule $j$ & Saving & $m_j$\\\midrule
\endhead
1 & Fourth-rank cap replacement & 3507783546602\\
2 & Fifth-rank parent subdivision & 1005163763573\\
3 & Disjoint fifth-rank region & 136378695993\\
4 & Sixth-rank cap replacement & 96281257455\\
5 & Second-rank fallback above $\mu_0$ & 4209783419\\
6 & Third-rank fallback above $\mu_0$ & 3908979560\\
7 & Sixth-rank logarithmic moments & 1525398138\\
8 & Sixth-rank descending region & 187855687\\
9 & Sixth-rank ascending cap difference & 387157212\\
10 & Sixth-rank descending cap difference & 236377744\\
11 & Sixth-rank clipped lower window & 846201484\\
12 & Sixth-rank descending lower window & 500857555\\
13 & Sixth-rank local cap corrections & 50721206\\
14 & Sixth-rank nonsingleton regions & 11743733\\
15 & Eighth-rank radial bound & 94105689\\
16 & Eighth-rank ascending regions & 1543640141\\
17 & Eighth-rank descending regions & 590010094\\
18 & Eighth-rank unequal-parent moments & 1304343681\\
19 & Eighth-rank third-coordinate moments & 357686854\\
20 & Eighth-rank prefix-dependent moments & 37235787\\
21 & Eighth-rank curvature increment & 97705\\
22 & Eighth-rank mixed increment & 607514\\
\bottomrule
\end{longtable}

The upper enclosures use $3699127$ fourth-rank boxes,
$292633$ fifth-rank parents subdivided into $18728512$ children,
$1763420$ boxes in the disjoint fifth-rank region, and $632104$
sixth-rank cubes. The eighth-rank regions have disjoint families
of sizes $30$, $15$, $13$, and $29$. Across ranks the cap
differences telescope. On a parent region take the maximum of its
valid saving bounds, comparing its own bound with the complete
child sum.

An upper bound for the four base terms, in the same denominator
$10^{14}$, is respectively
\[
 668896321070235,\quad76020940667,\quad
 11960259849166,\quad12313775.
\]
Let $E_{\rm up,14}$ be their sum minus the sum of Table~\ref{num:upper-table},
divided by $10^{14}$.  Exact integer arithmetic gives
\[
 E_{\rm up}\le E_{\rm up,14},\qquad E_{\rm up,14}+2u<K,
\]
where $u$ is given in \eqref{eq:current-u}; indeed the last difference
is greater than $9.65\cdot10^{-8}$.  The unrounded rational sum
gives \eqref{eq:current-E}.  The two copies of $u$ pay separately
for finite approximation and coefficient realization.  Quadrature errors are included in the one-sided integral endpoints.

\subsection{The finite defect recurrence}
Put $\mathsf Q=2^{80}$ and
\[
 C_i=[1+i/50,1+(i+1)/50],\qquad 0\le i<200.
\]
All table entries in this subsection are integer numerators at
denominator $\mathsf Q$.  A four-step seed is obtained by starting at
$J_0=1/50$ and using
\[
 J_{n+1}=\frac{20249}{10^6}+\frac{13}{100}J_n.
\]
The initial defect and the implication represented by this recurrence
are supplied by Lemma~\ref{sv:finite-seed}, with its marked-state and
physical-cutoff conditions. The ancillary \texttt{seed\_certificate}
contains the separate ordinary and feedback lower endpoints used
there. Thus the arithmetic starts from an analytic saving valid on
the required family. Then $J_4=0.023273777353>2327/100000$. Initialize
\begin{equation}
 H_i^{(0)}=
 \begin{cases}
 \lfloor \mathsf Q\,2327/100000\rfloor,&1+(i+1)/50\le11/5,\\
 0,&\text{otherwise},
 \end{cases}
 \qquad h_i^{(0)}=0.
 \label{num:table-seed}
\end{equation}

For completeness the exact integer transition is specified here.
An action $\alpha$ consists of an integer $b_\alpha$, a divisor
$d_\alpha\in\{1,5\}$, point pairs $(j,m)$, nonnegative integer
weights $w_{\alpha j},v_{\alpha j}$, and a target set $T_\alpha$.
Its candidate value is
\begin{equation}
 b_\alpha+\left\lfloor\frac1{d_\alpha}
 \left\{\sum_{(j,m)}mH_j+
 \sum_j\left\lfloor\frac{w_{\alpha j}h_j}{\mathsf Q}\right\rfloor+
 \sum_j\left\lfloor\frac{v_{\alpha j}H_j}{\mathsf Q}\right\rfloor
 \right\}\right\rfloor.
 \label{num:action-map}
\end{equation}
The $70$ one-step actions have $d_\alpha=1$ and the $36$ two-step
actions have $d_\alpha=5$.  They are evaluated at the upper
constant $K_*=6.761767>K$.  The action parameters, weights,
constants and target sets are given individually in the
\texttt{actions} array.  The target condition for action parameter
$s_\alpha$ is $1+(i+1)/50\le s_\alpha$.

Let $\ell_j$ be the integer lower bound for
$\mathsf Q\log((1+(j+1)/50)/(1+j/50))$ in the data, and set
\[
 S_H(j)=\sum_{k=j}^{199}\left\lfloor\frac{H_k\ell_k}{\mathsf Q}\right\rfloor,
 \qquad
 S_h(j)=\sum_{k=j}^{199}\left\lfloor\frac{h_k\ell_k}{\mathsf Q}\right\rfloor.
\]
For $50\le i<200$ the Buchstab candidates are
$S_h(\max(50,i-49))$ for $H_i$ and $S_H(i-49)$ for $h_i$.
Each new entry is the maximum of its preceding value and all
applicable candidates.  Every right-hand side is evaluated using
only the preceding arrays.  This is a simultaneous finite
recurrence, with no assumption of convergence to a fixed point.

The analytic action inequalities and their closure under descendants
imply by induction that the rational arrays are lower bounds for
the defect functions on the whole closed cells.  Floors in
\eqref{num:action-map} preserve this direction because the
feedback weights are nonnegative.  Applying the map $430$ times
to \eqref{num:table-seed} gives the $400$ entries in
\texttt{defect\_tables.430}.  The arrays at depth $248$ are
bounded below by the corresponding finite iteration.  The
depth-$430$ arrays are used only in the two occurrences
$U_3,U_4$; the other designated occurrences use the depth-$248$
arrays.  Thus a change of table is charged only at its actual
occurrence in the lower functional.

\subsection{The integrals attached to the eleven counts}
\label{num:eleven-integrals}
The correspondence with the arithmetic terms $Y_j$ is fixed as follows.
Write $A_+(s)=sF(s)/(2e^\gamma)$,
$A_-(s)=sf(s)/(2e^\gamma)$, and put
$\nu=249/448$, $w_*=25/224$. The reference one-prefix terms are
\begin{align*}
 L_1&=\frac4{a_{\rm s}}\overline f(\nu/{a_{\rm s}}),\\
 L_2&=4\left\{(1/{h_{\rm s}}+1)\overline f\left(\frac{1+{h_{\rm s}}}{2{h_{\rm s}}}\right)
                 -\overline f(249/50)\right\}+H_{\rm rem},
\end{align*}
where $H_{\rm rem}$ is defined in~\eqref{num:removed-defect}.
No defect table is inserted at the unmarked root of $L_1$:
its ratio $1/\nu$ is less than $2$.

For the two-prefix lower terms put
\[
 \mathcal P_5=\{{a_{\rm s}}<u<t<{h_{\rm s}}\},\qquad
 \mathcal P_6=\{{a_{\rm s}}<u<{h_{\rm s}}<t<{d_{\rm s}}\}.
\]
On the enhanced part $t<\mu_0$, let
$V(t,u)$ be the maximum of the admitted candidates
$v(t),v(t+u),\Lambda_{\rm cap}-3{a_{\rm s}}/8$; on $t>\mu_0$ take $V=v(t)$.
Here admission retains the full and selected squares, the selected
largest slots and the physical prime ceiling, as in
Appendix~\ref{app:coefficients}. Set $L=V-t-u$. The lower integrals are
\begin{equation}
 L_j=4\iint_{\mathcal P_j}
        \frac{A_-(L/{a_{\rm s}})+h_{248,\rm adm}(t,u)}{tuL}\,du\,dt,
 \qquad j=5,6.
 \label{num:pair-lower-integrals}
\end{equation}
The defect term equals $h_{248}(L/{a_{\rm s}})$ only on a fully admitted state
with $2\le L/{a_{\rm s}}\le5$, and is zero otherwise. The ordinary lower term
is also zero when $L/{a_{\rm s}}<2$. Thus an uncertain defect lookup can be
discarded without changing an ordinary lower estimate.

For $m=248,312,430$, the one-prefix upper integrals are
\begin{equation}
 U_j^{(m)}=4\int_{I_j}
       \frac{A_+(s_j(t))-H_m(s_j(t))}{tL_j(t)}\,dt,
 \quad I_3=({a_{\rm s}},1/3),\quad I_4=({a_{\rm s}},{d_{\rm s}}),
 \label{num:one-prefix-upper-integrals}
\end{equation}
with $L_j,s_j$ given in~\eqref{num:table-change-integral} below.
The reference values in the eleven-term sum use $m=248$.
The two subsequent table improvements are inserted only through
\eqref{num:successive-table-gains}.

The last two costs use decreasing coordinates $x>y>z>w$ on
\[
 \mathcal D_{10}=\{{a_{\rm s}}<w<z<y<x<{h_{\rm s}}\},\qquad
 \mathcal D_{11}=\{{a_{\rm s}}<w<z<y<{h_{\rm s}}<x<{d_{\rm s}}\}.
\]
The integration measure is
\[
 {d_{\rm s}}\sigma_4=\frac{dx\,dy\,dz\,dw}{xyzw}.
\]
Put $S=x+y+z+w$. The switched option has kernel
$K_0\omega((1-S)/z)/z$. Each admitted ordinary branch with residual
exponent $L=\lambda-S$ supplies the kernel
$4\{1-H_{248}(\max(2,L/w))\}/L$; all its full and selected support conditions
must hold on the whole cell. The chosen cell majorant bounds the
minimum of these options from above. Its integral defines the upper
enclosure for $U_{10}+U_{11}$. The cell choices, residual intervals,
closed-cell defect minima and complete ordered cover are the data
of the four-prefix certificate. The switched option remains available
when an ordinary branch fails its support conditions. Since ${d_{\rm s}}+3{h_{\rm s}}<1$, these two
domains have no unit-quotient contribution.
The remaining switched costs $U_7,U_8,U_9$ are defined in
Section~\ref{num:cost-slope}.

\subsection{Definitions of the lower correction terms}
\label{num:lower-definitions}
We specify the correction terms before giving their numerical
combination. Put $\nu=249/448$, $w_*=25/224$ and write
$\overline F=F/(2e^\gamma)$, $\overline f=f/(2e^\gamma)$.
The difference between the ordinary start at ${h_{\rm s}}$ and the ordinary
part of the least-prime decomposition is
\begin{equation}
 B_{\rm ord}=\frac4{h_{\rm s}}\overline f(\nu/{h_{\rm s}})
 -4\left\{(1/{h_{\rm s}}+1)\overline f\left(\frac{1+{h_{\rm s}}}{2{h_{\rm s}}}\right)
                   -\overline f(249/50)\right\}.
 \label{num:ordinary-start-difference}
\end{equation}
The defect part removed from that decomposition is the whole integral
\begin{equation}
 H_{\rm rem}=4\int_{w_*}^{{h_{\rm s}}}
   \frac{H_{248}((1-t)/(2t))}{t(1-t)/2}\,dt.
 \label{num:removed-defect}
\end{equation}
The finite ordinary estimate at ${h_{\rm s}}$ is inserted only after this
entire defect term has been removed. On a defect cell $[s_l,s_r]$,
the contributing interval in~\eqref{num:removed-defect} is
$(w_*,{h_{\rm s}})\cap[1/(2s_r+1),1/(2s_l+1)]$; its primitive is
$2\log(t/(1-t))$, before multiplication by the factor $4$.

For $w={a_{\rm s}},{h_{\rm s}}$, let $C_3,C_4,C_5$ be the successive caps of the same
lower policy in Appendix~\ref{app:lower}. A cap is inherited or
selected from the finitely specified subproduct branches subject to
all full and selected-complement support conditions. In particular
$\nu\le C_3\le C_4\le C_5\le\Lambda_{\rm cap}$. For decreasing
$t_1,\ldots,t_r$, write $S_r=\sum_{j\le r}t_j$ and
\[
 {d_{\rm s}}\mu_r=\frac{dt_1\cdots dt_r}{t_1\cdots t_{r-1}t_r^2}.
\]
The three retained domains are
\begin{align*}
 \mathcal E_3(w)&=\{0<t_3<t_2<t_1<w:t_1+3t_2<\nu\},\\
 \mathcal E_4(w)&=\{(t_1,\ldots,t_4):(t_1,t_2,t_3)\in\mathcal E_3(w),
          \ 0<t_4<t_3,\ S_3+3t_4<C_3\},\\
 \mathcal E_5(w)&=\{(t_1,\ldots,t_5):(t_1,\ldots,t_4)\in\mathcal E_4(w),
                                      \ 0<t_5<t_4\}.
\end{align*}
Thus the three nonnegative density gains are
\begin{align}
 I_3(w)&=4\int_{\mathcal E_3(w)}
  \left\{\overline F\left(\frac{\nu-S_3}{t_3}\right)
        -\overline F\left(\frac{C_3-S_3}{t_3}\right)\right\}\,{d_{\rm s}}\mu_3,
                                      \label{num:third-lower-gain}\\
 I_4(w)&=4\int_{\mathcal E_4(w)}
  \left\{\overline f\left(\frac{C_4-S_4}{t_4}\right)
        -\overline f\left(\frac{C_3-S_4}{t_4}\right)\right\}\,{d_{\rm s}}\mu_4,
                                      \label{num:fourth-lower-gain}\\
 I_5(w)&=4\int_{\mathcal E_5(w)}
  \left\{\overline F\left(\frac{C_4-S_5}{t_5}\right)
        -\overline F\left(\frac{C_5-S_5}{t_5}\right)\right\}\,{d_{\rm s}}\mu_5.
                                      \label{num:fifth-lower-gain}
\end{align}
The fifth square follows from
$S_4+2t_5<S_3+3t_4<C_3\le C_4$.
The cap choices, branch labels and parent cubes are fixed jointly;
an integral evaluated with a different fourth cap cannot be appended
to~\eqref{num:fifth-lower-gain}. The rational certificates may integrate
only a subset of these nonnegative gains, so their endpoints are
lower bounds, rather than assertions that the selected boxes exhaust
the ideal domains.

For the fourth-node restoration, let $\mathcal R_{h_{\rm s}}$ be the disjoint
union of the rational slanted slabs specified by
\eqref{sv:restore-gates}. On each slab a cap $\lambda_4$ is chosen
before the fourth cube, and
$C_3\le S_3+3t_4<\lambda_4$, with the full and selected support
conditions retained. Its additional gain is
\begin{equation}
 T_{h_{\rm s}}=4\int_{\mathcal R_{h_{\rm s}}}
       \overline f\left(\frac{\lambda_4-S_4}{t_4}\right)\,{d_{\rm s}}\mu_4.
 \label{num:restoration-integral}
\end{equation}
The exact slab volume is~\eqref{sv:slab-volume}. This term enters the
Chen functional with coefficient $1/4$, and its domain is disjoint
from the other gains already counted.

\subsection{The fixed upper-cost slope and table changes}
\label{num:cost-slope}
The coefficient $A_{\rm cost}$ used below is a safe slope of a specified affine
upper enclosure of five costs. It need not be the derivative of a
minimum of kernels. To state its components, put $K_0=6.7619$ and
$A_+(s)=s\overline F(s)$. The three one-dimensional switched
costs are
\begin{align*}
 \frac{U_7}{K_0}&=\int_{d_{\rm s}}^{1/3}
  \frac{\log((1-2t)/t)+A_+(1/t-1)-1}{t(1-t)}\,dt,\\
 \frac{U_8}{K_0}&=\int_{a_{\rm s}}^{1/3}\frac{\log(2-3t)}{t(1-t)}\,dt,\\
 \frac{U_9}{K_0}&=\int_{h_{\rm s}}^{d_{\rm s}}\frac{\log((1-t-{d_{\rm s}})/{d_{\rm s}})}{t(1-t)}\,dt.
\end{align*}
The extra ordinary-function contribution in the first integral can
be evaluated without interpolation:
\[
 \int_{d_{\rm s}}^{1/4}\frac{A_+(1/t-1)-1}{t(1-t)}\,dt
 =\int_2^{1/{d_{\rm s}}-2}\frac{\log(u-1)}u
                  \log\frac{1-{d_{\rm s}}}{{d_{\rm s}}(u+1)}\,du.
\]
Let $\underline U_j$ denote the rational lower endpoints for these
$K_0$-linear costs from the fourth-derivative Simpson enclosures
specified in \texttt{five\_cost\_slope\_certificate}. The root cost
upper endpoints may use the separate midpoint enclosures; an upper
endpoint and a lower endpoint have different roles here.
Their contribution to a safe slope is
$A_{{\rm cost},789}=(2\underline U_7+\underline U_8+\underline U_9)/(4K_0)$.
For a retained four-prefix cell $\mathscr B$ outside the $40000$ priority
parents, write $A_{\mathscr B}$ for its assigned upper mass at $K_0$ and $s_{\mathscr B}$
for the supplied nonnegative affine slope. The required cellwise
conditions are
\begin{equation}
 0\le s_{\mathscr B}\le A_{\mathscr B}/K_0,\qquad A_{\mathscr B}\ge K_0\mathcal T_{\mathscr B},
 \label{num:affine-majorant-slope}
\end{equation}
where $\mathcal T_{\mathscr B}$ is the exact nonnegative switched mass. They give
\[
 A_{\mathscr B}-(K_0-K)s_{\mathscr B}\ge(K/K_0)A_{\mathscr B}\ge K\mathcal T_{\mathscr B}.
\]
No inequality $s_{\mathscr B}\le\mathcal T_{\mathscr B}$ is required.
At denominator $\mathsf Q=2^{80}$ the retained row slopes are
$s_{\mathscr B}=10000d_{\mathscr B}/\mathsf Q$. Their sums for $U_{10},U_{11}$ are
$1277985388931956822$ and $7976183963523684871$, respectively.
Consequently
\begin{equation}
 A_{\rm cost}=A_{{\rm cost},789}+
 \frac{10000(1277985388931956822+7976183963523684871)}{4\mathsf Q}.
 \label{num:five-cost-slope}
\end{equation}
The complete retained-cell inequalities are indexed by the
\texttt{feedback\_original\_slope} entry of the certificate map.

For two valid defect envelopes $H_m\le H_n$, define
\begin{equation}
 \mathcal T(H_m,H_n)=
 \int_{a_{\rm s}}^{1/3}\frac{(H_n-H_m)(s_3(t))}{tL_3(t)}\,dt
 +\int_{a_{\rm s}}^{d_{\rm s}}\frac{(H_n-H_m)(s_4(t))}{tL_4(t)}\,dt,
 \label{num:table-change-integral}
\end{equation}
where $s_j(t)=\max(2,L_j(t)/{a_{\rm s}})$, with zero contribution outside
$[1,5]$, and, for $e=57/224$,
\[
 L_3(t)=\begin{cases}
 (1-t)/2,&{a_{\rm s}}<t<\mu_0,\\
 57/64-5t/2,&\mu_0<t<e,\\
 1/2-t,&e<t<1/3.
 \end{cases}
\]
The function $L_4$ uses the first two branches restricted to $({a_{\rm s}},{d_{\rm s}})$.
The factors $4$ in the two upper counts cancel the single Chen
factor $1/4$. Thus
\begin{equation}
 T_{312}=\mathcal T(H_{248},H_{312}),\qquad
 T_{430}=\mathcal T(H_{312},H_{430}).
 \label{num:successive-table-gains}
\end{equation}
At a knot the lookup uses the smaller adjacent value; on the
positive-measure branch $s_j=2$ this means the minimum of cells $49$
and $50$. On an affine piece $L=A_0-B_0t$ the exact inverse-cell
primitive is
\[
 \int_l^r\frac{dt}{t(A_0-B_0t)}
 =\frac1{A_0}\log\frac{r(A_0-B_0l)}{l(A_0-B_0r)}.
\]
This determines each table-change interval and its rational logarithm
enclosure without numerical interpolation across a knot.

\subsection{Mathematical objects and certificate fields}
\label{num:certificate-chain}
Table~\ref{num:certificate-table} links the count estimates, their
integral bounds and the ancillary certificates. Paths in the last
column are targets in \filepath{certificate-map.json}; each target has
the source filename, exact JSON pointer, direction and any rational
extraction formula. The map also includes
the enclosure and coverage assertions for each primitive family and
the corresponding full-witness filenames.
The map's \texttt{domain\_constructions} entries identify the actual
$C_3,C_4,C_5$ region generators and their selected records, including
the two third-rank starts, the fourth-rank chord and oblique regions,
the fifth-rank child increments, and the slanted restoration.
For $U_{10},U_{11}$ they identify the generator
\filepath{witnesses/generator-0041.py} and the complete ordered witness
\filepath{witnesses/record-0638.json}.
The witness records specify the parent policy, method choices,
generation-$248$ profile, and full facet descriptions of these domains.

\begingroup
\small
\setlength{\tabcolsep}{4pt}
\begin{longtable}{@{}>{\raggedright\arraybackslash}p{.15\textwidth}>{\raggedright\arraybackslash}p{.43\textwidth}>{\raggedright\arraybackslash}p{.35\textwidth}@{}}
\caption{Definitions, directions and certificate targets.}
\label{num:certificate-table}\\
\toprule Object & Mathematical estimate & Certificate target\\
\midrule
\endfirsthead
\toprule Object & Mathematical estimate & Certificate target\\
\midrule
\endhead
$G_j$, $1\le j\le22$ &
The specified child--parent density differences in the upper tree;
nonnegative savings may use a subdomain. &
\filepath{/upper_saving_bounds/}\,$(j-1)$; lower endpoint.\\[3pt]
Finite actions &
\eqref{sv:one-step-operator}, \eqref{sv:two-step-base},
\eqref{sv:two-step-feedback}; complete $\Gamma_5$ with separate payment
on $D_5^-$. &
\filepath{/actions/}\,$k$; lower base and feedback weights;
\filepath{/gamma5_cost_decompositions}.\\[3pt]
Seed and tables &
Lemma~\ref{sv:finite-seed} and the simultaneous map
\eqref{num:action-map}; closed-cell estimates. &
\filepath{/seed_certificate}, \filepath{/defect_tables/}\,$m$;
lower endpoints, $m=248,312,430$.\\[3pt]
$L_i,U_i$ &
The integrals in Section~\ref{num:eleven-integrals}; complete retained
domain for each negative cost, lower estimate for each positive term. &
\filepath{/lower_root_bounds}; lower for $L_i$, upper for $U_i$.\\[3pt]
$B_{\rm ord},H_{\rm rem}$ &
\eqref{num:ordinary-start-difference}, \eqref{num:removed-defect};
remove the entire generation-$248$ defect. &
\filepath{/lower_auxiliary_bounds}; lower for $B_{\rm ord}$, upper for
$H_{\rm rem}$.\\[3pt]
$I_r(w),T_{h_{\rm s}}$ &
\eqref{num:third-lower-gain}--\eqref{num:restoration-integral}; one
common lower policy and disjoint restoration slabs. &
\filepath{/lower_integral_bounds} and
\filepath{/lower_auxiliary_bounds/restoration_L2_lower}; lower.\\[3pt]
$A_{\rm cost}$, $T_{312},T_{430}$ &
\eqref{num:affine-majorant-slope}--\eqref{num:successive-table-gains};
upper-cost slope and successive table changes. &
\filepath{/five_cost_slope_certificate},
\filepath{/lower_auxiliary_bounds}; safe lower coefficients.\\[3pt]
Paired deletion &
Assigned positive majorant and discarded lower estimator on
$\mathscr E_0\cup \mathscr E_1$; \eqref{num:paired-extra}, \eqref{num:saved-loss}. &
\filepath{/paired_rectangles}; exact geometry, mass intervals and lower
estimator integers.\\[3pt]
$E_{\rm up},R_{\rm low}$ and reserves &
The unique combinations in~\eqref{num:lower-sum} and the upper
density expression. &
\filepath{/scalar_bounds}; exact sums and strict surpluses.\\
\bottomrule
\end{longtable}
\endgroup

For each negative cost, the mapped witnesses specify a complete
cover of its retained domain. For a positive saving, they specify
the selected nonnegative subfamily on which its lower bound is
evaluated.

\subsection{The lower functional}
For the signed eleven-term expression, use a lower endpoint for
each positive term and an upper endpoint for each negative term.
The following integers, divided by $10^{14}$, give such endpoints:
\begin{center}
\begin{tabular}{@{}l r r@{}}
\toprule Term & Coefficient before division by $4$ & Endpoint numerator\\
\midrule
$L_1$ & $3$ & $1346406373358053$\\
$L_2$ & $1$ & $710262601191162$\\
$L_5$ & $1$ & $237450056438739$\\
$L_6$ & $1$ & $301984961684469$\\
$U_3$ & $-1$ & $2031242249325404$\\
$U_4$ & $-1$ & $1405868301028074$\\
$U_7$ & $-2$ & $138343028785273$\\
$U_8$ & $-1$ & $405707594761181$\\
$U_9$ & $-1$ & $325703365221150$\\
$U_{10}+U_{11}$ & $-1$ & $67568703438104$\\
\bottomrule
\end{tabular}
\end{center}
The exact rational endpoints, rather than these shortened ones,
define $C_{11}$.  The four-prefix minimum covers all
\[
 \binom{51}{4}+48\binom{50}{3}=1190700
\]
ordered rows, including the tied-coordinate simplices.
The priority subdivision replaces $40000$ parents by maxima
with their complete sums of $516779$ children.  Its extra
coefficient saving is
\[
 P=\frac{106520738437439293219}{604462909807314587353088}.
\]

The terms $B_{\rm ord},H_{\rm rem}$, $I_r(a_{\rm s}),I_r(h_{\rm s})$, $T_{h_{\rm s}}$ are
defined in~\eqref{num:ordinary-start-difference}--\eqref{num:restoration-integral}.
Their exact bounds are
specified by \texttt{lower\_auxiliary\_bounds} and
\texttt{lower\_integral\_bounds}.  At $K_0=6.7619$ put
\[
 C_{K_0}=C_{11}+P+\frac14\left(
 B_{\rm ord}-H_{\rm rem}
 +\sum_{r=3}^5\{3I_r({a_{\rm s}})+I_r({h_{\rm s}})\}+T_{h_{\rm s}}\right).
\]
The third-rank integrals cover both cutoff regions.  The
fourth-rank integrals retain all specified masks and the
disjoint additional bands.  The fifth-rank integration uses
$158251$ floors and $23662$ additional slab or tied cells on
the ${h_{\rm s}}$ side and $679$ integrals on the ${a_{\rm s}}$ side.  The slanted
restoration uses $98668$ boxes and $789344$ layers.  Only
nonnegative gains may omit boundary strips.  An asymmetric
kernel is bounded on its ordered simplex; no factorial
symmetry is assumed for such a kernel.

With $A_{\rm cost}$ as in~\eqref{num:five-cost-slope} and
$T_{312},T_{430}$ as in~\eqref{num:successive-table-gains}, the
initial lower endpoint is bounded below by
\begin{equation}
 C_{K_0}+A_{\rm cost}(K_0-K)+T_{312}+T_{430}
 +0.002095161+0.003799514-0.002372094.
 \label{num:lower-sum}
\end{equation}
Direct rational substitution is at least $R_{\rm base}$ in
\eqref{eq:base-R}.  The paired-cell refinement below adds
exactly $\Delta$ from \eqref{eq:current-refinement-delta}.
After allocating the two reserves in \eqref{eq:current-B},
the resulting lower endpoint exceeds $c$ by
\[
 \frac{13457260926863633771775678584711245753}
 {2742952509978644604793778369418231808000000000}>0.
\]
Together with the upper calculation, this proves the scalar
inequalities in Lemma~\ref{lem:numerical-endpoints} from the
stated finite enclosures.

\section{A paired-cell refinement}\label{app:refinement}

The paired rectangles are specified without a geometric
rounding convention as follows.  Put
\[
 t_i=\mu_0+i\frac{{d_{\rm s}}-\mu_0}{512},\qquad
 u_j={a_{\rm s}}+j\frac{{h_{\rm s}}-{a_{\rm s}}}{512}\qquad(0\le i,j\le512).
\]
Each row of \texttt{paired\_rectangles} specifies a family
$\mathscr E_0$ or $\mathscr E_1$ and the rectangle
$[t_i,t_{i+1}]\times[u_j,u_{j+1}]$.  The families have
$6091$ and $11481$ cells, respectively, with no common
cell.  Their boundaries have zero integration measure.
For a rectangle $\mathscr B=[t_l,t_r]\times[u_l,u_r]$, define
\[
 Z_{\mathscr B}=\min\left(u_l,\frac{3(1-t_r-u_r)}{13}\right).
\]
The rational data satisfy $Z_{\mathscr B}>{a_{\rm s}}$ and
$1-t_r-u_r-2Z_{\mathscr B}\ge(7/3)Z_{\mathscr B}$ in every row.  Hence
\[
 \frac{1-t-u-z-w}{z}\ge\frac73
 \qquad({a_{\rm s}}<w<z<Z_{\mathscr B},\ (t,u)\in \mathscr B).
\]

On each rectangle, Lemma~\ref{lem:current-omega-refinement} gives
$\omega\ge1/2$ throughout the fibre and $\omega\ge11/20$
on its specified inner part. The second estimate improves a
subset of the first contribution and is used as an increment.

The positive parent majorant can therefore be bounded below
by the contribution from $1/2$ on its whole fibre plus
$1/20$ on this inner fibre.  These are simultaneous
nonnegative contributions.  After integration and the
single Chen factor $1/4$, the additional contribution is
\begin{equation}
 G_{\mathscr B}=\frac K{80}\log\frac{t_r}{t_l}\log\frac{u_r}{u_l}
 \left\{\frac1{a_{\rm s}}-\frac{1+\log(Z_{\mathscr B}/{a_{\rm s}})}{Z_{\mathscr B}}\right\}.
 \label{num:paired-extra}
\end{equation}
Indeed the $w$ integral is $\log(z/{a_{\rm s}})$ and
\[
 \int_{a_{\rm s}}^{Z_{\mathscr B}}\frac{\log(z/{a_{\rm s}})}{z^2}\,dz
 =\frac1{a_{\rm s}}-\frac{1+\log(Z_{\mathscr B}/{a_{\rm s}})}{Z_{\mathscr B}}.
\]
The normalization agrees with~\eqref{eq:current-extra-fibre}.

Every row of the data contains rational lower and upper
endpoints for \eqref{num:paired-extra}, represented by
integer numerators at denominator $2^{160}$.  Applying
\eqref{num:log-enclosure} to its three logarithms and
summing the lower endpoints gives
\begin{equation}
 \sum_{\mathscr B\in \mathscr E_0\cup \mathscr E_1}G_{\mathscr B}>\frac{558583}{10^9}.
 \label{num:gross-bound}
\end{equation}
The sums on the two families are approximately
$1.9318091250\cdot10^{-4}$ and
$3.6540308174\cdot10^{-4}$, respectively. Only the rational
lower bound \eqref{num:gross-bound} is used.

\section{Discarded lower estimators and the final reserves}
\label{app:refinement-review}

The paired deletion removes a positive term together with its
assigned lower estimator. To bound the resulting loss, we give
both the exact finite sum of that estimator and a direct rational
upper enclosure.

On a cell write
\[
 R(t,u)=\frac{57}{64}-\frac52t-u,\qquad
 r=R(t_r,u_r),\quad R_0=R(t_l,u_l),\quad
 M=\log(t_r/t_l)\log(u_r/u_l).
\]
Cells of $\mathscr E_0$ have $R_0/{a_{\rm s}}\le2$ and zero assigned lower estimator.
There are $227$ cells of $\mathscr E_1$ with $r/{a_{\rm s}}<2$; their
assigned lower estimator is also zero.  On the remaining $11254$ cells,
$2\le r/{a_{\rm s}}\le R_0/{a_{\rm s}}\le9/4$.  Put
\[
 f_l=\frac{\log(r/{a_{\rm s}}-1)}{r/{a_{\rm s}}},\qquad
 f_h=\frac{\log(R_0/{a_{\rm s}}-1)}{R_0/{a_{\rm s}}},
\]
and
\[
 M_R=\frac{57}{64}M
 -\frac52(t_r-t_l)\log(u_r/u_l)
 -(u_r-u_l)\log(t_r/t_l).
\]
Since $r\le R(t,u)\le R_0$, the two assigned ordinary
estimators, in Chen-quarter units, are bounded above by
\begin{equation}
 \frac{f_lM}{{a_{\rm s}}},\qquad
 \frac{f_l(R_0M-M_R)+f_h(M_R-rM)}{a_{\rm s}(R_0-r)}.
 \label{num:ordinary-proxy}
\end{equation}
The coefficients of $f_l,f_h$ are nonnegative.
Replacing lower endpoints for these functions by their
exact values therefore increases the estimator.  The
maximum of the two expressions in
\eqref{num:ordinary-proxy} remains an upper bound for
the maximum of the two assigned alternatives.

Let $h_{\mathscr B}^{\min}$ be the minimum of the depth-$248$ lower table
entries over all closed cells meeting $[r/{a_{\rm s}},R_0/{a_{\rm s}}]$.
Both adjoining cells are included at a knot.  The exact
support is contained in the outward-rounded lookup
support of the assigned estimator, so $h_{\mathscr B}^{\min}$ is at least
the minimum used there.  Thus the additional term
$h_{\mathscr B}^{\min}M/R_0$ bounds its assigned defect estimator from above.
This comparison uses minima of finite lower estimators,
not an upper estimate for the mathematical function $h$.
Adding it to the maximum in
\eqref{num:ordinary-proxy} gives the required cell proxy.

The data record the assigned lower estimator on every cell as
an integer at denominator $\mathsf Q=2^{80}$.  Their sum is
$11290806845175652020270$.  After division by $4\mathsf Q$,
\begin{equation}
 L_{\mathscr E}
 =\frac{5645403422587826010135}
 {2417851639229258349412352}
 <\lambda_{\rm loss}:=\frac{466977}{200000000}.
 \label{num:saved-loss}
\end{equation}
An outward evaluation of the proxies above gives an
upper endpoint below the same $\lambda_{\rm loss}$.  For this
evaluation one may reduce logarithm arguments to
$[1,3/2]$, use the $320$-term alternating series for
$\log(1+x)$ and bound its remainder by the next term.
All its divisors have positive interval lower endpoints.
The resulting proxy sum is enclosed near
$0.0023348841363598131014744273487904$.

The retained whole $L_6$ endpoint before the deletion is
$3650774173156900914101685/\mathsf Q$.  Its nonenhanced half
has the cell decomposition just used.  Consequently
subtracting $4\lambda_{\rm loss}$ from that endpoint bounds the
retained $L_6$ from below; the enhanced half is unaffected.
The decrease of the paid loss relative to the value
in \eqref{num:lower-sum} is exactly
\[
 0.002372094-\lambda_{\rm loss}=\frac{37209}{10^9}.
\]
Combining this with \eqref{num:gross-bound} gives the
raw increment
\[
 \Delta=\frac{558583+37209}{10^9}
       =\frac{37237}{62500000}.
\]
Allocate $\Delta/8$ to each of the finite-approximation
and coefficient-realization reserves.  The net retained
increment is therefore $3\Delta/4=0.000446844$.
All estimates refer to this same fixed set of rectangles,
and each density contribution appears once.  They give
$R_{\rm low}=R_{\rm base}+\Delta$ and
$B_{\rm low}=B_{\rm base}+\Delta/8$, as in
\eqref{eq:current-R} and \eqref{eq:current-B}.

\section*{Acknowledgment of AI assistance}
OpenAI ChatGPT/Codex assisted with mathematical exploration, proof development
and assessment, parameter searches, computational verification, and editorial
preparation. The author is responsible for the mathematical content and final
manuscript.

\begin{thebibliography}{99}
\bibitem{ABL}
E. Assing, V. Blomer and J. Li,
\emph{Uniform Titchmarsh divisor problems},
\href{https://arxiv.org/abs/2005.13915v1}{arXiv:2005.13915v1} (2020).

\bibitem{ChanSupplement}
S. Chan,
\emph{Theorem $(1+1.75)$ on the Goldbach Conjecture --- Numerical Supplement},
version 20261106, Zenodo.
\href{https://doi.org/10.5281/zenodo.23755453}{doi:10.5281/zenodo.23755453}.

\bibitem{Drappeau}
S. Drappeau,
\emph{Sums of Kloosterman sums in arithmetic progressions, and the error
term in the dispersion method}, Proc. Lond. Math. Soc. (3)
\textbf{114} (2017), no.~4, 684--732.
\href{https://doi.org/10.1112/plms.12022}{doi:10.1112/plms.12022}.

\bibitem{Kedlaya}
K. S. Kedlaya,
\emph{Analytic Number Theory}, Chapter~18, Theorem~18.2,
\url{https://kskedlaya.org/ant/chapter-18.html}.

\bibitem{KimSarnak}
H. Kim and P. Sarnak,
\emph{Refined estimates towards the Ramanujan and Selberg conjectures},
J. Amer. Math. Soc. \textbf{16} (2003), 175--181,
Appendix~2 to H. Kim, \emph{Functoriality for the exterior square of
$\mathrm{GL}_4$ and the symmetric fourth of $\mathrm{GL}_2$}.
\href{https://doi.org/10.1090/S0894-0347-02-00410-1}{doi:10.1090/S0894-0347-02-00410-1}.

\bibitem{LiLiu}
J. Li and J. Liu,
\emph{Theorem $(1+1.9)$ on the Goldbach Conjecture},
\href{https://arxiv.org/abs/2606.05224v2}{arXiv:2606.05224v2} (2026).

\bibitem{RunboLi}
R. Li,
\emph{On Chen's theorem, Goldbach's conjecture and almost prime twins II},
Math. Rep. (Bucur.) \textbf{28(78)} (2026), nos.~1--2, 39--61.
\href{https://doi.org/10.59277/mrar.2026.28.78.1.2.39}{doi:10.59277/mrar.2026.28.78.1.2.39};
\href{https://arxiv.org/abs/2405.05727v4}{arXiv:2405.05727v4}.

\bibitem{Lichtman}
J. D. Lichtman,
\emph{Primes in arithmetic progressions to large moduli, and Goldbach
beyond the square-root barrier},
\href{https://arxiv.org/abs/2309.08522v1}{arXiv:2309.08522v1} (2023).

\bibitem{MaynardI}
J. Maynard,
\emph{Primes in arithmetic progressions to large moduli I: Fixed residue
classes},
\href{https://arxiv.org/abs/2006.06572v2}{arXiv:2006.06572v2} (2021).

\bibitem{MaynardII}
J. Maynard,
\emph{Primes in arithmetic progressions to large moduli II: Well-factorable
estimates},
\href{https://arxiv.org/abs/2006.07088v1}{arXiv:2006.07088v1} (2020).

\bibitem{Pascadi}
A. Pascadi,
\emph{Large sieve inequalities for exceptional Maass forms and the
greatest prime factor of $n^2+1$}, Forum Math. Pi \textbf{14} (2026), e8.
\href{https://doi.org/10.1017/fmp.2026.10025}{doi:10.1017/fmp.2026.10025};
\href{https://arxiv.org/abs/2404.04239v3}{arXiv:2404.04239v3}.

\bibitem{Shiu}
P. Shiu,
\emph{A Brun--Titchmarsh theorem for multiplicative functions},
J. Reine Angew. Math. \textbf{313} (1980), 161--170.
\href{https://doi.org/10.1515/crll.1980.313.161}{doi:10.1515/crll.1980.313.161}.

\bibitem{Suzuki}
Y. Suzuki,
\emph{The Rosser--Iwaniec sieve}, unpublished notes, 14 March 2023,
102 pp., Proposition~11.8 and Lemma~14.7.
\url{https://www.suzuyu1729.com/notes/NagoyaSeminar2022/RosserIwaniec.pdf}.

\bibitem{TaoSW}
T. Tao,
\emph{254A, Notes 2: Complex-analytic multiplicative number theory},
9 December 2014, Exercise~64,
\url{https://terrytao.wordpress.com/2014/12/09/254a-notes-2-complex-analytic-multiplicative-number-theory/}.

\bibitem{Wu2004}
J. Wu,
\emph{Chen's double sieve, Goldbach's conjecture and the twin prime
problem}, Acta Arith. \textbf{114} (2004), 215--273.
\href{https://doi.org/10.4064/aa114-3-2}{doi:10.4064/aa114-3-2}.

\bibitem{Wu}
J. Wu,
\emph{Chen's double sieve, Goldbach's conjecture and the twin prime
problem, 2}, Acta Arith. \textbf{131} (2008), 367--387.
\href{https://arxiv.org/abs/0709.3764v1}{arXiv:0709.3764v1}.

\bibitem{Young}
M. P. Young,
\emph{Explicit calculations with Eisenstein series},
\href{https://arxiv.org/abs/1710.03624v2}{arXiv:1710.03624v2} (2017).
\end{thebibliography}

\end{document}
